\chapter{Deret Fourier dan PDE} \label{FS:chapter}

%%%%%%%%%%%%%%%%%%%%%%%%%%%%%%%%%%%%%%%%%%%%%%%%%%%%%%%%%%%%%%%%%%%%%%%%%%%%%%

\section{Masalah nilai batas} \label{bvp:section}

%mbxINTROSUBSECTION

\sectionnotes{2 kuliah\EPref{, serupa dengan \S3.8 dalam \cite{EP}}\BDref{,
\S10.1 dan \S11.1 dalam \cite{BD}}}

\subsection{Masalah nilai batas}

Sebelum membahas deret Fourier, kita mempelajari
apa yang disebut
\emph{masalah nilai batas\index{masalah nilai batas}}
(atau \emph{masalah titik ujung\index{masalah titik ujung}}).  Perhatikan
\begin{equation*}
x'' + \lambda x = 0, \quad x(a) = 0, \quad x(b) = 0,
\end{equation*}
untuk suatu konstanta $\lambda$, dengan $x(t)$ didefinisikan untuk $t$ dalam interval
$[a,b]$.
Sebelumnya kita menentukan nilai solusi dan turunannya
di satu titik.  Sekarang kita menentukan nilai solusi di dua titik
yang berbeda.  Karena $x=0$ adalah solusi, keberadaan
solusi bukanlah masalah.  Ketunggalan solusi merupakan persoalan lain.
Solusi umum untuk $x'' + \lambda x = 0$ memiliki dua
konstanta sembarang\footnote{%
Lihat \subsectionvref{subsection:fourfundamental} atau \examplevref{example:expsecondorder} dan
\examplevref{example:sincossecondorder}.}.
Oleh karena itu,
wajar (tetapi keliru) untuk meyakini bahwa menetapkan dua
syarat menjamin solusi tunggal.

\begin{example}
Ambil $\lambda = 1$,
$a=0$, $b=\pi$.  Artinya,
\begin{equation*}
x'' + x = 0, \quad x(0) = 0, \quad x(\pi) = 0.
\end{equation*}
Maka $x = \sin t$ adalah solusi lain (selain $x=0$) yang memenuhi kedua syarat
batas.  Masih ada solusi lainnya.  Tuliskan solusi umum
persamaan diferensial tersebut, yaitu $x= A \cos t + B \sin t$.
Syarat $x(0) = 0$ memaksa $A=0$.  Menetapkan $x(\pi) = 0$ tidak
memberi kita informasi tambahan karena $x = B \sin t$ sudah memenuhi kedua
syarat batas.
Jadi, terdapat tak hingga banyak solusi berbentuk $x = B \sin t$,
dengan $B$ suatu konstanta sembarang.
\end{example}

\begin{example}
Di sisi lain, perhatikan $\lambda = 2$.  Artinya,
\begin{equation*}
x'' + 2 x = 0, \quad x(0) = 0, \quad x(\pi) = 0.
\end{equation*}
Maka solusi umumnya adalah
$x= A \cos \bigl( \sqrt{2}\,t\bigr)
+ B \sin \bigl( \sqrt{2}\,t\bigr)$.
Menetapkan $x(0) = 0$ tetap memaksa $A = 0$.
Kita menerapkan syarat kedua untuk memperoleh
$0=x(\pi) = B \sin \bigl( \sqrt{2}\,\pi\bigr)$.
Karena $\sin \bigl( \sqrt{2}\,\pi\bigr) \neq 0$, kita memperoleh
$B = 0$.  Oleh karena itu, $x=0$ adalah solusi tunggal masalah ini.
\end{example}

Apa yang sedang terjadi?  Kita akan mencari konstanta
$\lambda$ mana yang memungkinkan solusi tak nol, dan kita akan mencari
solusi-solusi tersebut.  Masalah ini merupakan analog dari pencarian
nilai eigen dan vektor eigen matriks.

\subsection{Masalah nilai eigen}

Untuk teori dasar deret Fourier, kita memerlukan
tiga masalah nilai eigen berikut:
\begin{equation} \label{bv:eq1}
x'' + \lambda x = 0, \quad x(a) = 0, \quad x(b) = 0 ,
\end{equation}
\begin{equation} \label{bv:eq2}
x'' + \lambda x = 0, \quad x'(a) = 0, \quad x'(b) = 0 ,
\end{equation}
dan
\begin{equation} \label{bv:eq3}
x'' + \lambda x = 0, \quad x(a) = x(b), \quad x'(a) = x'(b) .
\end{equation}
Suatu bilangan $\lambda$ disebut
\emph{nilai eigen\index{nilai eigen dari masalah nilai batas}}
dari \eqref{bv:eq1}
(berturut-turut,\ \eqref{bv:eq2} atau \eqref{bv:eq3}) jika dan hanya jika
terdapat solusi tak nol (tidak identik nol) untuk \eqref{bv:eq1}
(berturut-turut,\ \eqref{bv:eq2} atau \eqref{bv:eq3})
untuk $\lambda$ tertentu tersebut.  Suatu
solusi tak nol disebut
\emph{\myindex{fungsi eigen}}\index{fungsi eigen yang bersesuaian} yang bersesuaian.
Kita akan mempertimbangkan persamaan dan syarat batas yang lebih umum,
tetapi pembahasan ini kita tunda hingga
\chapterref{SL:chapter}.

Perhatikan kemiripannya dengan nilai eigen dan vektor eigen matriks.  Kemiripan
tersebut bukan sekadar kebetulan.  Jika kita memandang persamaan-persamaan itu sebagai
operator diferensial, kita sebenarnya melakukan hal yang persis sama.
Pandang fungsi $x(t)$
sebagai vektor dengan tak hingga banyak komponen (satu untuk setiap $t$).
Misalkan $L = -\frac{d^2}{{dt}^2}$ adalah operator linear.
Maka pasangan nilai eigen/fungsi eigen seharusnya berupa $\lambda$ dan
$x$ tak nol sedemikian sehingga $Lx = \lambda x$.
Dengan kata lain,
kita mencari fungsi tak nol $x$
yang memenuhi syarat titik ujung tertentu dan menyelesaikan
$(L- \lambda)x = 0$.  Banyak formalisme dari aljabar linear tetap
berlaku di sini, meskipun kita tidak akan menelusuri jalur penalaran ini terlalu jauh.

\begin{example} \label{bvp:eig1ex}
Carilah nilai eigen dan fungsi eigen dari
\begin{equation*}
x'' + \lambda x = 0, \quad x(0) = 0, \quad x(\pi) = 0 .
\end{equation*}

Kita harus menangani
kasus $\lambda > 0$, $\lambda = 0$, dan $\lambda < 0$ secara terpisah.
Pertama, andaikan $\lambda > 0$.  Maka
solusi umum untuk $x''+\lambda x = 0$ adalah
\begin{equation*}
x = A \cos \bigl( \sqrt{\lambda}\, t\bigr)
+ B \sin \bigl( \sqrt{\lambda}\, t\bigr).
\end{equation*}
Syarat $x(0) = 0$ langsung mengakibatkan $A = 0$.
Selanjutnya
\begin{equation*}
0 = x(\pi) = B \sin \bigl( \sqrt{\lambda}\, \pi \bigr) .
\end{equation*}
Jika $B$ nol, maka $x$ bukan solusi tak nol.  Jadi, untuk memperoleh solusi
tak nol, kita harus memiliki $\sin \bigl( \sqrt{\lambda}\, \pi\bigr) = 0$.  Oleh karena itu,
$\sqrt{\lambda}\, \pi$ harus merupakan kelipatan bulat dari $\pi$.
Dengan kata lain, $\sqrt{\lambda} = k$ untuk suatu bilangan bulat positif $k$.
Jadi, nilai eigen positifnya adalah
$k^2$ untuk semua bilangan bulat $k \geq 1$.  Fungsi eigen yang bersesuaian
dapat diambil sebagai $x=\sin (k t)$.  Sama seperti untuk vektor eigen, kelipatan
konstan suatu fungsi eigen juga merupakan fungsi eigen,
sehingga kita hanya perlu memilih satu.

Sekarang, andaikan $\lambda = 0$.  Dalam kasus ini, persamaannya adalah $x'' = 0$,
dan solusi umumnya $x = At + B$.  Syarat $x(0) = 0$ mengakibatkan
$B=0$, dan $x(\pi) = 0$ mengakibatkan $A = 0$.  Ini berarti bahwa $\lambda
= 0$ \emph{bukan} nilai eigen.

Terakhir, andaikan $\lambda < 0$.  Dalam kasus ini, kita memiliki solusi
umum\footnote{Ingat bahwa
$\cosh s = \frac{1}{2}(e^s+e^{-s})$
dan
$\sinh s = \frac{1}{2}(e^s-e^{-s})$.  Sebagai latihan,
cobalah perhitungannya dengan solusi umum yang ditulis sebagai
$x = A e^{\sqrt{-\lambda}\, t} + B e^{-\sqrt{-\lambda}\, t}$ (tentu saja untuk
$A$ dan $B$ yang berbeda).}
\begin{equation*}
x = A \cosh \bigl( \sqrt{-\lambda}\, t\bigr)
+ B \sinh \bigl( \sqrt{-\lambda}\, t \bigr) .
\end{equation*}
Menetapkan $x(0) = 0$ mengakibatkan $A = 0$ (ingat $\cosh 0 = 1$ dan $\sinh 0 =
0$).  Jadi solusi kita haruslah $x = B \sinh \bigl( \sqrt{-\lambda}\, t \bigr)$ dan memenuhi
$x(\pi) = 0$.  Hal ini hanya mungkin jika $B$ nol.  Mengapa?  Karena
$\sinh \xi$ hanya nol ketika $\xi=0$.  Anda sebaiknya menggambar grafik $\sinh$ untuk melihat
fakta ini.
Kita juga dapat melihatnya dari definisi $\sinh$.
Kita memperoleh $0 = \sinh \xi = \frac{e^\xi -
e^{-\xi}}{2}$.  Jadi $e^\xi = e^{-\xi}$, yang mengakibatkan $\xi = -\xi$, dan ini hanya
benar jika $\xi=0$.  Jadi tidak ada nilai eigen negatif.

Ringkasnya, nilai eigen dan fungsi eigen yang bersesuaian adalah
\begin{equation*}
\lambda_k = k^2 \qquad \text{dengan fungsi eigen} \qquad x_k = \sin (k t)
\qquad \text{untuk semua bilangan bulat } k \geq 1 .
\end{equation*}
\end{example}

\begin{example}
Hitunglah nilai eigen dan fungsi eigen dari
\begin{equation*}
x'' + \lambda x = 0, \quad x'(0) = 0, \quad x'(\pi) = 0 .
\end{equation*}

Sekali lagi kita harus menangani kasus $\lambda > 0$, $\lambda = 0$, dan $\lambda
< 0$ secara terpisah.
Pertama, andaikan $\lambda > 0$.
Solusi umum untuk $x''+\lambda x = 0$ adalah
$x = A \cos \bigl( \sqrt{\lambda}\, t\bigr) + B \sin \bigl( \sqrt{\lambda}\,
t\bigr)$.
Jadi
\begin{equation*}
x' = -A\sqrt{\lambda}\, \sin \bigl( \sqrt{\lambda}\, t\bigr) +
B\sqrt{\lambda}\, \cos \bigl(\sqrt{\lambda}\, t\bigr) .
\end{equation*}
Syarat $x'(0) = 0$ langsung mengakibatkan $B = 0$.
Selanjutnya
\begin{equation*}
0 = x'(\pi) = -A\sqrt{\lambda}\, \sin \bigl( \sqrt{\lambda}\, \pi\bigr) .
\end{equation*}
Sekali lagi, $A$ tidak boleh nol jika $\lambda$ hendak menjadi nilai eigen,
dan $\sin \bigl( \sqrt{\lambda}\, \pi\bigr)$ hanya nol
jika
$\sqrt{\lambda} = k$ untuk suatu bilangan bulat positif $k$.
Jadi, nilai eigen positifnya kembali berupa
$k^2$ untuk semua bilangan bulat $k \geq 1$.  Dan fungsi eigen yang bersesuaian
dapat diambil sebagai $x=\cos (k t)$.

Sekarang andaikan $\lambda = 0$.  Dalam kasus ini, persamaannya adalah $x'' = 0$
dan solusi umumnya $x = At + B$, sehingga $x' = A$.  Syarat
$x'(0) = 0$ mengakibatkan
$A=0$.  Syarat $x'(\pi) = 0$ juga mengakibatkan $A=0$.
Jadi $B$ dapat berupa apa saja (mari kita ambil 1).  Dengan demikian, $\lambda = 0$
merupakan nilai eigen dan $x=1$ merupakan fungsi eigen yang bersesuaian.

Terakhir, misalkan $\lambda < 0$.  Dalam kasus ini, solusi umumnya adalah
$x = A \cosh \bigl( \sqrt{-\lambda}\, t\bigr)
+ B \sinh \bigl( \sqrt{-\lambda}\, t\bigr)$
dan
\begin{equation*}
x' = A\sqrt{-\lambda}\, \sinh \bigl( \sqrt{-\lambda}\, t\bigr)
+ B\sqrt{-\lambda}\, \cosh \bigl( \sqrt{-\lambda}\, t \bigr) .
\end{equation*}
Kita sudah melihat (dengan peran $A$ dan $B$ dipertukarkan) bahwa agar ekspresi ini
bernilai nol di $t=0$ dan $t=\pi$, kita harus memiliki $A=B=0$.
Oleh karena itu, tidak ada
nilai eigen negatif.

Ringkasnya, nilai eigen dan fungsi eigen yang bersesuaian adalah
\begin{equation*}
\lambda_k = k^2 \qquad \text{dengan fungsi eigen} \qquad x_k = \cos (k t)
\qquad \text{untuk semua bilangan bulat } k \geq 1 ,
\end{equation*}
dan terdapat satu nilai eigen lain
\begin{equation*}
\lambda_0 = 0 \qquad \text{dengan fungsi eigen} \qquad x_0 = 1.
\end{equation*}
\end{example}

Masalah berikut adalah masalah yang menghasilkan deret Fourier
umum.

\begin{example} \label{bvp-periodic:example}
Hitunglah
nilai eigen dan fungsi eigen dari
\begin{equation*}
x'' + \lambda x = 0, \quad x(-\pi) = x(\pi), \quad x'(-\pi) = x'(\pi) .
\end{equation*}
Kita tidak menentukan nilai atau turunannya
di titik-titik ujung, tetapi menetapkan bahwa keduanya sama di awal dan
di akhir interval.

Kita lewati $\lambda < 0$.  Perhitungannya sama seperti sebelumnya,
dan kembali kita memperoleh
bahwa tidak ada nilai eigen negatif.

Untuk $\lambda = 0$, solusi umumnya adalah $x = At + B$.  Syarat
$x(-\pi) = x(\pi)$ mengakibatkan $A=0$ ($A\pi + B = -A\pi +B$ mengakibatkan $A=0$).
Syarat kedua $x'(-\pi) = x'(\pi)$ tidak mengatakan apa pun tentang $B$, sehingga
$\lambda=0$ merupakan nilai eigen dengan fungsi eigen yang bersesuaian $x=1$.

Untuk $\lambda > 0$ kita memperoleh
$x = A \cos \bigl( \sqrt{\lambda}\, t \bigr)
+ B \sin \bigl( \sqrt{\lambda}\, t\bigr)$.
Sekarang
\begin{equation*}
\underbrace{A \cos \bigl(-\sqrt{\lambda}\, \pi\bigr)
+ B \sin \bigl(-\sqrt{\lambda}\,
\pi\bigr)}_{x(-\pi)}
=
\underbrace{A \cos \bigl(  \sqrt{\lambda}\, \pi \bigr)
+ B \sin \bigl( \sqrt{\lambda}\,
\pi\bigr)}_{x(\pi)} .
\end{equation*}
Kita ingat bahwa $\cos (- \theta) = \cos (\theta)$ dan
$\sin (-\theta) = - \sin (\theta)$.  Oleh karena itu,
\begin{equation*}
A \cos \bigl(\sqrt{\lambda}\, \pi\bigr)
- B \sin \bigl( \sqrt{\lambda}\, \pi\bigr)
=
A \cos \bigl(\sqrt{\lambda}\, \pi\bigr)
+ B \sin \bigl( \sqrt{\lambda}\, \pi\bigr).
\end{equation*}
Jadi, $B=0$ atau $\sin \bigl( \sqrt{\lambda}\, \pi\bigr) = 0$.
Demikian pula (latihan), jika kita mendiferensialkan $x$ dan memasukkan syarat
kedua, kita memperoleh bahwa $A=0$ atau $\sin \bigl( \sqrt{\lambda}\, \pi\bigr) = 0$.
Oleh karena itu, kecuali jika kita menginginkan $A$ dan $B$ keduanya nol (dan kita tidak menginginkannya),
kita harus memiliki $\sin \bigl( \sqrt{\lambda}\, \pi \bigr) = 0$.  Jadi, $\sqrt{\lambda}$
adalah bilangan bulat dan nilai eigennya kembali berupa $\lambda = k^2$ untuk
bilangan bulat $k \geq 1$.  Namun, dalam kasus ini,
$x = A \cos (k t) + B \sin (k t)$ adalah fungsi eigen untuk sembarang $A$ dan sembarang $B$.
Jadi kita memiliki dua fungsi eigen yang bebas linear, yaitu $\sin (kt)$ dan $\cos (kt)$.
Ingat bahwa untuk suatu matriks, kita juga dapat memiliki dua vektor eigen
yang bersesuaian dengan satu nilai eigen jika nilai eigen tersebut berulang.

Ringkasnya, nilai eigen dan fungsi eigen yang bersesuaian adalah
\begin{align*}
& \lambda_k = k^2 & & \text{dengan fungsi-fungsi eigen} & &
\cos (k t) \quad \text{dan}\quad  \sin (k t)
 & & \text{untuk semua bilangan bulat } k \geq 1 , \\
& \lambda_0 = 0 & & \text{dengan fungsi eigen} & & x_0 = 1.
\end{align*}
\end{example}

\subsection{Keortogonalan fungsi eigen}

Sifat yang akan sangat berguna pada bagian berikutnya adalah
\emph{\myindex{keortogonalan}} fungsi-fungsi eigen. Sifat ini merupakan analog
dari fakta berikut tentang vektor eigen suatu matriks.  Suatu matriks
disebut
\emph{simetris\index{matriks simetris}}
jika $A = A^T$ (matriks tersebut sama dengan transposnya).
\emph{Vektor eigen untuk dua nilai eigen yang berbeda dari suatu matriks
simetris saling ortogonal.}
Operator
diferensial yang kita hadapi bertindak sangat menyerupai matriks simetris.
Oleh karena itu, kita memperoleh teorema berikut.

%\medskip
%
%Andaikan $\lambda_1$ dan $\lambda_2$ adalah dua nilai eigen yang berbeda dari $A$
%dan $\vec{v}_1$ dan $\vec{v}_2$ adalah vektor eigen yang bersesuaian.  Maka
%tentu saja kita memiliki $A \vec{v}_1 = \lambda_1 \vec{v}_1$ dan
%$A \vec{v}_2 = \lambda_2 \vec{v}_2$.
%\begin{equation*}
%\langle A \vec{v}_1 , \vec{v}_2 \rangle = \lambda_1 \langle \vec{v}_1 , \vec{v}_2 \rangle
%\qquad
%\langle A \vec{v}_2 , \vec{v}_1 \rangle = \lambda_2 \langle \vec{v}_2 , \vec{v}_1 \rangle
%\end{equation*}
%
%\begin{equation*}
%\langle A \vec{v}_1 , \vec{v}_2 \rangle -
%\langle A \vec{v}_2 , \vec{v}_1 \rangle 
%=
%(\lambda_1 - \lambda_2 ) \langle \vec{v}_1 , \vec{v}_2 \rangle
%\end{equation*}
%
%\begin{equation*}
%\langle (A-A^T) \vec{v}_1 , \vec{v}_2 \rangle
%=
%(\lambda_1 - \lambda_2 ) \langle \vec{v}_1 , \vec{v}_2 \rangle
%\end{equation*}

\begin{theorem} \label{bvp:orthogonaleigen}
Andaikan $x_1(t)$ dan $x_2(t)$ adalah dua fungsi eigen dari masalah
\eqref{bv:eq1}, \eqref{bv:eq2}, atau \eqref{bv:eq3}
untuk dua nilai eigen yang berbeda,
$\lambda_1$ dan $\lambda_2$.  Maka keduanya
\emph{ortogonal\index{ortogonal!fungsi}}
dalam arti bahwa
\begin{equation*}
\int_a^b x_1(t) x_2(t) \,dt = 0 .
\end{equation*}
\end{theorem}

Terminologi ini berasal dari fakta bahwa integral tersebut merupakan suatu jenis
hasil kali dalam.  Kita akan menguraikan hal ini lebih lanjut pada bagian berikutnya.  Teorema tersebut
memiliki pembuktian yang sangat singkat, elegan, dan mencerahkan, sehingga kita sajikan di sini.
Pertama, kita memiliki dua persamaan berikut.
\begin{equation*}
x_1'' + \lambda_1 x_1 = 0
\qquad \text{dan} \qquad
x_2'' + \lambda_2 x_2 = 0.
\end{equation*}
Kalikan persamaan pertama dengan $x_2$ dan persamaan kedua dengan $x_1$, lalu kurangkan untuk memperoleh
\begin{equation*}
(\lambda_1 - \lambda_2) x_1 x_2 = x_2'' x_1 - x_2 x_1'' .
\end{equation*}
Integralkan kedua ruas persamaan:
\begin{equation*}
\begin{split}
(\lambda_1 - \lambda_2) \int_a^b x_1(t) x_2(t) \,dt
& =
\int_a^b \bigl( x_2''(t) x_1(t) - x_2(t) x_1''(t) \bigr) \,dt \\
& =
\int_a^b \frac{d}{dt} \Bigl( x_2'(t) x_1(t) - x_2(t) x_1'(t) \Bigr) \,dt \\
& =
\Bigl[ x_2'(t) x_1(t) - x_2(t) x_1'(t) \Bigr]_{t=a}^b
= 0 .
\end{split}
\end{equation*}
Kesamaan terakhir berlaku karena syarat-syarat batas.  Sebagai contoh, jika
kita mempertimbangkan \eqref{bv:eq1}, kita memiliki $x_1(a) = x_1(b) = x_2(a) = x_2(b) = 0$
sehingga $x_2' x_1 - x_2 x_1'$ bernilai nol di $a$ maupun $b$.
Karena $\lambda_1 \neq \lambda_2$, teorema tersebut terbukti.

\begin{exercise}[mudah]
Selesaikan pembuktian teorema tersebut (periksa kesamaan terakhir dalam pembuktian) untuk kasus
\eqref{bv:eq2} dan \eqref{bv:eq3}.
\end{exercise}

Fungsi $\sin (n t)$ adalah fungsi eigen untuk masalah
$x''+\lambda x = 0$, $x(0) = 0$, $x(\pi) = 0$.
Oleh karena itu, untuk bilangan bulat positif
$n$ dan $m$, kita memiliki integral
\begin{equation*}
\int_{0}^\pi \sin (mt) \sin (nt) \,dt = 0 ,
\quad
\text{ketika } m \neq n.
\end{equation*}
Demikian pula, dengan tetap mengandaikan bahwa $m$ dan $n$ adalah bilangan bulat
positif,
\begin{equation*}
\int_{0}^\pi \cos (mt) \cos (nt) \,dt = 0 ,
\quad
\text{ketika } m \neq n,
\qquad \text{dan} \qquad
\int_{0}^\pi  \cos (nt) \,dt = 0 .
\end{equation*}
Terakhir, kita juga memperoleh
\begin{equation*}
\int_{-\pi}^\pi \sin (mt) \sin (nt) \,dt = 0 ,
\quad
\text{ketika } m \neq n, 
\qquad \text{dan} \qquad
\int_{-\pi}^\pi  \sin (nt) \,dt = 0 ,
\end{equation*}
\begin{equation*}
\int_{-\pi}^\pi \cos (mt) \cos (nt) \,dt = 0 ,
\quad
\text{ketika } m \neq n,
\qquad \text{dan} \qquad
\int_{-\pi}^\pi  \cos (nt) \,dt = 0 ,
\end{equation*}
dan
\begin{equation*}
\int_{-\pi}^\pi \cos (mt) \sin (nt) \,dt = 0 
\qquad \text{(bahkan jika $m=n$).}
\end{equation*}

%\medskip
%
%Teorema tersebut juga benar ketika syarat batas yang berbeda diterapkan.
%Sebagai contoh, jika kita mensyaratkan $x'(a) = x'(b) = 0$, atau
%$x(a) = x'(b) = 0$, atau
%$x'(a) = x(b) = 0$.  Lihat pembuktiannya.


%Berdasarkan apa yang telah kita lihat sebelumnya, kita menerapkan teorema untuk memperoleh integral
%\begin{equation*}
%\int_{-\pi}^\pi \sin (mt) \sin (nt) \,dt = 0 \qquad \text{dan} \qquad
%\int_{-\pi}^\pi \cos (mt) \cos (nt) \,dt = 0 ,
%\end{equation*}
%ketika $m \neq n$, dan
%\begin{equation*}
%\int_{-\pi}^\pi \sin (mt) \cos (nt) \,dt = 0 ,
%\end{equation*}
%untuk semua $m$ dan $n$.

\subsection{Alternatif Fredholm}

Sekarang kita membahas secara singkat sebuah teorema yang sangat berguna dalam teori persamaan
diferensial.  Teorema ini berlaku dalam konteks yang lebih umum daripada yang akan
kita nyatakan, tetapi untuk keperluan kita pernyataan berikut sudah
memadai.  Kita akan memberikan versi yang sedikit lebih umum dalam
\chapterref{SL:chapter}.

\begin{theorem}[Alternatif Fredholm%
\footnote{Dinamai menurut matematikawan Swedia
\href{https://en.wikipedia.org/wiki/Fredholm}{Erik Ivar Fredholm}
(1866--1927).}]\index{alternatif Fredholm!kasus sederhana}
\label{thm:fredholmsimple}
Tepat satu dari pernyataan berikut berlaku.
Entah
\begin{equation} \label{simpfredhomeq}
x'' + \lambda x = 0, \quad x(a) = 0, \quad x(b) = 0
\end{equation}
memiliki solusi tak nol, atau
\begin{equation} \label{simpfrednonhomeq}
x'' + \lambda x = f(t), \quad x(a) = 0, \quad x(b) = 0
\end{equation}
memiliki solusi tunggal untuk setiap fungsi $f$ yang kontinu pada $[a,b]$.
\end{theorem}

Teorema tersebut juga benar untuk jenis-jenis lain dari
syarat batas yang kita pertimbangkan.
Teorema ini berarti bahwa jika $\lambda$ bukan nilai eigen, persamaan tak homogen
\eqref{simpfrednonhomeq} memiliki solusi tunggal untuk setiap ruas
kanan.  Di sisi lain, jika $\lambda$ adalah nilai eigen, maka
\eqref{simpfrednonhomeq} belum tentu memiliki solusi untuk setiap $f$,
dan lebih lanjut,
bahkan jika kebetulan memiliki solusi, solusinya tidak
tunggal.

Kita juga ingin menegaskan gagasan bahwa operator diferensial linear memiliki
banyak kesamaan dengan matriks.  Jadi tidak mengherankan bahwa
terdapat pula versi berdimensi hingga dari alternatif Fredholm untuk matriks.
Misalkan $A$ adalah matriks $n \times n$.  Alternatif Fredholm kemudian
menyatakan bahwa entah $(A-\lambda I) \vec{x}
= \vec{0}$ memiliki solusi taktrivial, atau $(A-\lambda I) \vec{x} = \vec{b}$
memiliki solusi tunggal untuk setiap $\vec{b}$.

Banyak intuisi dari aljabar linear dapat diterapkan pada operator diferensial
linear, tetapi tentu saja kita harus berhati-hati.  Sebagai contoh, salah satu
perbedaan yang sudah kita lihat adalah bahwa secara umum suatu operator diferensial
memiliki tak hingga banyak nilai eigen, sedangkan sebuah matriks hanya memiliki berhingga banyak.

\subsection{Penerapan}

Mari kita perhatikan penerapan fisik suatu masalah titik ujung.
Andaikan kita memiliki tali elastis teregang kuat yang berputar cepat
dengan kerapatan linear seragam $\rho$, misalnya dalam
$\unitfrac{kg}{m}$.
Mari kita tempatkan masalah ini pada bidang $xy$, dengan $x$ dan $y$
keduanya dalam meter.  Sumbu $x$ menyatakan
posisi pada tali.  Tali berputar dengan kecepatan sudut $\omega$,
dalam $\unitfrac{radians}{s}$.
Bayangkan seluruh bidang $xy$ berputar dengan kecepatan sudut $\omega$.
Dengan demikian, tali tetap berada dalam bidang $xy$ ini dan $y$
mengukur simpangannya dari posisi kesetimbangan, $y=0$, pada sumbu $x$.
Jadi grafik $y$ memberikan bentuk tali.
Kita mempertimbangkan tali ideal yang
tidak memiliki volume, hanya berupa kurva matematis.
Kita mengandaikan tegangan pada tali adalah konstanta $T$ dalam Newton.
%Jika kita mengambil ruas kecil dan melihat tegangan pada titik-titik ujungnya, kita
%melihat bahwa gaya ini tangensial dan kita akan mengandaikan bahwa besarnya
%sama pada kedua titik ujung.  Oleh karena itu, besarnya
%konstan di mana-mana dan akan kita
%sebut $T$.
Dengan mengandaikan simpangannya kecil,
kita dapat menggunakan hukum kedua Newton (kita lewati penurunannya) untuk memperoleh persamaan
\begin{equation*}
T y'' + \rho \omega^2 y = 0 .
\end{equation*}
Untuk memeriksa satuannya, perhatikan bahwa satuan $y''$ adalah $\unitfrac{m}{m^2}$, karena turunannya
terhadap $x$.

Misalkan $L$ adalah panjang tali (dalam meter) dan tali tersebut
terikat di kedua titik ujung.
Jadi, $y(0) = 0$ dan $y(L) = 0$.  Lihat
\figurevref{bvp:whirstringfig}.

\begin{myfig}
\capstart
\inputpdft{bvp-whirstring}{%
Diagram bidang xy dengan tali yang digambar sebagai kurva cekung ke bawah,
terikat pada sumbu x di x sama dengan 0 dan x sama dengan L, serta panah melingkar
yang menunjukkan rotasi mengelilingi sumbu x.  Panah lain menunjukkan bahwa
jarak tali dari sumbu x di setiap titik merupakan fungsi y dari x.}
\caption{Tali yang berputar.\label{bvp:whirstringfig}}
\end{myfig}

Kita menulis ulang persamaan tersebut sebagai
$y'' + \frac{\rho \omega^2}{T} y = 0$.
Susunannya serupa dengan \examplevref{bvp:eig1ex}, kecuali
panjang intervalnya $L$, bukan $\pi$.  Kita mencari nilai eigen
dari $y'' + \lambda y = 0, y(0) = 0, y(L) = 0$ dengan
$\lambda = \frac{\rho \omega^2}{T}$.  Seperti sebelumnya,
tidak ada nilai eigen tak positif.  Dengan $\lambda > 0$,
solusi umum persamaan tersebut adalah
$y = A \cos \bigl(  \sqrt{\lambda} \,x \bigr)
+ B \sin \bigl( \sqrt{\lambda} \,x \bigr)$.
Syarat $y(0) = 0$ mengakibatkan $A = 0$ seperti
sebelumnya.  Syarat $y(L) = 0$ mengakibatkan
$\sin \bigl( \sqrt{\lambda} \, L\bigr) = 0$ dan dengan demikian
$\sqrt{\lambda} \, L = k \pi$ untuk suatu bilangan bulat $k > 0$, sehingga
\begin{equation*}
\frac{\rho \omega^2}{T} = \lambda = \frac{k^2 \pi^2}{L^2} .
\end{equation*}

Apa makna hasil ini bagi bentuk tali?  Hasil tersebut mengatakan bahwa untuk
semua parameter $\rho$, $\omega$, $T$ yang tidak memenuhi persamaan di atas, tali
berada dalam posisi kesetimbangan, $y=0$.  Ketika
$\frac{\rho \omega^2}{T} = \frac{k^2 \pi^2}{L^2}$, tali akan
\myquote{mencuat keluar} sejauh $B$.  Kita tidak dapat menghitung $B$
dengan informasi yang kita miliki.

Mari kita andaikan bahwa $\rho$ dan $T$ tetap dan kita mengubah $\omega$.
Untuk sebagian besar nilai $\omega$, tali berada dalam keadaan kesetimbangan.  Ketika
kecepatan sudut $\omega$ mencapai nilai
$\omega = \frac{k \pi \sqrt{T}}{L\sqrt{\rho}}$, tali
mencuat keluar dan berbentuk gelombang sinus yang memotong
sumbu $x$ sebanyak $k-1$ kali di antara kedua titik ujung.
Sebagai contoh, pada $k=1$, tali tidak memotong sumbu $x$
dan bentuknya tampak seperti pada \figurevref{bvp:whirstringfig}.
Di sisi lain, ketika $k=3$, tali memotong sumbu $x$
2 kali; lihat \figurevref{bvp:whirstring2fig}.
Ketika $\omega$ berubah lagi, tali kembali ke
posisi kesetimbangan.  Semakin tinggi kecepatan sudutnya,
semakin sering tali memotong sumbu $x$ ketika mencuat keluar.

\begin{myfig}
\capstart
\inputpdft{bvp-whirstring2}{%
Diagram bidang xy dengan tali yang digambar sebagai kurva sinusoidal
yang terikat pada sumbu x di x sama dengan 0 dan x sama dengan L, tetapi juga melewati
sumbu x dua kali lagi antara 0 dan L.  Sebuah panah melingkar
menunjukkan rotasi mengelilingi sumbu x.}
\caption{Tali berputar pada nilai eigen ketiga ($k=3$).\label{bvp:whirstring2fig}}
\end{myfig}

Sebagai contoh lain, jika Anda memiliki tali lompat yang berputar (maka $k=1$ karena tali
sepenuhnya \myquote{mencuat keluar}) dan Anda
menarik ujung-ujungnya untuk meningkatkan tegangan, maka kecepatannya juga meningkat
agar tali tetap \myquote{mencuat keluar}.


\subsection{Latihan}

Petunjuk untuk latihan-latihan berikut: Perhatikan bahwa
jika $\lambda > 0$, maka
$\cos \bigl( \sqrt{\lambda}\, (t - a) \bigr)$
dan $\sin  \bigl( \sqrt{\lambda}\, (t - a) \bigr)$
juga merupakan solusi dari persamaan homogen
tersebut.

\begin{exercise}
Hitunglah semua
nilai eigen dan fungsi eigen dari
$x'' + \lambda x = 0$, $x(a) = 0$, $x(b) = 0$ (andaikan $a < b$).
\end{exercise}

\begin{exercise}
Hitunglah semua
nilai eigen dan fungsi eigen dari
$x'' + \lambda x = 0$, $x'(a) = 0$, $x'(b) = 0$ (andaikan $a < b$).
\end{exercise}

\begin{exercise}
Hitunglah semua
nilai eigen dan fungsi eigen dari
$x'' + \lambda x = 0$, $x'(a) = 0$, $x(b) = 0$ (andaikan $a < b$).
\end{exercise}

\begin{exercise}
Hitunglah semua
nilai eigen dan fungsi eigen dari
$x'' + \lambda x = 0$, $x(a) = x(b)$, $x'(a) = x'(b)$ (andaikan $a < b$).
\end{exercise}

\begin{exercise}
Kita melewati kasus $\lambda < 0$ untuk
masalah nilai batas
$x'' + \lambda x = 0$, $x(-\pi) = x(\pi)$, $x'(-\pi) = x'(\pi)$.
Selesaikan perhitungannya dan tunjukkan bahwa tidak ada nilai eigen negatif.
\end{exercise}

\setcounter{exercise}{100}

\begin{exercise}
Perhatikan tali berputar dengan panjang 2, kerapatan linear 0.1, dan tegangan 3.
Carilah kecepatan sudut terkecil ketika tali mencuat keluar.
\end{exercise}
\exsol{%
$\omega = \pi \sqrt{\frac{15}{2}}$
}

\begin{exercise}
Andaikan $x'' + \lambda x = 0$ dan $x(0)=1$, $x(1) = 1$.
Carilah semua $\lambda$ yang menghasilkan lebih dari
satu solusi.  Carilah pula solusi yang bersesuaian (hanya untuk
nilai eigennya).
\end{exercise}
\exsol{%
$\lambda_k = 4 k^2 \pi^2$ untuk $k = 1,2,3,\ldots$
\quad
$x_k =  \cos (2k\pi t) + B \sin (2k\pi t)$ \quad (untuk sembarang $B$)
}

\begin{exercise}
Andaikan $x'' + x = 0$ dan $x(0)=0$, $x'(\pi) = 1$.
Carilah semua solusi jika ada.
\end{exercise}
\exsol{%
$x(t) = - \sin(t)$
}

\begin{exercise}
Perhatikan
$x' + \lambda x = 0$ dan $x(0)=0$, $x(1) = 0$.  Mengapa persamaan tersebut tidak
memiliki nilai eigen?  Mengapa setiap persamaan orde pertama dengan dua syarat titik
ujung seperti di atas tidak memiliki nilai eigen?
\end{exercise}
\exsol{%
Solusi umumnya adalah $x = C e^{-\lambda t}$.  Karena $x(0) = 0$, maka $C=0$, sehingga $x(t) = 0$.
Oleh karena itu,
solusinya selalu identik nol.  Satu syarat selalu
cukup untuk menjamin solusi tunggal bagi persamaan orde pertama.
}

\begin{exercise}[menantang]
Andaikan $x''' + \lambda x = 0$ dan $x(0)=0$, $x'(0) = 0$, $x(1) = 0$.
Andaikan bahwa $\lambda > 0$.  Carilah suatu persamaan yang harus dipenuhi oleh semua
nilai eigen tersebut.
Petunjuk: Perhatikan bahwa $-\sqrt[3]{\lambda}$ adalah akar
dari $r^3+\lambda = 0$.
\end{exercise}
\exsol{%
$\frac{\sqrt{3}}{3} e^{\frac{-3}{2}\sqrt[3]{\lambda}}
- \frac{\sqrt{3}}{3} \cos \bigl( \frac{\sqrt{3}\, \sqrt[3]{\lambda}}{2} \bigr)
+ \sin \bigl( \frac{\sqrt{3}\, \sqrt[3]{\lambda}}{2}\bigr) = 0$
}

%%%%%%%%%%%%%%%%%%%%%%%%%%%%%%%%%%%%%%%%%%%%%%%%%%%%%%%%%%%%%%%%%%%%%%%%%%%%%%

\sectionnewpage
\section{Deret trigonometri} \label{ts:section}

%mbxINTROSUBSECTION

\sectionnotes{2 kuliah\EPref{, \S9.1 dalam \cite{EP}}\BDref{,
\S10.2 dalam \cite{BD}}}

\subsection{Fungsi periodik dan motivasi}

Sebagai motivasi untuk mempelajari deret Fourier, perhatikan masalah
\begin{equation} \label{ts:deq}
x'' + \omega_0^2 x = f(t) ,
\end{equation}
untuk suatu fungsi periodik $f(t)$.
Dalam \sectionref{forcedo:section}, kita menemukan solusi umum untuk
\begin{equation} \label{ts:deqcos}
x'' + \omega_0^2 x = F_0 \cos ( \omega t) .
\end{equation}
Salah satu cara menyelesaikan \eqref{ts:deq} adalah dengan
menguraikan $f(t)$ sebagai jumlah kosinus (dan sinus), lalu
menyelesaikan banyak masalah berbentuk \eqref{ts:deqcos}.  Selanjutnya kita menggunakan
prinsip superposisi untuk menjumlahkan semua solusi ini
guna memperoleh solusi untuk \eqref{ts:deq}.

Sebelum melanjutkan, mari kita membahas fungsi periodik
sedikit lebih terperinci.
Suatu fungsi disebut \emph{\myindex{periodik}} dengan periode $P$ jika
$f(t) = f(t+P)$ untuk semua $t$.  Untuk ringkasnya, kita mengatakan $f(t)$ bersifat $P$-periodik.
Perhatikan bahwa fungsi $P$-periodik juga bersifat $2P$-periodik, $3P$-periodik,
dan seterusnya.
Sebagai contoh, $\cos (t)$ dan $\sin (t)$ bersifat
$2\pi$-periodik.  Demikian pula $\cos (kt)$ dan $\sin (kt)$ untuk semua bilangan bulat $k$.  Fungsi
konstan merupakan contoh ekstrem.  Fungsi tersebut periodik untuk sembarang periode
(latihan).

Biasanya, kita memulai dengan fungsi $f(t)$ yang didefinisikan pada suatu interval $[-L,L]$,
dan kita ingin
\emph{memperluas $f(t)$ secara periodik}\index{memperluas secara periodik}\index{perluasan periodik}
agar menjadi
fungsi $2L$-periodik.  Kita melakukan perluasan ini
dengan mendefinisikan fungsi baru $F(t)$
sedemikian sehingga untuk $t$ dalam $[-L,L]$, $F(t) = f(t)$.  Untuk $t$ dalam $[L,3L]$,
kita definisikan $F(t) = f(t-2L)$; untuk $t$ dalam $[-3L,-L]$, $F(t) = f(t+2L)$; dan
seterusnya.
Agar cara ini berfungsi, kita memerlukan $f(-L) = f(L)$.
Kita juga dapat memulai dengan $f$
yang hanya didefinisikan pada interval setengah terbuka $(-L,L]$, kemudian mendefinisikan $f(-L) = f(L)$.

\begin{example}
Definisikan $f(t) = 1-t^2$ pada $[-1,1]$.  Perluas $f(t)$ secara periodik menjadi
fungsi 2-periodik.
Untuk $1 \leq t \leq 3$, kita memperoleh
$f(t) = 1-{(t-2)}^2$.
Untuk $-3 \leq t \leq -1$, kita memperoleh
$f(t) = 1-{(t+2)}^2$.
Untuk $3 \leq t \leq 5$, kita memperoleh
$f(t) = 1-{(t-4)}^2$.  Dan seterusnya.
Lihat \figurevref{ts:perextofinvertedparabolafig}.
\begin{myfig}
\capstart
\diffyincludegraphics{width=3in}{width=4.5in}{ts-perextofinvertedparabola}{%
Grafik suatu fungsi antara t sama dengan minus 3 dan t sama dengan 3
yang tersusun atas tiga parabola terbuka ke bawah.
Parabola pertama dimulai pada sumbu t di t sama dengan minus 3, naik, mencapai
nilai 1, lalu turun kembali ke sumbu t di t sama dengan minus 1.  Kemudian sebuah
parabola identik melakukan hal yang sama dan berakhir di t sama dengan 1 pada sumbu t, lalu
parabola identik lainnya berakhir di t sama dengan 3.}
\caption{Perluasan periodik fungsi
$1-t^2$.\label{ts:perextofinvertedparabolafig}}
\end{myfig}
\end{example}

Anda harus berhati-hati membedakan $f(t)$ dari perluasannya.  Kesalahan yang umum
adalah mengandaikan bahwa rumus untuk $f(t)$ berlaku bagi perluasannya.  Hal ini
dapat membingungkan ketika rumus $f(t)$ bersifat periodik, tetapi mungkin dengan
periode yang berbeda.

\begin{exercise}
Definisikan $f(t) = \cos t$ pada $[\nicefrac{-\pi}{2},\nicefrac{\pi}{2}]$.  Ambil perluasan
$\pi$-periodiknya dan gambarkan grafiknya.  Bagaimana perbandingannya dengan grafik
$\cos t$?
\end{exercise}

\subsection{Hasil kali dalam dan penguraian vektor eigen}

Andaikan $A$ adalah \emph{\myindex{matriks simetris}},
yaitu $A^T = A$.  Seperti yang telah kita nyatakan sebelumnya,
vektor-vektor eigen $A$ kemudian saling ortogonal.  Di sini kata
\emph{ortogonal}\index{ortogonal!vektor} berarti
bahwa jika $\vec{v}$ dan $\vec{w}$ adalah dua
vektor eigen $A$ untuk nilai eigen yang berbeda,
maka $\langle \vec{v} , \vec{w} \rangle = 0$.
Dalam kasus ini, hasil kali dalam $\langle \vec{v} , \vec{w} \rangle$
adalah \emph{\myindex{hasil kali titik}},
yang dapat dihitung sebagai $\vec{v}^T\vec{w}$.

Untuk menguraikan vektor $\vec{v}$ dalam suku-suku vektor yang saling ortogonal,
$\vec{w}_1$ dan $\vec{w}_2$, kita menulis
\begin{equation*}
\vec{v} = a_1 \vec{w}_1  + a_2 \vec{w}_2 .
\end{equation*}
Mari kita cari rumus untuk $a_1$ dan $a_2$.  Kita hitung,
\begin{equation*}
\langle \vec{v} , \vec{w_1} \rangle
=
\langle a_1 \vec{w}_1  + a_2 \vec{w}_2 , \vec{w_1} \rangle
=
a_1 \langle \vec{w}_1 , \vec{w_1} \rangle
+
a_2 \underbrace{\langle \vec{w}_2 , \vec{w_1} \rangle}_{=0}
=
a_1 \langle \vec{w}_1 , \vec{w_1} \rangle .
\end{equation*}
Oleh karena itu,
\begin{equation*}
a_1 = 
\frac{\langle \vec{v} , \vec{w_1} \rangle}{
\langle \vec{w}_1 , \vec{w_1} \rangle} .
\end{equation*}
Demikian pula,
\begin{equation*}
a_2 = 
\frac{\langle \vec{v} , \vec{w_2} \rangle}{
\langle \vec{w}_2 , \vec{w_2} \rangle} .
\end{equation*}
Anda mungkin mengingat rumus ini dari kalkulus vektor.

\begin{example}
Tuliskan
$\vec{v} = \left[ \begin{smallmatrix} 2 \\ 3 \end{smallmatrix} \right]$
sebagai kombinasi linear dari
$\vec{w_1} = \left[ \begin{smallmatrix} 1 \\ -1 \end{smallmatrix} \right]$
dan
$\vec{w_2} = \left[ \begin{smallmatrix} 1 \\ 1 \end{smallmatrix} \right]$.

Perhatikan bahwa $\vec{w}_1$ dan $\vec{w}_2$ saling ortogonal
karena $\langle \vec{w}_1 , \vec{w}_2 \rangle = 1(1) + (-1)1 = 0$.
Maka
\begin{align*}
& a_1 = 
\frac{\langle \vec{v} , \vec{w_1} \rangle}{
\langle \vec{w}_1 , \vec{w_1} \rangle}
=
\frac{2(1) + 3(-1)}{1(1) + (-1)(-1)} = \frac{-1}{2} ,
\\
& a_2 = 
\frac{\langle \vec{v} , \vec{w_2} \rangle}{
\langle \vec{w}_2 , \vec{w_2} \rangle}
=
\frac{2 + 3}{1 + 1} = \frac{5}{2} .
\end{align*}
Jadi,
\begin{equation*}
\begin{bmatrix} 2 \\ 3 \end{bmatrix}
=
\frac{-1}{2}
\begin{bmatrix} 1 \\ -1 \end{bmatrix}
+
\frac{5}{2}
\begin{bmatrix} 1 \\ 1 \end{bmatrix} .
\end{equation*}
\end{example}

\subsection{Deret trigonometri}

Alih-alih menguraikan vektor dalam suku-suku vektor eigen suatu matriks,
kita menguraikan fungsi dalam suku-suku fungsi eigen dari suatu
masalah nilai eigen tertentu.  Masalah nilai eigen yang kita gunakan untuk
deret Fourier adalah
\begin{equation*}
x'' + \lambda x = 0, \quad x(-\pi) = x(\pi), \quad x'(-\pi) = x'(\pi) .
\end{equation*}
Kita telah menghitung bahwa fungsi-fungsi eigennya adalah 1, $\cos (k t)$,
$\sin (k t)$.  Artinya, kita ingin menemukan representasi suatu
fungsi $2\pi$-periodik $f(t)$ sebagai
\begin{equation*}
\mybxbg{~~
\begin{aligned}
f(t) &=
\frac{a_0}{2} +
\sum_{n=1}^\infty
\Bigl[ a_n \cos (n t) + b_n \sin (n t) \Bigr]
\\
& = 
\frac{a_0}{2}
+ a_1 \cos (t) + b_1 \sin (t) 
+ a_2 \cos (2 t) + b_2 \sin (2 t) 
%+ a_3 \cos (3 t) + b_3 \sin (3 t) 
+ \cdots
\end{aligned}
~~}
\end{equation*}
Deret ini disebut \emph{\myindex{deret Fourier}}%
\footnote{Dinamai menurut matematikawan Prancis
\href{https://en.wikipedia.org/wiki/Joseph_Fourier}{Jean Baptiste Joseph Fourier}
(1768--1830).} atau
\emph{\myindex{deret trigonometri}} untuk $f(t)$.
Suku
$a_n \cos (n t) + b_n \sin (n t)$ terkadang disebut
\emph{\myindex{harmonik}} ke-$n^{\text{}}$.
Untuk kemudahan, kita menulis koefisien fungsi eigen 1 sebagai $\frac{a_0}{2}$.
Kita juga dapat memandang $1 = \cos (0t)$, sehingga
kita hanya perlu memperhatikan $\cos (kt)$ dan $\sin (kt)$.

Seperti untuk matriks, kita ingin menemukan \emph{\myindex{proyeksi}}
$f(t)$ pada subruang-subruang yang diberikan oleh fungsi-fungsi eigen.  Jadi kita ingin
mendefinisikan \emph{\myindex{hasil kali dalam fungsi}}.  Sebagai contoh, untuk
menemukan $a_n$,
kita ingin menghitung $\langle \, f(t) \, , \, \cos (nt) \, \rangle$.
Kita mendefinisikan hasil kali dalam sebagai
\begin{equation*}
\langle \, f(t)\, , \, g(t) \, \rangle \overset{\text{def}}{=}
\int_{-\pi}^\pi f(t) \, g(t) \, dt .
\end{equation*}
Dengan hasil kali dalam ini,
kita melihat pada bagian sebelumnya bahwa fungsi-fungsi eigen $\cos (kt)$
(termasuk fungsi eigen konstan), dan
$\sin (kt)$ bersifat \emph{ortogonal\index{ortogonal!fungsi}}, yaitu
\begin{align*}
\langle \, \cos (mt)\, , \, \cos (nt) \, \rangle = 0 & \qquad \text{untuk } m \neq n , \\
\langle \, \sin (mt)\, , \, \sin (nt) \, \rangle = 0 & \qquad \text{untuk } m \neq n , \\
\langle \, \sin (mt)\, , \, \cos (nt) \, \rangle = 0 & \qquad \text{untuk semua } m \text{ dan } n .
\end{align*}
Untuk $n=1,2,3,\ldots$,
kita memiliki
\begin{align*}
\langle \, \cos (nt) \, , \, \cos (nt) \, \rangle &=
\int_{-\pi}^\pi \cos(nt)\cos(nt) \, dt
=
\pi,
\\
\langle \, \sin (nt) \, , \, \sin (nt) \, \rangle &=
\int_{-\pi}^\pi \sin(nt)\sin(nt) \, dt
=
\pi,
\end{align*}
menurut kalkulus elementer.  Untuk fungsi konstannya, kita memperoleh
\begin{equation*}
\langle \, 1 \, , \, 1 \, \rangle
=
\int_{-\pi}^\pi 1 \cdot 1 \, dt
 = 2\pi.
\end{equation*}
Koefisien-koefisiennya diberikan oleh
\begin{equation*}
\mybxbg{~~
\begin{aligned}
& a_n =
\frac{\langle \, f(t) \, , \, \cos (nt) \, \rangle}{\langle \, \cos (nt) \, , \,
\cos (nt) \, \rangle}
= 
\frac{1}{\pi} \int_{-\pi}^\pi f(t) \cos (nt) \, dt , \\
& b_n =
\frac{\langle \, f(t) \, , \, \sin (nt) \, \rangle}{\langle \, \sin (nt) \, , \,
\sin (nt) \, \rangle}
= 
\frac{1}{\pi} \int_{-\pi}^\pi f(t) \sin (nt) \, dt .
\end{aligned}
~~}
\end{equation*}
Bandingkan ekspresi-ekspresi ini dengan contoh berdimensi hingga.
Untuk $a_0$, kita memperoleh rumus serupa
\begin{equation*}
\mybxbg{~~
a_0 = 2
\frac{\langle \, f(t) \, , \, 1 \, \rangle}{\langle \, 1 \, , \,
1 \, \rangle}
=
\frac{1}{\pi} \int_{-\pi}^\pi f(t) \, dt .
~~}
\end{equation*}

Mari kita periksa rumus-rumus tersebut melalui sifat keortogonalan.  Andaikan untuk
sementara bahwa
\begin{equation*}
f(t) = \frac{a_0}{2} + \sum_{n=1}^\infty a_n \cos (n t) + b_n
\sin (n t) .
\end{equation*}
Maka untuk $m \geq 1$, kita memiliki
\begin{equation*}
\begin{split}
\langle \, f(t)\,,\,\cos (mt) \, \rangle
& =
\Bigl\langle \, \frac{a_0}{2} + \sum_{n=1}^\infty
\Bigl[
a_n \cos (n t) + b_n \sin (n t)
\Bigr]
\,,\, \cos (mt) \, \Bigr\rangle \\
& =
\frac{a_0}{2}
\langle \, 1 \, , \, \cos (mt) \, \rangle
%mbxlatex \\
%mbxlatex & \phantom{=abcd}
+
\sum_{n=1}^\infty
\Bigl[
a_n \langle \, \cos (nt) \, , \, \cos (mt) \, \rangle +
b_n \langle \, \sin (n t) \, , \, \cos (mt) \, \rangle
\Bigr]
\\
& =
a_m \langle \, \cos (mt) \, , \, \cos (mt) \, \rangle .
\end{split}
\end{equation*}
Jadi,
$a_m =
\frac{\langle \, f(t) \, , \, \cos (mt) \, \rangle}{\langle \, \cos (mt) \, , \,
\cos (mt) \, \rangle}$.

\begin{exercise}
Lakukan perhitungan untuk $a_0$ dan $b_m$.
\end{exercise}

\begin{example}
Ambil fungsi
\begin{equation*}
f(t) = t
\end{equation*}
untuk $t$ dalam $(-\pi,\pi]$.  Perluas $f(t)$ secara periodik dan tuliskan
sebagai deret Fourier.  Fungsi ini disebut
\emph{\myindex{gelombang gigi gergaji}}, dan memiliki banyak penerapan,
misalnya dalam musik elektronik.

\begin{myfig}
\capstart
\diffyincludegraphics{width=3in}{width=4.5in}{ts-sawtooth}{%
Untuk t antara minus 2 pi dan 2 pi,
ditunjukkan grafik fungsi yang berupa garis diagonal berkemiringan 1
melalui titik asal antara minus pi dan pi.  Untuk t kurang dari minus pi,
grafiknya berupa garis berkemiringan 1 melalui titik (minus 2 pi,0),
dan untuk t lebih dari pi,
grafiknya berupa garis berkemiringan 1 melalui (2 pi,0).}
\caption{Grafik fungsi gelombang gigi gergaji.\label{ts:sawtoothfig}}
\end{myfig}

Grafik fungsi periodik yang telah diperluas diberikan dalam
\figurevref{ts:sawtoothfig}.
Mari kita hitung koefisien-koefisiennya.  Kita mulai dengan $a_0$,
\begin{equation*}
a_0 = \frac{1}{\pi} \int_{-\pi}^\pi t \,dt = 0 .
\end{equation*}
Kita akan sering menggunakan hasil kalkulus yang menyatakan bahwa integral fungsi
ganjil pada interval simetris bernilai nol.  Ingat bahwa
\emph{\myindex{fungsi ganjil}} adalah
fungsi $\varphi(t)$ sedemikian sehingga $\varphi(-t) = -\varphi(t)$.
Sebagai contoh,
fungsi $t$, $\sin t$, atau (yang penting bagi kita)
$t \cos (nt)$ semuanya merupakan fungsi ganjil.  Jadi
\begin{equation*}
a_n = \frac{1}{\pi} \int_{-\pi}^\pi t \cos (nt) \,dt = 0 .
\end{equation*}
Kita beralih ke $b_n$.  Fakta berguna lain dari kalkulus
adalah bahwa integral fungsi genap pada
interval simetris sama dengan
dua kali integral fungsi yang sama pada setengah interval tersebut.
Suatu \emph{\myindex{fungsi genap}}
adalah
fungsi $\varphi(t)$ sedemikian sehingga $\varphi(-t) = \varphi(t)$.  Sebagai contoh,
$t \sin (nt)$ bersifat genap.
\begin{equation*}
\begin{split}
b_n & = \frac{1}{\pi} \int_{-\pi}^\pi t \sin (nt) \,dt \\
& = \frac{2}{\pi} \int_{0}^\pi t \sin (nt) \,dt \\
& = \frac{2}{\pi} \left(
\left[ \frac{-t \cos (nt)}{n} \right]_{t=0}^{\pi}
+
\frac{1}{n}
\int_{0}^\pi \cos (nt) \,dt
\right)
\\
& = \frac{2}{\pi} \left(
\frac{-\pi \cos (n\pi)}{n}
+
0
\right) \\
& =  \frac{-2 \cos (n\pi)}{n}
=  \frac{2 \,{(-1)}^{n+1}}{n} .
\end{split}
\end{equation*}
Kita telah menggunakan fakta bahwa
\begin{equation*}
\cos (n\pi) = {(-1)}^n =
\begin{cases}
1 & \text{jika } n \text{ genap} , \\
-1 & \text{jika } n \text{ ganjil} .
\end{cases}
\end{equation*}
Oleh karena itu, deretnya adalah
\begin{equation*}
\sum_{n=1}^\infty
\frac{2 \,{(-1)}^{n+1}}{n} \,
\sin (n t) .
\end{equation*}

Secara lebih eksplisit,
tiga harmonik pertama dari deret untuk $f(t)$ adalah
\begin{equation*}
2 \, \sin (t)
- \sin (2t)
+\frac{2}{3} \sin (3t)
+ \cdots
\end{equation*}
Grafik tiga suku pertama deret ini, bersama dengan grafik
20 suku pertamanya, diberikan dalam
\figurevref{ts:sawtoothfsfig}.

\begin{myfig}
\capstart
%original files ts-sawtooth-fs3 ts-sawtooth-fs20
\diffyincludegraphics{width=6.24in}{width=9in}{ts-sawtooth-fs3-fs20}{%
Dua grafik fungsi gelombang gigi gergaji pada interval minus 2 pi hingga 2 pi.
Grafik sebelah kiri ditumpangi tiga harmonik pertama
(tiga suku tak nol pertama deret Fourier) dari deret Fourier,
yaitu garis berombak yang kira-kira mengikuti gelombang gigi gergaji sepanjang ruas
lurusnya, lalu turun cepat ke ruas lurus berikutnya.
Dalam versi ini hanya terdapat sedikit riak.  Di sebelah kanan,
grafik yang ditumpangkan adalah grafik 20 harmonik pertama.  Pada grafik ini,
garisnya beriak jauh lebih cepat, tetapi juga lebih dekat dengan ruas-ruas lurus.
Terlihat bahwa riaknya sedikit lebih menjauh dari gelombang gigi gergaji
di dekat titik-titik diskontinuitas.}
\caption{Tiga harmonik pertama (grafik kiri) dan 20 harmonik pertama (grafik kanan) dari fungsi
gelombang gigi gergaji.\label{ts:sawtoothfsfig}}
\end{myfig}
\end{example}

\begin{example}
Ambil fungsi
\begin{equation*}
f(t) =
\begin{cases}
0 & \text{jika } \;{-\pi} < t \leq 0 , \\
\pi & \text{jika } \;\phantom{-}0 < t \leq \pi .
\end{cases}
\end{equation*}
\nopagebreak[4]%
Perluas $f(t)$ secara periodik dan tuliskan
sebagai deret Fourier.  Fungsi ini atau variasinya sering muncul
dalam penerapan dan disebut
\emph{\myindex{gelombang persegi}}.  Fungsi ini merupakan sinyal
yang dihasilkan hanya dengan menyalakan dan mematikan sakelar secara periodik.

\begin{myfig}
\capstart
\diffyincludegraphics{width=3in}{width=4.5in}{ts-squarewave}{%
Grafik fungsi gelombang persegi periodik pada interval
minus 2 pi hingga 2 pi.
Fungsi tersebut bernilai pi pada interval minus 2 pi hingga minus pi,
lalu bernilai 0 pada interval minus pi hingga 0, kemudian bernilai pi dari 0 hingga pi,
lalu bernilai nol dari pi hingga 2 pi.}
\caption{Grafik fungsi gelombang persegi.\label{ts:squarewavefig}}
\end{myfig}

\pagebreak[2]
Grafik fungsi periodik yang telah diperluas diberikan dalam
\figurevref{ts:squarewavefig}.
Sekarang kita menghitung koefisien-koefisiennya.  Kita mulai dengan $a_0$
\begin{equation*}
a_0 = \frac{1}{\pi} \int_{-\pi}^\pi f(t) \,dt
= \frac{1}{\pi} \int_{0}^\pi \pi \,dt = \pi .
\end{equation*}
Selanjutnya,
\begin{equation*}
a_n = \frac{1}{\pi} \int_{-\pi}^\pi f(t) \cos (nt) \,dt 
= \frac{1}{\pi} \int_{0}^\pi \pi \cos (nt) \,dt = 0 .
\end{equation*}
Dan terakhir,
\begin{equation*}
\begin{split}
b_n & = \frac{1}{\pi} \int_{-\pi}^\pi f(t) \sin (nt) \,dt \\
& = \frac{1}{\pi} \int_{0}^\pi \pi \sin (nt) \,dt \\
& = \left[ \frac{- \cos (nt)}{n} \right]_{t=0}^\pi \\
& = \frac{1 - \cos (\pi n)}{n}
= \frac{1 - {(-1)}^n}{n}
=
\begin{cases}
\frac{2}{n} & \text{jika } n \text{ ganjil} , \\
0 & \text{jika } n \text{ genap} .
\end{cases}
\end{split}
\end{equation*}
Deret Fourier-nya adalah
\begin{equation*}
\frac{\pi}{2} +  \sum_{\substack{n=1\\n \text{ ganjil}}}^\infty
\frac{2}{n} 
\sin (n t)
=
\frac{\pi}{2} + \sum_{k=1}^\infty
\frac{2}{2k-1} 
\sin \bigl( (2k-1)\, t \bigr) .
\end{equation*}

Tiga harmonik pertama dari deret untuk $f(t)$ adalah
\begin{equation*}
\frac{\pi}{2}
+
2 \, \sin (t)
+
\frac{2}{3}  \sin (3t)
+ \cdots
\end{equation*}
Grafik ketiga suku pertama ini dan juga 20 suku pertama deret
diberikan dalam
\figurevref{ts:squarewavefsfig}.

\begin{myfig}
\capstart
%original files ts-squarewave-fs3 ts-squarewave-fs20
\diffyincludegraphics{width=6.24in}{width=9in}{ts-squarewave-fs3-fs20}{%
Dua grafik fungsi gelombang persegi pada interval minus 2 pi hingga 2 pi.
Grafik sebelah kiri ditumpangi tiga harmonik pertama deret Fourier,
yaitu garis berombak yang kira-kira mengikuti gelombang persegi sepanjang ruas-ruas
horizontal lurusnya, lalu bergerak cepat turun atau naik ke ruas lurus berikutnya.
Dalam versi ini hanya terdapat sedikit riak.  Di sebelah kanan,
grafik yang ditumpangkan adalah grafik 20 harmonik pertama.  Pada grafik ini,
garisnya beriak jauh lebih cepat, tetapi juga lebih dekat dengan ruas-ruas lurus.
Terlihat bahwa riaknya sedikit lebih menjauh dari gelombang persegi
di dekat titik-titik diskontinuitas.}
\caption{Tiga harmonik pertama (grafik kiri) dan 20 harmonik pertama (grafik kanan) dari
fungsi gelombang persegi.\label{ts:squarewavefsfig}}
\end{myfig}
\end{example}

Sejauh ini kita menghindari persoalan konvergensi.  Sebagai contoh,
jika $f(t)$ adalah fungsi gelombang persegi,
persamaan
\begin{equation*}
f(t) = 
\frac{\pi}{2} + \sum_{k=1}^\infty
\frac{2}{2k-1} 
\sin \bigl( (2k-1)\, t \bigr) .
\end{equation*}
hanya merupakan kesamaan untuk nilai $t$ yang membuat $f(t)$ kontinu.
Kita tidak memperoleh kesamaan untuk $t=-\pi,0,\pi$ dan semua titik diskontinuitas
lain dari $f(t)$.  Tidak sulit untuk melihat bahwa ketika $t$ merupakan kelipatan bulat dari
$\pi$ (yang memberikan semua titik diskontinuitas), maka
$\sin \bigl( (2k-1)\, t \bigr) = 0$, sehingga
\begin{equation*}
\frac{\pi}{2} + \sum_{k=1}^\infty
\frac{2}{2k-1} 
\sin \bigl( (2k-1)\, t \bigr) = \frac{\pi}{2} .
\end{equation*}
Kita mendefinisikan ulang $f(t)$ pada $[-\pi,\pi]$ sebagai
\begin{equation*}
f(t) =
\begin{cases}
0 & \text{jika } \; {-\pi} < t < 0 , \\
\pi & \text{jika } \; \phantom{-}0 < t < \pi , \\
\nicefrac{\pi}{2} & \text{jika } \; \phantom{-}t = -\pi, 
t = 0,\text{ atau }
t = \pi,
\end{cases}
\end{equation*}
dan memperluasnya secara periodik.
Deret tersebut sama dengan fungsi $f(t)$ baru yang telah diperluas ini di mana-mana, termasuk di
titik-titik diskontinuitas.
Secara umum kita tidak akan mempermasalahkan perubahan nilai fungsi
di beberapa (berhingga banyak) titik.

Kita akan membahas lebih lanjut tentang konvergensi pada bagian berikutnya.
Secara kasar,
jika $f(t)$ merupakan fungsi yang cukup baik, maka deretnya
konvergen ke $f(t)$ di setiap titik tempat $f(t)$ kontinu.  Namun, mari kita
sebutkan secara singkat suatu efek diskontinuitas.
Kita memperbesar tampilan di dekat titik diskontinuitas gelombang persegi
dan menggambar 100 harmonik pertama.  Lihat
\figurevref{ts:squarewavegibbsfig}.  Meskipun
deret tersebut merupakan aproksimasi yang sangat baik jauh dari titik-titik diskontinuitas, galat
(lewatannya)
di dekat titik diskontinuitas $t=\pi$ tampaknya tidak semakin kecil
ketika kita mengambil semakin banyak harmonik.
Perilaku ini dikenal sebagai \emph{\myindex{fenomena Gibbs}}.
Namun, daerah tempat galatnya besar memang mengecil ketika
kita mengambil lebih banyak suku deret.

\begin{myfig}
\capstart
\diffyincludegraphics{width=3in}{width=4.5in}{ts-squarewave-gibbs}{%
Tampilan dekat fenomena Gibbs di sekitar diskontinuitas fungsi gelombang
persegi.  Yang ditampilkan adalah tampilan dekat di sekitar satu sisi salah satu
titik diskontinuitas; khususnya, kita dapat melihat ruas tempat fungsi bernilai
pi sesaat hingga t sama dengan pi, sedangkan ruas tempat fungsi bernilai
0 kini berada di luar gambar.  Di atas grafik gelombang persegi
ditumpangkan grafik 100 harmonik pertama deret Fourier, berupa
garis yang sangat berombak dan cukup dekat dengan grafik gelombang persegi hingga
kita semakin mendekati t sama dengan pi dari sisi kiri.  Ketika semakin dekat, riaknya
semakin besar dan melewati fungsi gelombang persegi dengan
jarak yang berarti.  Di t sama dengan pi, grafik turun untuk menyusul tempat
gelombang persegi bernilai 0.}
\caption{Fenomena Gibbs sedang terjadi.\label{ts:squarewavegibbsfig}}
\end{myfig}

\medskip

Kita dapat memandang fungsi periodik sebagai \myquote{sinyal}
yang merupakan superposisi dari
banyak sinyal berfrekuensi murni.
Sebagai contoh, kita dapat memandang gelombang
persegi sebagai nada dengan frekuensi dasar tertentu
yang keluar melalui pengeras suara komputer Anda.
Frekuensi dasar ini disebut
\emph{\myindex{frekuensi fundamental}}
dan merupakan \myquote{nada musik} yang Anda kenali sebagai suara tersebut.
Namun, gelombang persegi juga merupakan
superposisi dari
banyak nada murni berbeda dengan frekuensi yang merupakan kelipatan
frekuensi fundamental.
Dalam musik,
frekuensi-frekuensi yang lebih tinggi disebut \emph{\myindex{nada atas}}.
Himpunan semua frekuensi yang muncul disebut
\emph{\myindex{spektrum}} sinyal tersebut.
Di sisi lain,
jika sinyalnya berupa
gelombang sinus sederhana, bukan gelombang persegi,
maka sinyal tersebut hanya terdiri atas nada murni dengan
frekuensi fundamental (tanpa nada atas).
Cara paling sederhana untuk membuat suara dengan komputer adalah menggunakan gelombang persegi,
dan suaranya sangat berbeda dari nada murni.
Jika Anda pernah memainkan permainan video
dari sekitar tahun 1980-an, Anda telah mendengar bunyi gelombang persegi.

\subsection{Latihan}

\begin{exercise}
Andaikan $f(t)$ didefinisikan pada $[-\pi,\pi]$ sebagai $\sin (5t) + \cos (3t)$.  Perluas
secara periodik dan hitunglah deret Fourier dari $f(t)$.
\end{exercise}

\begin{exercise}
Andaikan $f(t)$ didefinisikan pada $[-\pi,\pi]$ sebagai $\lvert t \rvert$.
Perluas secara periodik dan hitunglah deret Fourier dari $f(t)$.
\end{exercise}

\begin{exercise}
Andaikan $f(t)$ didefinisikan pada $[-\pi,\pi]$ sebagai $\lvert t \rvert^3$.
Perluas secara periodik dan hitunglah deret Fourier dari $f(t)$.
\end{exercise}

\begin{exercise}
Andaikan $f(t)$ didefinisikan pada $(-\pi,\pi]$ sebagai
\begin{equation*}
f(t) =
\begin{cases}
-1 & \text{jika } \; {-\pi} < t \leq 0 , \\
1 & \text{jika } \; \phantom{-}0 < t \leq \pi .
\end{cases}
\end{equation*}
Perluas secara periodik dan hitunglah deret Fourier dari $f(t)$.
\end{exercise}

\begin{exercise}
Andaikan $f(t)$ didefinisikan pada $(-\pi,\pi]$ sebagai $t^3$.
Perluas secara periodik dan hitunglah deret Fourier dari $f(t)$.
\end{exercise}

\begin{exercise}
Andaikan $f(t)$ didefinisikan pada $[-\pi,\pi]$ sebagai $t^2$.
Perluas secara periodik dan hitunglah deret Fourier dari $f(t)$.
\end{exercise}

Terdapat bentuk lain dari deret Fourier yang menggunakan eksponensial kompleks
$e^{nt}$ untuk $n=\ldots,-2,-1,0,1,2,\ldots$, bukan
$\cos(nt)$ dan $\sin(nt)$ untuk $n$ positif.  Bentuk ini terkadang lebih mudah
digunakan.  Bentuk ini tentu lebih ringkas untuk ditulis,
dan hanya ada satu rumus untuk koefisiennya.  Kekurangannya,
koefisien-koefisien tersebut merupakan bilangan kompleks.

\begin{exercise}
\begin{samepage}
Misalkan
\begin{equation*}
f(t) = \frac{a_0}{2} + \sum_{n=1}^\infty
\Bigl[ a_n \cos (n t) + b_n \sin (n t) \Bigr] .
\end{equation*}
Gunakan rumus Euler $e^{i\theta} = \cos (\theta) + i \sin (\theta)$ untuk
menunjukkan bahwa terdapat bilangan kompleks $c_m$ sedemikian sehingga
\begin{equation*}
f(t) = 
\sum_{m=-\infty}^\infty c_m e^{imt} .
\end{equation*}
Perhatikan bahwa penjumlahan sekarang mencakup semua bilangan bulat, termasuk yang negatif.
Jangan khawatirkan konvergensi dalam perhitungan ini.
Petunjuk: Mungkin lebih baik memulai dari bentuk eksponensial kompleks dan menulis
deret tersebut sebagai
\begin{equation*}
c_0 + \sum_{m=1}^\infty \Bigl[ c_m e^{imt} + c_{-m} e^{-imt} \Bigr] .
\end{equation*}
\end{samepage}
\end{exercise}

\setcounter{exercise}{100}

\begin{exercise}
Andaikan $f(t)$ didefinisikan pada $[-\pi,\pi]$ sebagai $f(t) = \sin(t)$.  Perluas
secara periodik dan hitunglah deret Fourier-nya.
\end{exercise}
\exsol{%
$\sin(t)$
}

\begin{exercise}
Andaikan $f(t)$ didefinisikan pada $(-\pi,\pi]$ sebagai $f(t) = \sin(\pi t)$.  Perluas
secara periodik dan hitunglah deret Fourier-nya.
\end{exercise}
\exsol{%
%$\sum\limits_{n=1}^\infty
%\frac{(\pi-n) \sin( \pi n+{\pi}^{2})
%+(\pi+n)\sin(\pi n-{\pi}^{2}) }{\pi {n}^{2}-{\pi}^{3}}
%\sin(nt)$
%Ini menyederhana menjadi:
$\sum\limits_{n=1}^\infty
\frac{2n {(-1)}^{n} \sin({\pi}^{2})}{\pi({\pi}^{2} - {n}^{2})}
\sin(nt)$
%$a_0 = \frac{1}{\pi} \int_{-\pi}^\pi \sin(\pi t) \, dt$
%\\
%$a_n =
%\frac{1}{\pi} \int_{-\pi}^\pi \sin(\pi t) \cos (nt) \, dt$
%\\
%$b_n =
%\frac{1}{\pi} \int_{-\pi}^\pi f(t) \sin (nt) \, dt$
}

\begin{exercise}
Andaikan $f(t)$ didefinisikan pada $(-\pi,\pi]$ sebagai $f(t) = \sin^2(t)$.
Perluas
secara periodik dan hitunglah deret Fourier-nya.
\end{exercise}
\exsol{%
$\frac{1}{2}-\frac{1}{2}\cos(2t)$
}

\begin{exercise}
Andaikan $f(t)$ didefinisikan pada $(-\pi,\pi]$ sebagai $f(t) = t^4$.
Perluas secara periodik dan hitunglah deret Fourier-nya.
\end{exercise}
\exsol{%
$\frac{\pi^4}{5} + \sum\limits_{n=1}^\infty
\frac{{(-1)}^{n} (8{\pi}^{2}{n}^{2}-48) }{{n}^{4}}
\cos(nt)$
}

%%%%%%%%%%%%%%%%%%%%%%%%%%%%%%%%%%%%%%%%%%%%%%%%%%%%%%%%%%%%%%%%%%%%%%%%%%%%%%

\sectionnewpage
\section{Lebih lanjut tentang deret Fourier}
\label{moreonfourier:section}

%mbxINTROSUBSECTION

\sectionnotes{2 kuliah\EPref{, \S9.2--\S9.3 dalam \cite{EP}}\BDref{,
\S10.3 dalam \cite{BD}}}

%Sebelum membaca kuliah ini, mungkin baik untuk terlebih dahulu mencoba
%Proyek IV (deret Fourier)\index{perangkat lunak IODE!Proyek IV} dari
%situs web IODE: \url{http://www.math.uiuc.edu/iode/}.  Setelah membaca
%kuliah ini, mungkin baik untuk melanjutkan dengan
%Proyek V (deret Fourier lagi)\index{perangkat lunak IODE!Proyek V}.

\subsection{Fungsi $2L$-periodik}

Kita telah menghitung deret Fourier untuk fungsi $2\pi$-periodik, tetapi bagaimana
dengan fungsi yang periodenya berbeda?  Tidak perlu khawatir, perhitungannya merupakan
kasus sederhana dari perubahan variabel.  Kita hanya menskalakan ulang sumbu
bebas.  Perhatikan fungsi $2L$-periodik $f(t)$.  Besaran $L$ disebut
\emph{\myindex{setengah periode}}.  Misalkan $s = \frac{\pi}{L}  t$.
Maka fungsi
\begin{equation*}
g(s) = f\left(\frac{L}{\pi} s \right)
\end{equation*}
bersifat $2\pi$-periodik, dan kita mengetahui apa yang harus dilakukan dengannya.
Kita juga harus menskalakan ulang semua sinus dan kosinus kita.
Dalam deret tersebut, kita menggunakan $\frac{\pi}{L} t$ sebagai variabel.  Artinya, kita
ingin menulis
\begin{equation*}
\mybxbg{~~
f(t) = 
\frac{a_0}{2} +
\sum_{n=1}^\infty
\Bigl[
a_n \cos \Bigl( \frac{n \pi}{L} t \Bigr)
+ b_n \sin \Bigl(\frac{n \pi}{L} t \Bigr)
\Bigr] .
~~}
\end{equation*}
Jika kita mengubah variabel menjadi $s$, kita melihat bahwa
\begin{equation*}
g(s) = 
\frac{a_0}{2} +
\sum_{n=1}^\infty
\Bigl[
a_n \cos (n s)
+ b_n \sin (n s)
\Bigr] .
\end{equation*}
Kita menghitung $a_n$ dan $b_n$ seperti sebelumnya.  Setelah menuliskan
integral-integralnya, kita mengubah variabel dari $s$ kembali ke $t$, dengan juga memperhatikan
bahwa $ds = \frac{\pi}{L} \, dt$.
\begin{equation*}
\mybxbg{~~
\begin{aligned}
& a_0 =
\frac{1}{\pi}
\int_{-\pi}^\pi
g(s) \, ds
=
\frac{1}{L}
\int_{-L}^L
f(t) \, dt , \\
& a_n =
\frac{1}{\pi}
\int_{-\pi}^\pi
g(s) \, \cos (n s) \, ds
=
\frac{1}{L}
\int_{-L}^L
f(t) \, \cos \left( \frac{n \pi}{L} t \right) \, dt , \\
& b_n =
\frac{1}{\pi}
\int_{-\pi}^\pi
g(s) \, \sin (n s) \, ds
=
\frac{1}{L}
\int_{-L}^L
f(t) \, \sin \left( \frac{n \pi}{L} t \right) \, dt .
\end{aligned}
~~}
\end{equation*}

Dua setengah periode yang paling umum muncul dalam contoh
adalah $\pi$ dan 1 karena rumusnya sederhana.  Perlu ditekankan bahwa kita
tidak melakukan matematika baru---kita hanya mengubah variabel.
Jika Anda memahami
deret Fourier untuk fungsi $2\pi$-periodik, Anda memahaminya untuk
fungsi $2L$-periodik.  Anda dapat memandangnya sebagai sekadar penggunaan
satuan waktu yang berbeda.  Yang kita lakukan hanyalah memindahkan beberapa konstanta,
tetapi seluruh matematikanya tetap sama.

\begin{example}
Misalkan
\begin{equation*}
f(t) =
\lvert t \rvert
\qquad \text{untuk } \; {-1} < t \leq 1,
\end{equation*}
yang diperluas secara periodik.  Grafik
perluasan periodiknya diberikan dalam \figurevref{gfs:sawcontfig}.
Hitunglah deret Fourier dari $f(t)$.

\begin{myfig}
\capstart
\diffyincludegraphics{width=3in}{width=4.5in}{gfs-sawcont}{%
Grafik fungsi 2-periodik pada interval minus 2 hingga 2, berupa
garis dari (minus 2,0) ke (minus 1,1), lalu dari (minus 1,1) ke (0,0), kemudian
dari (0,0) ke (1,1), dan selanjutnya dari (1,1) ke (2,0).  Garis tersebut kontinu,
tetapi memiliki sudut-sudut tajam.}
\caption{Perluasan periodik fungsi $f(t)$.\label{gfs:sawcontfig}}
\end{myfig}

Pertama, kita mengenali bahwa $f$ bersifat $2$-periodik sehingga $L=1$.
Kita ingin
menulis $f(t) = \frac{a_0}{2} + \sum_{n=1}^\infty a_n \cos (n \pi t) + b_n
\sin (n \pi t)$.  Kita mulai dengan $a_n$ untuk $n \geq 1$.
Kita perhatikan bahwa $\lvert t \rvert \cos (n \pi t)$
bersifat genap dan untuk $0 \leq t \leq 1$, $f(t) = \lvert t \rvert = t$.  Jadi,
\begin{equation*}
\begin{split}
a_n & = \int_{-1}^1 f(t) \cos (n \pi t) \, dt \\
& = 2 \int_{0}^1 t \cos (n \pi t) \, dt \\
 & = 2 \left[ \frac{t}{n \pi} \sin (n \pi t) \right]_{t=0}^1 -
2 \int_{0}^1 \frac{1}{n \pi} \sin (n \pi t) \, dt \\
& =  0 + \frac{2}{n^2 \pi^2} \Bigl[ \cos (n \pi t) \Bigr]_{t=0}^1
 =  \frac{2 \bigl( {(-1)}^n -1 \bigr) }{n^2 \pi^2}
=
\begin{cases}
0 & \text{jika } n \text{ genap} , \\
\frac{-4 }{n^2 \pi^2} & \text{jika } n \text{ ganjil}  .
\end{cases}
\end{split}
\end{equation*}
Selanjutnya kita mencari $a_0$:
\begin{equation*}
a_0 = \int_{-1}^1 \lvert t \rvert \, dt 
=
1 .
\end{equation*}
Anda seharusnya dapat menemukan integral ini dengan memandang integral tersebut
sebagai luas di bawah grafik tanpa melakukan perhitungan sama sekali.
Terakhir, kita mencari $b_n$.  Perhatikan bahwa
$\lvert t \rvert \sin (n \pi t)$ bersifat ganjil, sehingga
\begin{equation*}
b_n = \int_{-1}^1 f(t) \sin (n \pi t) \, dt = 0 .
\end{equation*}
Jadi,
deretnya adalah
\begin{equation*}
\frac{1}{2} + 
\sum_{\substack{n=1 \\ n \text{ ganjil}}}^\infty \frac{-4}{n^2 \pi^2} \cos (n \pi t) .
\end{equation*}

Beberapa suku pertama deret tersebut
hingga harmonik ke-$3^{\text{}}$ adalah
\begin{equation*}
\frac{1}{2} -
\frac{4}{\pi^2} \cos (\pi t)
-
\frac{4}{9 \pi^2} \cos (3 \pi t)
- \cdots
\end{equation*}
Grafik beberapa suku ini dan juga grafik hingga harmonik
ke-${20}^{\text{}}$ diberikan dalam
\figurevref{gfs:sawcontfsfig}.  Anda seharusnya memperhatikan betapa dekatnya grafik tersebut
dengan fungsi sebenarnya.  Anda juga seharusnya memperhatikan bahwa tidak terdapat
\myquote{fenomena Gibbs} karena tidak ada diskontinuitas.

\begin{myfig}
\capstart
%original files gfs-sawcontfs3 gfs-sawcont-fs20
\diffyincludegraphics{width=6.24in}{width=9in}{gfs-sawcont-fs3-fs20}{%
Grafik tiga dan 20 harmonik pertama dari deret Fourier
untuk f dari t.  Terlihat bahwa garis di sebelah kiri (tiga harmonik pertama)
sedikit membulat, tetapi mengikuti bentuknya dengan baik. Di sebelah kanan,
grafik 20 harmonik pertama hampir tidak dapat dibedakan dari
grafik f dari t pada pandangan pertama, kecuali grafiknya sedikit membulat
di sudut-sudut.}
\caption{Deret Fourier dari $f(t)$ hingga harmonik ke-$3^{\text{}}$ (grafik
kiri)
dan hingga harmonik ke-${20}^{\text{}}$ (grafik kanan).\label{gfs:sawcontfsfig}}
\end{myfig}
\end{example}

\subsection{Konvergensi}

Kita akan memerlukan limit satu sisi dari fungsi.
Kita akan menggunakan notasi berikut
\begin{equation*}
f(c-) = \lim_{t \uparrow c} f(t),
\qquad \text{dan} \qquad
f(c+) = \lim_{t \downarrow c} f(t).
\end{equation*}
Jika Anda tidak terbiasa dengan notasi ini,
$\lim_{t \uparrow c} f(t)$ berarti kita mengambil limit $f(t)$
ketika $t$ mendekati $c$ dari bawah (yaitu\ $t < c$), dan
$\lim_{t \downarrow c} f(t)$ berarti kita mengambil limit $f(t)$
ketika $t$ mendekati $c$ dari atas (yaitu\ $t > c$).
Sebagai contoh, untuk fungsi gelombang persegi
\begin{equation} \label{gfs:sqwaveeq}
f(t) =
\begin{cases}
0 & \text{jika } \; {-\pi} < t \leq 0 , \\
\pi & \text{jika } \; \phantom{-}0 < t \leq \pi ,
\end{cases}
\end{equation}
kita memiliki $f(0-) = 0$ dan $f(0+) = \pi$.

Misalkan $f(t)$ adalah fungsi yang didefinisikan pada interval $[a,b]$.  Andaikan
kita menemukan berhingga banyak titik
$a=t_0$, $t_1$, $t_2$, \ldots, $t_k=b$ dalam
interval tersebut, sedemikian sehingga $f(t)$ kontinu
pada interval-interval
$(t_0,t_1)$,
$(t_1,t_2)$, \ldots,
$(t_{k-1},t_k)$.
Andaikan pula bahwa semua limit satu sisinya ada, yaitu
semua nilai
$f(t_0+)$,
$f(t_1-)$,
$f(t_1+)$,
$f(t_2-)$,
$f(t_2+)$,
\ldots,
$f(t_k-)$
ada dan berhingga.
Maka
kita mengatakan bahwa $f(t)$ \emph{\myindex{kontinu sepotong-sepotong}}.

Jika, lebih lanjut, $f(t)$ terdiferensialkan kecuali di berhingga banyak titik,
dan $f'(t)$ kontinu sepotong-sepotong, maka
$f(t)$ disebut \emph{\myindex{mulus sepotong-sepotong}}.

\begin{example}
Fungsi gelombang persegi, yaitu \eqref{gfs:sqwaveeq} yang diperluas secara periodik,
mulus sepotong-sepotong pada $[-\pi,\pi]$ atau interval berhingga lainnya, sehingga
kita cukup mengatakan bahwa $f(t)$
mulus sepotong-sepotong tanpa menyebutkan interval.
\end{example}

\begin{example}
Fungsi $f(t) = \lvert t \lvert$
mulus sepotong-sepotong.
\end{example}

\begin{example}
Fungsi $f(t) = \frac{1}{t}$ tidak mulus sepotong-sepotong pada
$[-1,1]$ (atau interval lain yang memuat nol).  Bahkan, fungsi tersebut tidak
kontinu sepotong-sepotong.
\end{example}

\begin{example}
Fungsi $f(t) = \sqrt[3]{t}$ tidak mulus sepotong-sepotong pada
$[-1,1]$ (atau interval lain yang memuat nol).
Fungsi $f(t)$ kontinu, tetapi
turunannya $f'(t) = \frac{1}{3t^{2/3}}$ tak terbatas di dekat nol sehingga tidak kontinu
sepotong-sepotong.
\end{example}

Fungsi mulus sepotong-sepotong memberikan jawaban sederhana mengenai konvergensi
deret Fourier.

\begin{theorem}
Andaikan $f(t)$ adalah fungsi $2L$-periodik yang mulus sepotong-sepotong.
Misalkan
\begin{equation*}
\frac{a_0}{2} + \sum_{n=1}^\infty
\Bigl[
a_n \cos \Bigl( \frac{n \pi}{L} t \Bigr)
+ b_n \sin \Bigl( \frac{n \pi}{L} t \Bigr)
\Bigr]
\end{equation*}
adalah deret Fourier untuk $f(t)$.  Maka deret tersebut konvergen
untuk semua $t$.  Jika $f(t)$ kontinu
di $t$, maka
\begin{equation*}
f(t) = \frac{a_0}{2} + \sum_{n=1}^\infty
\Bigl[
a_n \cos \Bigl( \frac{n \pi}{L} t \Bigr)
+ b_n \sin \Bigl( \frac{n \pi}{L} t \Bigr)
\Bigr] .
\end{equation*}
Jika tidak,
\begin{equation*}
\frac{f(t-)+f(t+)}{2} =
\frac{a_0}{2} + \sum_{n=1}^\infty
\Bigl[
a_n \cos \Bigl( \frac{n \pi}{L}  t \Bigr)
+ b_n \sin \Bigl( \frac{n \pi}{L} t \Bigr)
\Bigr] .
\end{equation*}
\end{theorem}

Jika kebetulan kita memiliki
$f(t) = \frac{f(t-)+f(t+)}{2}$ di semua titik diskontinuitas, deret Fourier
konvergen ke $f(t)$ di mana-mana.  Kita selalu dapat mendefinisikan ulang $f(t)$
dengan mengubah nilai di setiap titik diskontinuitas secara sesuai.  Kemudian kita dapat menulis
tanda sama dengan antara $f(t)$ dan deretnya tanpa kekhawatiran.
Kita telah menyebutkan fakta ini
secara singkat pada akhir bagian sebelumnya.

Teorema tersebut tidak menyatakan seberapa cepat deretnya konvergen.
Ingat kembali pembahasan fenomena Gibbs pada bagian sebelumnya.
Semakin dekat Anda ke titik diskontinuitas, semakin banyak suku yang perlu diambil
untuk memperoleh aproksimasi fungsi yang akurat.

\subsection{Diferensiasi dan integrasi deret Fourier}

Deret Fourier tidak hanya konvergen dengan baik, tetapi juga mudah didiferensialkan
dan diintegralkan.  Kita dapat melakukannya hanya dengan mendiferensialkan atau
mengintegralkan suku demi suku.

\begin{theorem}
Andaikan
\begin{equation*}
f(t) = \frac{a_0}{2} +
\sum_{n=1}^\infty
\Bigl[
a_n \cos \Bigl( \frac{n \pi}{L} t \Bigr)
+ b_n \sin \Bigl( \frac{n \pi}{L} t \Bigr)
\Bigr]
\end{equation*}
adalah fungsi kontinu yang mulus sepotong-sepotong dan turunannya $f'(t)$
mulus sepotong-sepotong.  Maka turunannya dapat
diperoleh dengan mendiferensialkan suku demi suku,
\begin{equation*}
f'(t) = \sum_{n=1}^\infty
\Bigl[
\frac{-a_n n \pi}{L} 
\sin \Bigl( \frac{n \pi}{L} t \Bigr)
+ \frac{b_n n \pi}{L} \cos \Bigl( \frac{n \pi}{L} t \Bigr)
\Bigr] .
\end{equation*}
\end{theorem}

Penting bahwa fungsinya kontinu.  Fungsi tersebut boleh memiliki sudut, tetapi tidak boleh memiliki
lompatan.  Jika tidak, deret yang didiferensialkan akan gagal konvergen.  Sebagai
latihan, ambil deret yang diperoleh untuk gelombang persegi dan cobalah
mendiferensialkan deret tersebut.
Serupa dengan diferensiasi,
integrasi deret Fourier juga dilakukan suku demi suku.

\begin{theorem}
Andaikan
\begin{equation*}
f(t) = \frac{a_0}{2} + \sum_{n=1}^\infty
\Bigl[
a_n \cos \Bigl( \frac{n \pi}{L} t \Bigr)
+ b_n \sin \Bigl( \frac{n \pi}{L} t \Bigr)
\Bigr]
\end{equation*}
adalah fungsi yang mulus sepotong-sepotong.  Maka antiturunannya
diperoleh dengan mengantidiferensialkan suku demi suku, sehingga
\begin{equation*}
F(t) = \frac{a_0 t}{2} + C + \sum_{n=1}^\infty
\Bigl[
\frac{a_n L}{n \pi} \sin \Bigl( \frac{n \pi}{L} t \Bigr)
+ \frac{-b_n L}{n \pi}  \cos \Bigl( \frac{n \pi}{L} t \Bigr)
\Bigr] ,
\end{equation*}
dengan $F'(t) = f(t)$ dan $C$ suatu konstanta sembarang.
\end{theorem}

Perhatikan bahwa deret untuk $F(t)$ bukan lagi deret Fourier karena memuat
suku $\frac{a_0 t}{2}$.  Antiturunan
fungsi periodik tidak harus periodik, sehingga kita tidak seharusnya
mengharapkan deret Fourier.  Kecuali, tentu saja, jika $a_0 = 0$.

\subsection{Laju konvergensi dan kemulusan}

Kita mempertimbangkan contoh fungsi periodik yang dapat didiferensialkan satu kali, tetapi tidak
dua kali.

\begin{example}
Ambil fungsi
\begin{equation*}
f(t) =
\begin{cases}
(t+1)\,t & \text{jika } \; {-1} < t \leq 0 , \\
(1-t)\,t & \text{jika } \; \phantom{-}0 < t \leq 1 ,
\end{cases}
\end{equation*}
dan perluas menjadi
fungsi 2-periodik.
Turunan $f'(t)$ ada untuk semua $t$, tetapi $f''(t)$
tidak ada jika $t$ merupakan bilangan bulat.
Khususnya,
$f'(t) = 2t+1$ jika $-1 \leq t \leq 0$ dan $f'(t) = 1-2t$ jika
$0 \leq t \leq 1$, sehingga $f'(t)$ memiliki \myquote{sudut} di $0$.
Lihat
\figurevref{gfs:smoothexfig}
untuk grafik $f(t)$ dan $f'(t)$.


\begin{myfig}
\capstart
\diffyincludegraphics{width=3in}{width=4.5in}{gfs-smoothex}{%
Grafik fungsi f dari t yang tampak kurang lebih sinusoidal, tetapi tersusun
atas ruas-ruas parabola yang disambungkan sedemikian rupa sehingga
turunan pada titik penyambungan tersebut terdefinisi.  Grafik turunannya
ditumpangkan, dan turunannya hanya tersusun atas garis-garis lurus serta
memiliki sudut (tetapi kontinu).  Lebih tepatnya, untuk t dari minus 2 hingga 2,
fungsi f dari t mula-mula
naik dan melengkung ke bawah kembali ke sumbu t, lalu turun dan melengkung
naik kembali ke sumbu t, dan seterusnya.  Turunannya dimulai dari 1 ketika
t sama dengan minus 2, lalu turun sebagai garis lurus hingga mencapai
minus 1 ketika t sama dengan minus 1, kemudian
naik kembali sebagai garis lurus hingga mencapai 1 di t sama dengan 0, dan seterusnya.}
\caption{Fungsi 2-periodik yang mulus (garis mulus) beserta turunannya yang tidak mulus
(garis bergerigi).\label{gfs:smoothexfig}}
\end{myfig}


\begin{exercise}
Hitunglah $f''(0+)$ dan $f''(0-)$.
\end{exercise}

Mari kita hitung koefisien deret Fourier.  Perhitungan lengkapnya
melibatkan beberapa integrasi parsial dan diserahkan kepada pembaca.
\begin{align*}
a_0 & = 
\int_{-1}^1
f(t) \, dt = 
\int_{-1}^0
(t+1)\,t \, dt +
\int_0^1
(1-t)\,t \, dt = 0 , \\
a_n & = 
\int_{-1}^1
f(t) \, \cos (n\pi t) \, dt = 
\int_{-1}^0
(t+1)\,t
\, \cos (n \pi t) \, dt +
\int_0^1
(1-t)\,t
\, \cos (n \pi t) \, dt = 0, \\
b_n & = 
\int_{-1}^1
f(t) \, \sin (n\pi t) \, dt = 
\int_{-1}^0
(t+1)\,t
\, \sin (n \pi t) \, dt +
\int_0^1
(1-t)\,t
\, \sin (n \pi t) \, dt \\
& =
\frac{4 ( 1-{(-1)}^n)}{\pi^3 n^3} 
=
\begin{cases}
\frac{8}{\pi^3 n^3} & \text{jika } n \text{ ganjil} , \\
0 & \text{jika } n \text{ genap} .
\end{cases}
\end{align*}
Artinya, deretnya adalah
\begin{equation*}
\sum_{\substack{n=1 \\ n \text{ ganjil}}}^\infty \frac{8}{\pi^3 n^3} \sin (n \pi t) .
\end{equation*}

Deret ini konvergen dengan sangat cepat.
Jika kita menggambar hingga
harmonik ketiga, yaitu
\begin{equation*}
\frac{8}{\pi^3} \sin (\pi t) + 
\frac{8}{27 \pi^3} \sin (3 \pi t) ,
\end{equation*}
grafiknya hampir tidak dapat dibedakan dari grafik $f(t)$ dalam
\figurevref{gfs:smoothexfig}.
Bahkan, koefisien
$\frac{8}{27 \pi^3}$ sudah hanya sebesar $0.0096$ (kira-kira).
Alasan perilaku ini adalah suku $n^3$ pada penyebut.
Koefisien $b_n$ menuju nol secepat
$\nicefrac{1}{n^3}$ menuju
nol.
\end{example}

Untuk fungsi yang disusun sepotong-sepotong dari polinomial seperti di atas,
secara umum berlaku bahwa
jika Anda memiliki satu turunan tetapi tidak memiliki dua turunan,
koefisien Fourier
akan menuju nol kira-kira seperti $\nicefrac{1}{n^3}$.
Jika Anda
hanya memiliki fungsi kontinu yang tidak terdiferensialkan,
maka koefisien Fourier akan menuju
nol seperti $\nicefrac{1}{n^2}$.  Jika terdapat diskontinuitas, maka
koefisien Fourier akan menuju nol kira-kira seperti $\nicefrac{1}{n}$.
Untuk fungsi yang lebih umum, keadaannya agak lebih rumit, tetapi
gagasan yang sama tetap berlaku: Semakin banyak turunan yang dimiliki, semakin cepat koefisiennya
menuju nol.  Penalaran serupa berlaku dalam arah sebaliknya.  Jika koefisiennya menuju
nol seperti $\nicefrac{1}{n^2}$ (atau lebih cepat),
Anda selalu memperoleh fungsi kontinu.  Jika
koefisiennya menuju nol seperti $\nicefrac{1}{n^3}$ (atau lebih cepat),
Anda memperoleh fungsi yang terdiferensialkan
di mana-mana.
%Oleh karena itu, kita dapat mengetahui banyak hal tentang kemulusan suatu fungsi dengan melihat
%koefisien Fourier-nya.

Untuk membenarkan perilaku ini, ambil sebagai contoh fungsi yang didefinisikan oleh
deret Fourier
\begin{equation*}
f(t) = \sum_{n=1}^\infty \frac{1}{n^3} \sin (n t) .
\end{equation*}
Ketika kita mendiferensialkan suku demi suku, kita memperoleh
\begin{equation*}
f'(t) = \sum_{n=1}^\infty \frac{1}{n^2} \cos (n t) .
\end{equation*}
Oleh karena itu, koefisiennya sekarang turun seperti $\nicefrac{1}{n^2}$, yang
berarti bahwa kita memiliki fungsi kontinu.
Turunan
dari $f'(t)$ terdefinisi
di sebagian besar titik, tetapi terdapat titik-titik tempat $f'(t)$ tidak terdiferensialkan.
Fungsi tersebut memiliki sudut, tetapi tidak memiliki lompatan.
Jika kita
mendiferensialkan lagi (di tempat yang memungkinkan), kita memperoleh bahwa fungsi
$f''(t)$
sekarang tidak kontinu (memiliki lompatan)
\begin{equation*}
f''(t) = \sum_{n=1}^\infty \frac{-1}{n} \sin (n t) .
\end{equation*}
Fungsi ini serupa dengan gelombang gigi gergaji.  Jika kita mencoba mendiferensialkan
deret tersebut sekali lagi, kita akan memperoleh
\begin{equation*}
\sum_{n=1}^\infty -\cos (n t) ,
\end{equation*}
yang tidak konvergen!

\begin{exercise}
Gunakan komputer untuk menggambar deret yang kita peroleh untuk $f(t)$, $f'(t)$, dan
$f''(t)$.  Artinya, gambarkan, misalnya, lima harmonik pertama dari fungsi-fungsi tersebut.  Di
titik mana saja $f''(t)$ memiliki diskontinuitas?
\end{exercise}

\subsection{Latihan}

\begin{exercise}
Misalkan
\begin{equation*}
f(t) =
\begin{cases}
0 & \text{jika } \; {-1} < t \leq 0 , \\
t & \text{jika } \; \phantom{-}0 < t \leq  1 ,
\end{cases}
\end{equation*}
yang diperluas secara periodik.
\begin{tasks}
\task Hitunglah deret Fourier untuk $f(t)$.
\task Tuliskan deret tersebut secara eksplisit hingga harmonik ke-$3^{\text{}}$.
\end{tasks}
\end{exercise}

\begin{exercise}
Misalkan
\begin{equation*}
f(t) =
\begin{cases}
-t & \text{jika } \; {-1} < t \leq 0 , \\
t^2 & \text{jika } \; \phantom{-}0 < t \leq  1 ,
\end{cases}
\end{equation*}
yang diperluas secara periodik.
\begin{tasks}
\task Hitunglah deret Fourier untuk $f(t)$.
\task Tuliskan deret tersebut secara eksplisit hingga harmonik ke-$3^{\text{}}$.
\end{tasks}
\end{exercise}

\begin{exercise}
Misalkan
\begin{equation*}
f(t) =
\begin{cases}
\frac{-t}{10} & \text{jika } \; {-10} < t \leq 0 , \\
\frac{t}{10} & \text{jika } \; \phantom{-1}0 < t \leq  10 ,
\end{cases}
\end{equation*}
yang diperluas secara periodik (periodenya 20).
\begin{tasks}
\task Hitunglah deret Fourier untuk $f(t)$.
\task Tuliskan deret tersebut secara eksplisit hingga harmonik ke-$3^{\text{}}$.
\end{tasks}
\end{exercise}

\begin{exercise}
Misalkan $f(t) = \sum_{n=1}^\infty \frac{1}{n^3} \cos (n t)$.  Apakah $f(t)$
kontinu dan terdiferensialkan di mana-mana?  Carilah turunannya (jika turunannya ada
di mana-mana),
atau berikan alasan mengapa $f(t)$ tidak terdiferensialkan di mana-mana.
\end{exercise}

\begin{exercise}
Misalkan $f(t) = \sum_{n=1}^\infty \frac{{(-1)}^n}{n} \sin (n t)$.  Apakah $f(t)$
terdiferensialkan di mana-mana?  Carilah turunannya (jika turunannya ada di mana-mana), atau
berikan alasan mengapa $f(t)$ tidak terdiferensialkan di mana-mana.
\end{exercise}

\begin{exercise}
Misalkan
\begin{equation*}
f(t) =
\begin{cases}
0 & \text{jika } \; {-2} < t \leq 0, \\
t & \text{jika } \; \phantom{-}0 < t \leq 1, \\
-t+2 & \text{jika } \; \phantom{-}1 < t \leq 2,
\end{cases}
\end{equation*}
yang diperluas secara periodik.
\begin{tasks}
\task Hitunglah deret Fourier untuk $f(t)$.
\task Tuliskan deret tersebut secara eksplisit hingga harmonik ke-$3^{\text{}}$.
\end{tasks}
\end{exercise}

\begin{exercise}
Misalkan
\begin{equation*}
f(t) = e^t \qquad \text{untuk } \; {-1} < t \leq 1
\end{equation*}
yang diperluas secara periodik.
\begin{tasks}
\task Hitunglah deret Fourier untuk $f(t)$.
\task Tuliskan deret tersebut secara eksplisit hingga harmonik ke-$3^{\text{}}$.
\task Ke nilai apa deret tersebut konvergen pada $t=1$.
\end{tasks}
\end{exercise}

\begin{exercise}
\pagebreak[2]
Misalkan
\begin{equation*}
f(t) = t^2 \qquad \text{untuk } \; {-1} < t \leq 1
\end{equation*}
yang diperluas secara periodik.
\begin{tasks}
\task Hitunglah deret Fourier untuk $f(t)$.
\task Dengan memasukkan $t=0$,
hitunglah $\displaystyle \sum_{n=1}^\infty \frac{{(-1)}^n}{n^2} = -1 + \frac{1}{4} -
\frac{1}{9} + \cdots$.
\task Sekarang hitunglah $\displaystyle \sum_{n=1}^\infty \frac{1}{n^2} = 1 + \frac{1}{4} +
\frac{1}{9} + \cdots$.
\end{tasks}
\end{exercise}

\begin{exercise}
Misalkan
\begin{equation*}
f(t) =
\begin{cases}
0 & \text{jika } \; {-3} < t \leq 0, \\
t & \text{jika } \; \phantom{-}0 < t \leq 3,
\end{cases}
\end{equation*}
yang diperluas secara periodik.  Andaikan $F(t)$ adalah fungsi yang diberikan
oleh deret Fourier dari $f$.  Tanpa menghitung deret Fourier,
hitunglah
\begin{tasks}(3)
\task $F(2)$
\task $F(-2)$
\task $F(4)$
\task $F(-4)$
\task $F(3)$
\task $F(-9)$
\end{tasks}
\end{exercise}

\setcounter{exercise}{100}

\begin{exercise}
Misalkan
\begin{equation*}
f(t) = t^2 \qquad \text{untuk } \; {-2} < t \leq 2
\end{equation*}
yang diperluas secara periodik.
\begin{tasks}
\task Hitunglah deret Fourier untuk $f(t)$.
\task Tuliskan deret tersebut secara eksplisit hingga harmonik ke-$3^{\text{}}$.
\end{tasks}
\end{exercise}
\exsol{%
a) $\frac{8}{6} +
\sum\limits_{n=1}^\infty
\frac{16{(-1)}^n}{\pi^2 n^2}
\cos\bigl(\frac{n\pi}{2} t\bigr)$
\quad
b) $\frac{8}{6}
-
\frac{16}{\pi^2 }
\cos\bigl(\frac{\pi}{2} t\bigr)
+
\frac{4}{\pi^2}
\cos\bigl(\pi t\bigr)
-
\frac{16}{9\pi^2}
\cos\bigl(\frac{3\pi}{2} t\bigr) + \cdots$
}

\begin{exercise}
Misalkan
\begin{equation*}
f(t) = t \qquad \text{untuk } \; {-\lambda} < t \leq \lambda \; \text{ (untuk suatu } \lambda > 0 \text{)}
\end{equation*}
yang diperluas secara periodik.
\begin{tasks}
\task Hitunglah deret Fourier untuk $f(t)$.
\task Tuliskan deret tersebut secara eksplisit hingga harmonik ke-$3^{\text{}}$.
\end{tasks}
\end{exercise}
\exsol{%
a)
$\sum\limits_{n=1}^\infty
\frac{{(-1)}^{n+1}2\lambda}{n \pi}
\sin\bigl(\frac{n\pi}{\lambda} t\bigr)$
\quad
b)
$\frac{2\lambda}{\pi}
\sin\bigl(\frac{\pi}{\lambda} t\bigr)
-
\frac{\lambda}{\pi}
\sin\bigl(\frac{2\pi}{\lambda} t\bigr)
+
\frac{2\lambda}{3\pi}
\sin\bigl(\frac{3\pi}{\lambda} t\bigr) - \cdots$
}

\begin{exercise}
Misalkan
\begin{equation*}
f(t) = \frac{1}{2} + \sum_{n=1}^\infty
\frac{1}{n(n^2+1)}
\sin(n\pi t) .
\end{equation*}
Hitunglah $f'(t)$.
\end{exercise}
\exsol{%
$f'(t) = \sum\limits_{n=1}^\infty
\frac{\pi}{n^2+1}
\cos(n\pi t)$
}

\begin{exercise}
\pagebreak[2]
Misalkan
\begin{equation*}
f(t) = \frac{1}{2} + \sum_{n=1}^\infty
\frac{1}{n^3}
\cos(n t) .
\end{equation*}
\begin{tasks}
\task Carilah antiturunannya.
\task Apakah antiturunannya periodik?
\end{tasks}
\end{exercise}
\exsol{%
a)
$F(t) = \frac{t}{2} + C + \sum\limits_{n=1}^\infty
\frac{1}{n^4}
\sin(nt)$
\qquad
b) tidak
}

\begin{exercise}
Misalkan
\begin{equation*}
f(t) = \nicefrac{t}{2} \qquad \text{untuk } \; {-\pi} < t < \pi
\end{equation*}
yang diperluas secara periodik.
\begin{tasks}
\task Hitunglah deret Fourier untuk $f(t)$.
\task Masukkan $t=\nicefrac{\pi}{2}$ untuk menemukan representasi deret
bagi $\nicefrac{\pi}{4}$.
\task Dengan menggunakan empat suku pertama hasil dari bagian b), aproksimasikan
$\nicefrac{\pi}{4}$.
\end{tasks}
\end{exercise}
\exsol{%
a)
$\sum\limits_{n=1}^\infty
\frac{{(-1)}^{n+1}}{n} \sin(nt)$
\qquad
b) $f$ kontinu pada $t=\nicefrac{\pi}{2}$ sehingga
deret Fourier konvergen ke $f(\nicefrac{\pi}{2}) = \nicefrac{\pi}{4}$.
Diperoleh
$\nicefrac{\pi}{4} = \sum\limits_{n=1}^\infty
\frac{{(-1)}^{n+1}}{2n-1} = 1 - \nicefrac{1}{3} + \nicefrac{1}{5}-
\nicefrac{1}{7} + \cdots$.
\qquad
c) Dengan menggunakan empat suku pertama, diperoleh $\nicefrac{76}{105}\approx 0.72$ (aproksimasi yang cukup
buruk; Anda harus mengambil sekitar 50 suku agar mulai mencapai galat
kurang dari $0.01$ terhadap $\nicefrac{\pi}{4}$).
}

\begin{exercise}
Misalkan
\begin{equation*}
f(t) = 
\begin{cases}
0 & \text{jika } \; {-2} < t \leq 0, \\
2 & \text{jika } \; \phantom{-}0 < t \leq 2,
\end{cases}
\end{equation*}
yang diperluas secara periodik.  Andaikan $F(t)$ adalah fungsi yang diberikan
oleh deret Fourier dari $f$.  Tanpa menghitung deret Fourier,
hitunglah
\begin{tasks}(3)
\task $F(0)$
\task $F(-1)$
\task $F(1)$
\task $F(-2)$
\task $F(4)$
\task $F(-9)$
\end{tasks}
\end{exercise}
\exsol{%
a) $F(0) = 1$, 
b) $F(-1) = 0$, 
c) $F(1) = 2$, 
d) $F(-2) = 1$, 
e) $F(4) = 1$, 
f) $F(-9) = 0$ 
}

%%%%%%%%%%%%%%%%%%%%%%%%%%%%%%%%%%%%%%%%%%%%%%%%%%%%%%%%%%%%%%%%%%%%%%%%%%%%%%

\sectionnewpage
\section{Deret sinus dan kosinus}
\label{sec:scs}

%mbxINTROSUBSECTION

\sectionnotes{2 kuliah\EPref{, \S9.3 dalam \cite{EP}}\BDref{,
\S10.4 dalam \cite{BD}}}

\subsection{Fungsi periodik ganjil dan genap}

Mungkin sekarang Anda telah memperhatikan bahwa fungsi ganjil tidak memiliki suku kosinus dalam 
deret Fourier dan fungsi genap tidak memiliki suku sinus dalam deret Fourier.
Pengamatan ini bukan kebetulan.  Mari kita telaah fungsi periodik genap dan ganjil
secara lebih terperinci.

Ingat bahwa fungsi $f(t)$ disebut \emph{ganjil}\index{fungsi ganjil} jika $f(-t) =
-f(t)$.  Fungsi $f(t)$ disebut \emph{genap}\index{fungsi genap} jika
$f(-t) = f(t)$.  Sebagai contoh, $\cos (n t)$ genap dan $\sin (n t)$ ganjil.
Demikian pula, fungsi $t^k$ genap jika $k$ genap dan ganjil jika $k$ ganjil.

\begin{exercise}
Ambil dua fungsi $f(t)$ dan $g(t)$ dan definisikan hasil kalinya $h(t) =
f(t)g(t)$.
\begin{tasks}
\task Andaikan $f(t)$ dan $g(t)$ keduanya ganjil.  Apakah $h(t)$ ganjil atau genap?
\task Andaikan yang satu genap dan yang lain ganjil.  Apakah $h(t)$ ganjil atau genap?
\task Andaikan keduanya genap.  Apakah $h(t)$ ganjil atau genap?
\end{tasks}
\end{exercise}

Jika $f(t)$ dan $g(t)$ keduanya ganjil, maka $f(t)+g(t)$ ganjil.  Hal serupa berlaku untuk
fungsi genap.  Di sisi lain,
jika $f(t)$ ganjil dan $g(t)$ genap, kita tidak dapat mengatakan apa pun tentang
jumlah
$f(t) + g(t)$.  Bahkan, deret Fourier setiap fungsi merupakan jumlah dari
fungsi ganjil (suku-suku sinus) dan fungsi genap (suku-suku kosinus).

Dalam bagian ini, kita meninjau fungsi periodik ganjil dan genap.
Sebelumnya kita telah mendefinisikan perluasan periodik-$2L$
dari suatu fungsi yang didefinisikan pada interval $[-L,L]$.  Kadang-kadang kita hanya
berkepentingan dengan fungsi pada interval $[0,L]$, dan akan lebih mudah
jika kita memiliki fungsi ganjil (berturut-turut genap).  Jika fungsinya ganjil (berturut-turut genap),
semua suku kosinus (berturut-turut sinus) menghilang.
Yang akan kita lakukan adalah
mengambil
perluasan ganjil (berturut-turut genap) dari fungsi ke $[-L,L]$, lalu 
memperluasnya secara periodik menjadi fungsi periodik-$2L$.

Ambil fungsi $f(t)$ yang didefinisikan pada $[0,L]$.  Pada $(-L,L]$ definisikan fungsi-fungsi
\begin{align*}
F_{\text{odd}}(t) & \overset{\text{def}}{=}
\begin{cases}
f(t) & \text{jika } \; \phantom{-}0 \leq t \leq L , \\
-f(-t) & \text{jika } \; {-L} < t < 0 ,
\end{cases}
\\
F_{\text{even}}(t) & \overset{\text{def}}{=}
\begin{cases}
f(t) & \text{jika } \; \phantom{-}0 \leq t \leq L , \\
f(-t) & \text{jika } \; {-L} < t < 0 .
\end{cases}
\end{align*}
Perluas $F_{\text{odd}}(t)$ dan $F_{\text{even}}(t)$ agar periodik-$2L$.
Kemudian
$F_{\text{odd}}(t)$ disebut
\emph{\myindex{perluasan periodik ganjil}} dari $f(t)$, dan
$F_{\text{even}}(t)$ disebut
\emph{\myindex{perluasan periodik genap}} dari $f(t)$.
Untuk perluasan ganjil, umumnya kita mengandaikan bahwa $f(0) = f(L) = 0$.

\begin{exercise}
Periksa bahwa $F_{\text{odd}}(t)$ ganjil dan $F_{\text{even}}(t)$ genap.
Untuk $F_{\text{odd}}$,
andaikan $f(0) = f(L) = 0$.
\end{exercise}

\begin{example}
Ambil fungsi $f(t) = t\,(1-t)$ yang didefinisikan pada $[0,1]$. 
\figurevref{scs:oddevenextfig}
menunjukkan grafik perluasan periodik ganjil dan genap dari $f(t)$.

\begin{myfig}
\capstart
%berkas asli scs-oddext scs-evenext
\diffyincludegraphics{width=6.24in}{width=9in}{scs-ext-odd-even}{%
Dua grafik untuk t antara minus 2 dan 2, yang merupakan perluasan periodik
ganjil (di kiri) dan genap (di kanan) dari t kali besaran 1 dikurangi t
yang didefinisikan untuk t antara 0 dan 1.  Artinya, fungsi pada interval dari
0 hingga 1 merupakan parabola terbalik yang bermula dan berakhir pada sumbu t.
Dalam perluasan ganjil, grafik hanya terdiri atas parabola terbalik dan tegak
yang berselang-seling, semuanya bermula dan berakhir pada sumbu t,
membentuk pola yang menyerupai sinusoid.  Di kanan, perluasan genap
hanyalah sederet parabola terbalik identik yang disambungkan.}
\caption{Perluasan periodik-2 ganjil dan genap dari $f(t) =
t\,(1-t)$, $0 \leq t \leq 1$.\label{scs:oddevenextfig}}
\end{myfig}
\end{example}

\subsection{Deret sinus dan kosinus}

Mari kita tuliskan deret Fourier untuk
fungsi periodik-$2L$ ganjil $f(t)$.
Pertama, kita menghitung koefisien $a_n$ (termasuk
$n=0$) dan memperoleh
\begin{equation*}
a_n = \frac{1}{L} \int_{-L}^L f(t) \cos \left( \frac{n \pi}{L} t \right)
\, dt = 0 .
\end{equation*}
Artinya, deret Fourier dari fungsi ganjil tidak memiliki suku kosinus.
Integralnya nol
karena $f(t) \cos \left( \frac{n \pi}{L} t \right)$
merupakan fungsi ganjil (hasil kali fungsi ganjil dan
fungsi genap adalah ganjil) dan integral fungsi ganjil pada interval
simetris adalah nol.
Fungsi $f(t) \sin \left( \frac{n \pi}{L} t \right)$ genap
karena merupakan hasil kali dua
fungsi ganjil.
Integral fungsi genap pada interval simetris
$[-L,L]$ adalah dua kali integral fungsi tersebut pada interval $[0,L]$.
Jadi
\begin{equation*}
b_n = 
\frac{1}{L} \int_{-L}^L f(t) \sin \left( \frac{n \pi}{L} t \right) \, dt =
\frac{2}{L} \int_{0}^L f(t) \sin \left( \frac{n \pi}{L} t \right) \, dt .
\end{equation*}
Deret Fourier dari $f(t)$ ganjil kemudian adalah
\begin{equation*}
\sum_{n=1}^\infty b_n \sin \left( \frac{n \pi}{L} t \right) .
\end{equation*}

Sekarang andaikan $f(t)$ merupakan fungsi periodik-$2L$ genap.
Dengan alasan yang persis
sama seperti di atas, kita memperoleh $b_n = 0$ dan
\begin{equation*}
a_n = 
\frac{2}{L} \int_{0}^L f(t) \cos \left( \frac{n \pi}{L} t \right) \, dt .
\end{equation*}
Rumus tersebut tetap berlaku untuk $n=0$, dan dalam hal itu menjadi
\begin{equation*}
a_0 = 
\frac{2}{L} \int_{0}^L f(t) \, dt .
\end{equation*}
Deret Fourier dari $f(t)$ genap kemudian adalah
\begin{equation*}
\frac{a_0}{2}
+
\sum_{n=1}^\infty a_n \cos \left( \frac{n \pi}{L} t \right) .
\end{equation*}

Konsekuensi yang menarik ialah bahwa koefisien deret Fourier dari
fungsi ganjil (atau genap) dapat dihitung hanya dengan mengintegralkan pada
setengah interval $[0,L]$.  Oleh karena itu, kita dapat menghitung deret Fourier dari
perluasan ganjil (atau genap) suatu
fungsi dengan menghitung integral-integral tertentu pada interval
tempat fungsi aslinya didefinisikan.

\begin{theorem}
Misalkan $f(t)$ suatu fungsi mulus bagian demi bagian yang didefinisikan pada $[0,L]$.
Maka perluasan periodik ganjil
dari $f(t)$ memiliki deret Fourier
\begin{equation*}
\mybxbg{~~
F_{\text{odd}}(t) = \sum_{n=1}^\infty b_n \sin \left( \frac{n \pi}{L} t
\right) ,
~~}
\end{equation*}
dengan
\begin{equation*}
\mybxbg{~~
b_n = 
\frac{2}{L} \int_{0}^L f(t)\, \sin \left( \frac{n \pi}{L} t \right) \, dt .
~~}
\end{equation*}
Perluasan periodik genap dari $f(t)$ memiliki deret Fourier
\begin{equation*}
\mybxbg{~~
F_{\text{even}}(t) = \frac{a_0}{2} + \sum_{n=1}^\infty a_n \cos \left(
\frac{n \pi}{L} t \right) ,
~~}
\end{equation*}
dengan
\begin{equation*}
\mybxbg{~~
a_n = 
\frac{2}{L} \int_{0}^L f(t)\, \cos \left( \frac{n \pi}{L} t \right) \, dt .
~~}
\end{equation*}
\end{theorem}

Kita menyebut deret $\sum_{n=1}^\infty b_n \sin \left( \frac{n \pi}{L} t\right)$ 
sebagai \emph{\myindex{deret sinus}} dari $f(t)$ dan deret
$\frac{a_0}{2} + \sum_{n=1}^\infty a_n \cos \left( \frac{n \pi}{L} t
\right)$
sebagai \emph{\myindex{deret kosinus}} dari $f(t)$.  
Sering kali kita sebenarnya tidak memedulikan apa yang terjadi di luar $[0,L]$.
Kita cukup memilih deret mana pun yang lebih sesuai dengan masalah kita.

Tidak perlu memulai dengan deret Fourier lengkap untuk memperoleh
deret sinus dan kosinus.
Deret sinus sebenarnya merupakan ekspansi fungsi eigen dari $f(t)$ dengan menggunakan 
fungsi eigen dari masalah nilai eigen $x''+\lambda x = 0$, $x(0) = 0$,
$x(L) = 0$.  Deret kosinus merupakan ekspansi fungsi eigen dari $f(t)$
dengan menggunakan 
fungsi eigen dari masalah nilai eigen $x''+\lambda x = 0$, $x'(0) = 0$,
$x'(L) = 0$.  Oleh karena itu, kita akan memperoleh rumus yang sama
dengan mendefinisikan hasil kali dalam
\begin{equation*}
\langle f(t), g(t) \rangle = \int_0^L f(t) g(t) \, dt ,
\end{equation*}
dan mengikuti prosedur dalam \sectionref{ts:section}.  Sudut pandang ini
berguna, karena kita lazimnya menggunakan deret tertentu yang muncul karena pertanyaan
yang mendasari kita 
mengarah pada suatu masalah nilai eigen.  Jika masalah nilai eigen 
bukan salah satu dari tiga masalah yang sejauh ini kita bahas, Anda tetap dapat melakukan
ekspansi fungsi eigen, yang menggeneralisasi hasil bab ini.  Kita akan
membahas generalisasi tersebut dalam \chapterref{SL:chapter}.

%f(t) = \frac{a_0}{2} + \sum_{n=1}^\infty a_n \cos \left( \frac{n \pi}{L} 
%t \right)
%+ b_n \sin \left( \frac{n \pi}{L} t \right) ,

\begin{example}
Temukan deret Fourier dari perluasan periodik genap dari 
fungsi $f(t) = t^2$ untuk $0 \leq t \leq \pi$.

Kita ingin menuliskan
\begin{equation*}
f(t) = \frac{a_0}{2} + \sum_{n=1}^\infty a_n \cos (n t) ,
\end{equation*}
dengan
\begin{equation*}
a_0 = \frac{2}{\pi}
\int_0^\pi t^2 \, dt = \frac{2 \pi^2}{3} ,
\end{equation*}
dan
\begin{equation*}
\begin{split}
a_n & = \frac{2}{\pi}
\int_0^\pi t^2 \cos (n t) \, dt
= \frac{2}{\pi} \left[ t^2 \frac{1}{n} \sin (nt) \right]_0^\pi -
\frac{4}{n\pi}
\int_0^\pi t \sin (n t) \, dt \\
& = 
\frac{4}{n^2\pi}
\Bigl[ t \cos (n t) \Bigr]_0^\pi
-
\frac{4}{n^2\pi}
\int_0^\pi \cos (n t) \, dt
= 
\frac{4{(-1)}^n}{n^2} .
\end{split}
\end{equation*}
Perhatikan bahwa kita \myquote{mendeteksi} kekontinuan perluasan tersebut karena
koefisiennya meluruh seperti $\frac{1}{n^2}$.  Artinya, perluasan periodik genap
dari $t^2$ tidak memiliki diskontinuitas lompatan.  Perluasan tersebut memang memiliki sudut, karena
turunannya, yang merupakan fungsi ganjil dan deret sinus, memiliki lompatan; turunannya memiliki
deret Fourier dengan koefisien yang hanya meluruh seperti $\frac{1}{n}$.

Secara eksplisit, beberapa suku pertama deret untuk $f(t)$ adalah
\begin{equation*}
\frac{\pi^2}{3} - 4 \cos (t) + \cos (2t) - \frac{4}{9} \cos (3t) + \cdots
\end{equation*}
\end{example}

\begin{exercise}
\leavevmode
\begin{tasks}
\task Hitung turunan dari perluasan periodik genap $f(t)$ di atas dan verifikasi bahwa turunan tersebut
memiliki diskontinuitas lompatan.  Gunakan definisi sebenarnya dari $f(t)$, bukan deret
kosinusnya!
\task Mengapa turunan dari perluasan periodik genap $f(t)$ merupakan
perluasan periodik ganjil dari $f'(t)$?
\end{tasks}
\end{exercise}

\subsection{Penerapan}

Deret Fourier berkaitan dengan masalah nilai batas yang
kita pelajari sebelumnya.
Perhatikan masalah nilai batas
\begin{equation*}
x''(t) + \lambda\, x(t) = f(t) , \quad 0 < t < L,
\end{equation*}
dengan \emph{\myindex{syarat batas Dirichlet}}
$x(0) = 0$, $x(L) = 0$.
Alternatif Fredholm (\thmvref{thm:fredholmsimple})
menyatakan bahwa
selama $\lambda$ bukan nilai eigen dari masalah homogen yang mendasarinya,
terdapat solusi tunggal.
Fungsi eigen dari masalah nilai eigen ini adalah fungsi-fungsi
$\sin \left( \frac{n \pi}{L} t \right)$.
Untuk menemukan solusi,
pertama kita mencari deret sinus Fourier untuk $f(t)$.
Kita juga menuliskan $x(t)$ sebagai deret sinus, tetapi dengan koefisien yang belum diketahui.  
Kita menyubstitusikan deret untuk $x(t)$
ke dalam persamaan dan menyelesaikan koefisien yang belum diketahui.
Jika kita memiliki
\emph{\myindex{syarat batas Neumann}}
$x'(0) = 0$, $x'(L) = 0$, kita melakukan prosedur yang sama dengan menggunakan deret
kosinus.
%
Mari kita lihat cara kerja metode ini pada contoh-contoh.

\begin{example}
Ambil masalah nilai batas
\begin{equation*}
x''(t) + 2 x(t) = f(t) ,
\quad
0 < t < 1,
\end{equation*}
dengan $f(t) = t$ pada $0 < t < 1$, dan 
memenuhi syarat batas Dirichlet
$x(0) = 0$, $x(1)=0$.
Kita menulis $f(t)$ sebagai deret sinus
\begin{equation*}
f(t) = \sum_{n=1}^\infty c_n \sin (n \pi t) .
\end{equation*}
Hitung
\begin{equation*}
c_n = 2 \int_0^1 t \sin (n \pi t) \,dt = \frac{2 \, {(-1)}^{n+1}}{n \pi} .
\end{equation*}
Kita menulis $x(t)$ sebagai
\begin{equation*}
x(t) = \sum_{n=1}^\infty b_n \sin (n \pi t) .
\end{equation*}
Kita menyubstitusikannya untuk memperoleh 
\begin{equation*}
\begin{split}
x''(t) + 2 x(t) & =
\underbrace{
\sum_{n=1}^\infty - b_n n^2 \pi^2 \sin (n \pi t)
}_{x''}
\,
+
\,
2
\underbrace{
\sum_{n=1}^\infty b_n \sin (n \pi t)
}_{x}
\\
& =
\sum_{n=1}^\infty b_n \bigl(2 - n^2 \pi^2 \bigr) \sin (n \pi t)
\\
& = f(t)
=
\sum_{n=1}^\infty  \frac{2\, {(-1)}^{n+1}}{n \pi} \sin (n \pi t) .
\end{split}
\end{equation*}
Oleh karena itu,
\begin{equation*}
b_n \bigl(2 - n^2 \pi^2\bigr)
=
\frac{2\,{(-1)}^{n+1}}{n \pi}
\end{equation*}
atau
\begin{equation*}
b_n
=
\frac{2\,{(-1)}^{n+1}}{n \pi (2 - n^2 \pi^2)} .
\end{equation*}
Fakta bahwa $2-n^2\pi^2$ tidak nol untuk setiap $n$, dan bahwa kita dapat
menyelesaikan $b_n$, terjadi tepat karena
$2$ bukan nilai eigen masalah tersebut.
Dengan demikian kita telah memperoleh deret Fourier untuk solusi
\begin{equation*}
x(t) = 
\sum_{n=1}^\infty
\frac{2\,{(-1)}^{n+1}}{n \pi \,(2 - n^2 \pi^2)}
\sin (n \pi t) .
\end{equation*}
Lihat \figurevref{bnd-dirich-graph:fig} untuk grafik solusi tersebut.
Perhatikan bahwa karena fungsi-fungsi eigen memenuhi syarat batas, 
dan $x$ dituliskan dalam fungsi-fungsi eigen, maka $x$
memenuhi syarat batas.
\begin{myfig}
\capstart
\diffyincludegraphics{width=3in}{width=4.5in}{bnd-dirich-graph}{%
Grafik suatu fungsi pada bidang tx untuk t antara 0 dan 1.
Grafik bermula di 0 ketika t sama dengan 0 dan turun cukup lurus untuk
sesaat.  Kemudian grafik melengkung ke atas, mula-mula mencapai minimum sekitar
minus 0,08 sebelum naik kembali ke 0 ketika t sama dengan 1.  Nilai minimum
sedikit berada di sebelah kanan pusat gambar.}
\caption{Grafik solusi dari $x''+2x=t$, $x(0)=0$, $x(1)=0$.%
\label{bnd-dirich-graph:fig}}
\end{myfig}
\end{example}

\begin{example}
Kita menangani syarat Neumann dengan deret kosinus.
Ambil masalah nilai batas
\begin{equation*}
x''(t) + 2 x(t) = f(t) , \quad 0 < t < 1,
\end{equation*}
dengan kembali $f(t) = t$ pada $0 < t < 1$, tetapi sekarang memenuhi
syarat batas Neumann
$x'(0) = 0$, $x'(1)=0$.
Kita menulis $f(t)$ sebagai deret kosinus
\begin{equation*}
f(t) = \frac{c_0}{2} + \sum_{n=1}^\infty c_n \cos (n \pi t) ,
\end{equation*}
dengan
\begin{equation*}
c_0 = 2 \int_0^1 t \,dt = 1 ,
\end{equation*}
dan
\begin{equation*}
c_n = 2 \int_0^1 t \cos (n \pi t) \,dt =
\frac{2\bigl({(-1)}^n-1\bigr)}{\pi^2 n^2} = 
\begin{cases}
\frac{-4}{\pi^2 n^2} & \text{jika } n \text{ ganjil} , \\
0 & \text{jika } n \text{ genap}.
\end{cases}
\end{equation*}
Kita menulis $x(t)$ sebagai deret kosinus
\begin{equation*}
x(t) = \frac{a_0}{2} + \sum_{n=1}^\infty a_n \cos (n \pi t) .
\end{equation*}
Kita menyubstitusikannya untuk memperoleh 
\begin{equation*}
\begin{split}
x''(t) + 2 x(t) & =
\left(
\sum_{n=1}^\infty - a_n n^2 \pi^2 \cos (n \pi t)
\right)
+
a_0 +
2
\left(
\sum_{n=1}^\infty a_n \cos (n \pi t)
\right)
\\
& =
a_0 +
\sum_{n=1}^\infty a_n \bigl(2 - n^2 \pi^2 \bigr) \cos (n \pi t)
\\
& = f(t)
=
\frac{1}{2} +
\sum_{\substack{n=1\\n~\text{ganjil}}}^\infty
\frac{-4}{\pi^2 n^2} \cos (n \pi t) .
\end{split}
\end{equation*}
Oleh karena itu,
$a_0 = \frac{1}{2}$,
$a_n = 0$ untuk $n$ genap ($n \geq 2$), dan untuk
$n$ ganjil, kita memiliki
\begin{equation*}
a_n \bigl(2 - n^2 \pi^2\bigr)
=
\frac{-4}{\pi^2 n^2} ,
\end{equation*}
atau
\begin{equation*}
a_n
=
\frac{-4}{n^2 \pi^2 (2 - n^2 \pi^2)} .
\end{equation*}
Deret Fourier untuk solusi $x(t)$ adalah
\begin{equation*}
x(t) = 
\frac{1}{4} +
\sum_{\substack{n=1\\n~\text{ganjil}}}^\infty
\frac{-4}{n^2 \pi^2 (2 - n^2 \pi^2)} 
\cos (n \pi t) .
\end{equation*}
\end{example}

\subsection{Latihan}

\begin{exercise}
Ambil $f(t) = {(t-1)}^2$ yang didefinisikan pada $0 \leq t \leq 1$.
\begin{tasks}
\task Sketsakan grafik perluasan periodik genap dari $f$.
\task Sketsakan grafik perluasan periodik ganjil dari $f$.
\end{tasks}
\end{exercise}

\begin{exercise}
Temukan deret Fourier dari perluasan periodik ganjil maupun genap dari 
fungsi $f(t) = {(t-1)}^2$ untuk $0 \leq t \leq 1$.
Dapatkah Anda menentukan perluasan mana yang kontinu dari koefisien
deret Fouriernya?
\end{exercise}

\begin{exercise}
Temukan deret Fourier dari perluasan periodik ganjil maupun genap dari 
fungsi $f(t) = t$ untuk $0 \leq t \leq \pi$.
\end{exercise}

\begin{exercise}
Temukan deret Fourier dari perluasan periodik genap dari 
fungsi $f(t) = \sin t$ untuk $0 \leq t \leq \pi$.
\end{exercise}

\begin{exercise}
\pagebreak[2]
Perhatikan
\begin{equation*}
x''(t) + 4 x(t) = f(t) ,
\end{equation*}
dengan $f(t) = 1$ pada $0 < t < 1$.
\begin{tasks}
\task Selesaikan untuk syarat Dirichlet $x(0)=0, x(1) = 0$.
\task Selesaikan untuk syarat Neumann $x'(0)=0, x'(1) = 0$.
\end{tasks}
\end{exercise}

\begin{exercise}
Perhatikan
\begin{equation*}
x''(t) + 9 x(t) = f(t) ,
\end{equation*}
untuk $f(t) = \sin (2\pi t)$ pada $0 < t < 1$.
\begin{tasks}
\task Selesaikan untuk syarat Dirichlet $x(0)=0, x(1) = 0$.
\task Selesaikan untuk syarat Neumann $x'(0)=0, x'(1) = 0$.
\end{tasks}
\end{exercise}

\begin{exercise}
Perhatikan
\begin{equation*}
x''(t) + 3 x(t) = f(t) , \quad x(0) = 0, \quad x(1) = 0,
\end{equation*}
dengan $f(t) = \sum_{n=1}^\infty b_n \sin (n \pi t)$.  Tuliskan solusi $x(t)$
sebagai deret Fourier, dengan koefisien yang diberikan dalam $b_n$.
\end{exercise}

\begin{exercise}
Misalkan $f(t) = t^2(2-t)$ untuk $0 \leq t \leq 2$.  Misalkan $F(t)$ merupakan perluasan periodik
ganjil.  Hitung $F(1)$, $F(2)$, $F(3)$, $F(-1)$, $F(\nicefrac{9}{2})$,
$F(101)$, $F(103)$.  Catatan: \textbf{Jangan} hitung deret sinusnya.
\end{exercise}

\setcounter{exercise}{100}

\begin{exercise}
Misalkan $f(t) = \nicefrac{t}{3}$ pada $0 \leq t < 3$.
\begin{tasks}
\task Temukan deret Fourier dari perluasan periodik genap.
\task Temukan deret Fourier dari perluasan periodik ganjil.
\end{tasks}
\end{exercise}
\exsol{%
a)
$\nicefrac{1}{2}
+
\sum\limits_{\substack{n=1\\n\text{ ganjil}}}^\infty
\frac{-4}{\pi^2 n^2}
\cos\bigl(\frac{n\pi}{3} t \bigr)$
\qquad
b) 
$\sum\limits_{n=1}^\infty
\frac{2{(-1)}^{n+1}}{\pi n}
\sin\bigl(\frac{n\pi}{3} t \bigr)$
}

\begin{exercise}
Misalkan $f(t) = \cos(2t)$ pada $0 \leq t < \pi$.
\begin{tasks}
\task Temukan deret Fourier dari perluasan periodik genap.
\task Temukan deret Fourier dari perluasan periodik ganjil.
\end{tasks}
\end{exercise}
\exsol{%
a)
$\cos(2t)$
\qquad
b) 
$\sum\limits_{\substack{n=1 \\n \text{ ganjil}}}^\infty
\frac{4n}{\pi n^2 - 4 \pi}
\sin(n t)$
}

\begin{exercise}
Misalkan $f(t)$ didefinisikan pada $0 \leq t < 1$.  Perhatikan
rata-rata kedua perluasan
$g(t) = \frac{F_{\text{odd}}(t)+ F_{\text{even}}(t)}{2}$.
\begin{tasks}(2)
\task Berapakah $g(t)$ jika $0 \leq t < 1$ (Berikan alasan!)
\task Berapakah $g(t)$ jika $-1 < t < 0$ (Berikan alasan!)
\end{tasks}
\end{exercise}
\exsol{%
a) $f(t)$
\qquad
b) $0$
}

\begin{exercise}
Misalkan $f(t) = \sum_{n=1}^\infty \frac{1}{n^2} \sin(nt)$.  Selesaikan
$x''- x = f(t)$ untuk syarat Dirichlet $x(0) = 0$
dan $x(\pi) = 0$.
\end{exercise}
\exsol{%
$\sum\limits_{n=1}^\infty \frac{-1}{n^2(1+n^2)} \sin(nt)$
}

\begin{exercise}[menantang]
Misalkan $f(t) = t + \sum_{n=1}^\infty \frac{1}{2^n} \sin(nt)$.  Selesaikan
$x'' + \pi x = f(t)$ untuk syarat Dirichlet $x(0) = 0$
dan $x(\pi) = 1$.  Petunjuk:  Perhatikan bahwa $\frac{t}{\pi}$ memenuhi
syarat Dirichlet yang diberikan.
\end{exercise}
\exsol{%
$\frac{t}{\pi} + \sum\limits_{n=1}^\infty \frac{1}{2^n(\pi-n^2)} \sin(nt)$
}


%%%%%%%%%%%%%%%%%%%%%%%%%%%%%%%%%%%%%%%%%%%%%%%%%%%%%%%%%%%%%%%%%%%%%%%%%%%%%%

\sectionnewpage
\section{Penerapan deret Fourier}
\label{appoffourier:section}

%mbxINTROSUBSECTION

\sectionnotes{2 kuliah\EPref{, \S9.4 dalam \cite{EP}}\BDref{,
tidak ada dalam \cite{BD}}}

\subsection{Osilasi yang dipaksa secara periodik}

\begin{mywrapfigsimp}{2.0in}{2.3in}
\noindent
\inputpdft{massfigforce}{%
%perhatikan bahwa diagram ini muncul lebih dari sekali dalam buku
Diagram yang menunjukkan massa berlabel m,
yang terhubung di sebelah kiri ke pegas berlabel k, yang terhubung ke dinding di
sebelah kiri.  Panah yang menunjukkan gaya ke kanan pada massa diberi label F dari t.
Panah di bawah massa yang menunjukkan perpindahan mengarah ke kanan dan
diberi label x.  `Lantai' diagram diberi label `redaman c'.}
\end{mywrapfigsimp}
Kita kembali ke osilasi paksa.  Perhatikan sistem massa-pegas seperti
sebelumnya, dengan massa $m$
pada pegas dengan konstanta pegas $k$,
redaman $c$, dan gaya $F(t)$ yang diterapkan pada massa.  Andaikan 
fungsi pemaksa $F(t)$ periodik-$2L$ untuk suatu $L > 0$.
Kita melihat
masalah ini dalam \chapterref{ho:chapter} dengan $F(t) = F_0 \cos (\omega t)$.  Persamaan
yang mengatur susunan khusus ini adalah
\begin{equation} \label{afs:eq}
mx''(t) + cx'(t) + kx(t) = F(t) .
\end{equation}

Solusi umum dari \eqref{afs:eq} terdiri atas solusi komplementer $x_c$, yang
menyelesaikan persamaan homogen terkait $mx'' + cx' + kx = 0$, dan
solusi khusus dari \eqref{afs:eq} yang kita sebut $x_p$.  Untuk $c > 0$,
solusi komplementer $x_c$ akan meluruh seiring berjalannya waktu.
Oleh karena itu,
kita terutama berkepentingan
dengan solusi khusus $x_p$ yang tidak meluruh
dan periodik dengan periode yang sama seperti $F(t)$.  Kita menyebut solusi khusus ini
\emph{\myindex{solusi periodik tunak}} dan menuliskannya sebagai $x_{sp}$ seperti sebelumnya.
Hal baru dalam bagian ini adalah bahwa kita meninjau fungsi pemaksa
sembarang $F(t)$ alih-alih kosinus sederhana.

Demi kesederhanaan, andaikan $c=0$.  Masalah dengan $c > 0$ sangat
serupa.
Persamaan
\begin{equation*}
mx'' + kx = 0 
\end{equation*}
memiliki solusi umum
\begin{equation*}
x(t) = A \cos (\omega_0 t) + 
B \sin (\omega_0 t) ,
\end{equation*}
dengan $\omega_0 = \sqrt{\frac{k}{m}}$.
Setiap solusi untuk
$mx''(t) + kx(t) = F(t)$ berbentuk
$A \cos (\omega_0 t) + B \sin (\omega_0 t) + x_{sp}$.
Solusi periodik
tunak $x_{sp}$ memiliki periode yang sama seperti $F(t)$.

Dengan semangat bagian terakhir dan gagasan koefisien tak tentu,
pertama kita menuliskan
\begin{equation*}
F(t) = \frac{c_0}{2} + \sum_{n=1}^\infty
\Bigl[
c_n \cos \Bigl( \frac{n \pi}{L} t \Bigr) +
d_n \sin \Bigl( \frac{n \pi}{L} t \Bigr)
\Bigr] .
\end{equation*}
Kemudian kita menuliskan usulan solusi periodik tunak $x$ sebagai
\begin{equation*}
x(t) = \frac{a_0}{2} + \sum_{n=1}^\infty
\Bigl[
a_n \cos \Bigl( \frac{n \pi}{L} t \Bigr) +
b_n \sin \Bigl( \frac{n \pi}{L} t \Bigr)
\Bigr],
\end{equation*}
dengan $a_n$ dan $b_n$ belum diketahui.
Kita menyubstitusikan $x$ ke dalam persamaan diferensial dan menyelesaikan $a_n$ dan
$b_n$ dalam $c_n$ dan $d_n$.  Proses ini
mungkin paling mudah dipahami melalui contoh.
\pagebreak[2]

\begin{example} \label{afs:steadyex}
Andaikan bahwa $k=2$ dan $m=1$.
Satuannya kembali merupakan satuan mks\index{satuan mks}
(meter-kilogram-detik).
Sebuah ransel jet diikatkan pada massa, yang menyala dengan gaya 1
newton selama 1
detik lalu mati selama 1 detik, dan seterusnya.  Kita ingin menemukan solusi periodik
tunak.

Oleh karena itu, persamaannya adalah
\begin{equation*}
x'' + 2 x = F(t) ,
\end{equation*}
dengan $F(t)$ fungsi tangga
\begin{equation*}
F(t) =
\begin{cases}
0 & \text{jika } \; {-1} < t < 0 , \\
1 & \text{jika } \; \phantom{-}0 < t < 1 ,
\end{cases}
\end{equation*}
yang diperluas secara periodik.
Kita menulis
\begin{equation*}
F(t) = \frac{c_0}{2} + \sum_{n=1}^\infty
\Bigl[
c_n \cos (n \pi t) +
d_n \sin (n \pi t)
\Bigr] .
\end{equation*}
Kita menghitung
\begin{align*}
c_n & = \int_{-1}^1 F(t) \cos (n \pi t) \, dt = 
\int_{0}^1 \cos (n \pi t) \, dt = 0 \qquad \text{untuk } \; n \geq 1,
\\
c_0 & = \int_{-1}^1 F(t) \, dt = 
\int_{0}^1 \, dt = 1 ,
\\
d_n & = \int_{-1}^1 F(t) \sin (n \pi t) \, dt
\\
& = \int_{0}^1 \sin (n \pi t) \, dt
\\
& = \left[ \frac{-\cos (n \pi t)}{n \pi} \right]_{t=0}^1
\\
& = \frac{1-{(-1)}^n}{\pi n} =
\begin{cases}
\frac{2}{\pi n} & \text{jika } n \text{ ganjil} , \\
0 & \text{jika } n \text{ genap} .
\end{cases}
\end{align*}
Jadi
\begin{equation*}
F(t) = \frac{1}{2} + \sum_{\substack{n=1 \\ n \text{ ganjil}}}^\infty
\frac{2}{\pi n} \sin (n \pi t) .
\end{equation*}

Kita ingin mencoba
\begin{equation*}
x(t) = \frac{a_0}{2} + \sum_{n=1}^\infty
\Bigl[
a_n \cos (n \pi t) +
b_n \sin (n \pi t)
\Bigr] .
\end{equation*}
Setelah kita menyubstitusikan $x$ ke dalam persamaan diferensial $x''+2x = F(t)$,
jelas bahwa $a_n = 0$ untuk $n \geq 1$ karena tidak ada suku yang bersesuaian
dalam deret untuk
$F(t)$.  Demikian pula, $b_n = 0$ untuk $n$ genap.  Maka kita mencoba
\begin{equation*}
x(t) = \frac{a_0}{2} +
\sum_{\substack{n=1 \\ n \text{ ganjil}}}^\infty
b_n \sin (n \pi t) .
\end{equation*}
Kita menyubstitusikannya ke dalam persamaan diferensial dan memperoleh
\begin{equation*}
\begin{split}
x'' + 2 x & =
\left(
\sum_{\substack{n=1 \\ n \text{ ganjil}}}^\infty
- b_n n^2 \pi^2 \sin (n \pi t)
\right)
+ 
a_0 +
2
\left(
\sum_{\substack{n=1 \\ n \text{ ganjil}}}^\infty
b_n \sin (n \pi t)
\right)
\\
& =
a_0 +
\sum_{\substack{n=1 \\ n \text{ ganjil}}}^\infty
b_n (2 - n^2 \pi^2 ) \sin (n \pi t)
\\
& =
F(t) = \frac{1}{2} + \sum_{\substack{n=1 \\ n \text{ ganjil}}}^\infty
\frac{2}{\pi n} \sin (n \pi t) .
\end{split}
\end{equation*}
Jadi $a_0 = \frac{1}{2}$, $b_n = 0$ untuk $n$ genap, dan untuk $n$ ganjil, kita
memperoleh
\begin{equation*}
b_n = 
\frac{2}{\pi n (2 - n^2 \pi^2 )} .
\end{equation*}

Solusi periodik tunak memiliki deret Fourier
\begin{equation*}
x_{sp}(t) = \frac{1}{4} + \sum_{\substack{n=1 \\ n \text{ ganjil}}}^\infty
\frac{2}{\pi n (2 - n^2 \pi^2 )}
\sin (n \pi t) .
\end{equation*}
Kita tahu bahwa ini merupakan solusi periodik tunak karena tidak memuat suku 
dari solusi komplementer dan periodik dengan periode yang sama seperti
$F(t)$ itu sendiri.  Lihat \figurevref{afs:steadyexfig} untuk grafik solusi ini.
\begin{myfig}
\capstart
\diffyincludegraphics{width=3in}{width=4.5in}{afs-steadyex}{%
Grafik fungsi periodik yang tampak seperti lintasan sinusoidal
yang berosilasi beberapa kali saat t bergerak dari 0 ke 10.  Grafik bermula di 0,25
ketika t sama dengan 0, turun sedikit ke bawah 0,2, lalu naik kembali sedikit ke atas 0,3,
dan seterusnya.  Grafik melakukan 5 osilasi pada interval tersebut.}
\caption{Grafik solusi periodik tunak $x_{sp}$ dari
\exampleref{afs:steadyex}.%
\label{afs:steadyexfig}}
\end{myfig}
\end{example}

\subsection{Resonansi}

Seperti ketika fungsi pemaksa berupa kosinus sederhana, kita dapat menjumpai
resonansi.  Kita mengandaikan $c=0$ sehingga hanya membahas resonansi murni.
Misalkan $F(t)$ periodik-$2L$ dan perhatikan
\begin{equation*}
m x''(t) + k x (t) = F(t) .
\end{equation*}
Ketika kita mengekspansi $F(t)$ dan mendapati bahwa beberapa sukunya berimpit dengan
solusi komplementer untuk $mx''+kx=0$, kita tidak dapat menggunakan suku-suku tersebut dalam
tebakan.  Sama seperti sebelumnya, suku-suku itu menghilang ketika kita menyubstitusikannya ke ruas kiri,
dan kita memperoleh persamaan yang bertentangan (seperti $0=1$).   Artinya,
andaikan
\begin{equation*}
x_c = A \cos (\omega_0 t) + B \sin (\omega_0 t), 
\end{equation*}
dengan $\omega_0 = \frac{N \pi}{L}$ untuk suatu bilangan bulat positif $N$.
Kita harus
memodifikasi tebakan dan mencoba
\begin{equation*}
x(t) = \frac{a_0}{2} +
t \Biggl(
a_N \cos \Bigl( \frac{N \pi}{L} t \Bigr) +
b_N \sin \Bigl( \frac{N \pi}{L} t \Bigr)
\Biggr) +
\sum_{\substack{n=1\\n \neq N}}^\infty
\Bigl[
a_n \cos \Bigl( \frac{n \pi}{L} t \Bigr) +
b_n \sin \Bigl( \frac{n \pi}{L} t \Bigr)
\Bigr].
\end{equation*}
Dengan kata lain, kita mengalikan suku yang bermasalah dengan $t$.  Sejak itu, kita
melanjutkan seperti sebelumnya.

Solusinya bukan deret Fourier (bahkan tidak
periodik) karena memuat suku-suku tersebut yang dikalikan dengan $t$.
Suku-suku
$t \left( a_N \cos \left( \frac{N \pi}{L} t \right) +
b_N \sin \left( \frac{N \pi}{L} t \right) \right)$ pada akhirnya mendominasi dan menyebabkan
osilasi liar.  Seperti sebelumnya, perilaku ini disebut \emph{\myindex{resonansi
murni}} atau cukup \emph{\myindex{resonansi}}.

Perhatikan bahwa kita mungkin mengenai frekuensi resonansi dengan tak berhingga
banyak suku (nada atas) dari fungsi pemaksa $F$.
Artinya, andaikan kita menggunakan \myquote{bentuk} yang sama untuk $F$, dan mengubah
frekuensi dasar (kita mengubah $L$).
Maka suku-suku yang berbeda
dari deret Fourier $F$ dapat berinterferensi dengan solusi komplementer,
yaitu, $\frac{n \pi}{L} = \omega_0$ untuk suatu $n$, dan mungkin
menyebabkan resonansi.
Secara teoretis, tak berhingga banyak frekuensi dasar
dapat menyebabkan resonansi.
Namun, perlu kita perhatikan bahwa karena segala sesuatu merupakan pendekatan terhadap
setiap penerapan dalam kehidupan nyata,
hanya beberapa
suku pertama $F$, dan karena itu hanya beberapa frekuensi semacam itu,
yang akan berpengaruh.

\begin{example}
Kita ingin menyelesaikan persamaan
\begin{equation} \label{afs:eq-resonance}
2 x'' + 18 \pi^2 x = F(t) ,
\end{equation}
dengan
\begin{equation*}
F(t) =
\begin{cases}
-1 & \text{jika } \; {-1} < t < 0 , \\
1 & \text{jika } \; \phantom{-}0 < t < 1 ,
\end{cases}
\end{equation*}
yang diperluas secara periodik.  Kita perhatikan bahwa
\begin{equation*}
F(t) =
\sum_{\substack{n=1 \\ n \text{ ganjil}}}^\infty
\frac{4}{\pi n}
\sin (n \pi t) . 
\end{equation*}

\begin{exercise}
Hitung deret Fourier dari $F$ untuk memverifikasi persamaan di atas.
\end{exercise}

Karena $\sqrt{\frac{k}{m}} = \sqrt{\frac{18\pi^2}{2}} = 3\pi$,
solusi untuk \eqref{afs:eq-resonance} adalah
\begin{equation*}
x(t) = c_1 \cos  (3\pi t) + c_2 \sin (3\pi t) + x_p (t)
\end{equation*}
untuk suatu solusi khusus $x_p$.

Jika kita sekadar mencoba $x_p$ yang diberikan sebagai deret Fourier dengan $\sin (n\pi t)$ seperti biasa,
persamaan komplementer, $2x''+18\pi^2x=0$, meniadakan harmonik ketiga kita.  Artinya, suku
dengan $\sin(3 \pi t)$
sudah
berada dalam solusi komplementer kita.
Oleh karena itu, kita keluarkan suku tersebut dan
mengalikannya dengan $t$.  Kita juga menambahkan suku kosinus agar semuanya tepat.
Artinya, kita mencoba
\begin{equation*}
x_p(t) =
a_3
t \cos (3 \pi t )
+
b_3
t \sin (3 \pi t)
+
\sum_{\substack{n=1 \\ n~\text{ganjil} \\ n \neq 3}}^\infty
b_n
\sin (n \pi t) . 
\end{equation*}
Kita menghitung turunan keduanya.
\begin{multline*}
x_p''(t) =
- 6 a_3
\pi \, \sin (3 \pi t) - 9\pi^2 a_3 \, t \, \cos (3 \pi t)
+
6 b_3
\pi \, \cos (3 \pi t) - 9\pi^2 b_3 \, t \, \sin (3 \pi t)
\\
{} +
\sum_{\substack{n=1 \\ n~\text{ganjil} \\ n \neq 3}}^\infty
(-n^2 \pi^2 b_n ) \,
\sin (n \pi t) . 
\end{multline*}
Sekarang kita menyubstitusikannya ke ruas kiri persamaan diferensial.
\begin{align*}
2x_p'' + 18\pi^2 x_p 
= & 
- 12 a_3 \pi \sin (3 \pi t)
- 18\pi^2 a_3 t \cos (3 \pi t)
+ 12 b_3 \pi \cos (3 \pi t)
- 18\pi^2 b_3 t \sin (3 \pi t)
\\
& \phantom{\, - 12 a_3 \pi \sin (3 \pi t)} ~
{} + 18 \pi^2 a_3 t \cos (3 \pi t)
\phantom{\, + 12 b_3 \pi \cos (3 \pi t)} ~
{} + 18 \pi^2 b_3 t \sin (3 \pi t)
\\
& {} + \sum_{\substack{n=1 \\ n~\text{ganjil} \\ n \neq 3}}^\infty
(-2n^2 \pi^2 b_n + 18\pi^2 b_n) \,
\sin (n \pi t) . 
\end{align*}
Kita sederhanakan,
\begin{equation*}
2x_p'' + 18\pi^2 x_p =
- 12 a_3
\pi \sin (3 \pi t)
+
12 b_3
\pi \cos (3 \pi t)
+
\sum_{\substack{n=1 \\ n~\text{ganjil} \\ n \neq 3}}^\infty
(-2n^2 \pi^2 b_n + 18\pi^2 b_n)
\sin (n \pi t) . 
\end{equation*}
Deret ini harus sama dengan deret untuk $F(t)$.
Kita menyamakan koefisien dan menyelesaikan $a_3$ serta $b_n$:
\begin{align*}
& a_3 = \frac{4/(3\pi)}{-12\pi} = \frac{-1}{9\pi^2} , \\
& b_3 = 0 , \\
& b_n = \frac{4}{n\pi(18\pi^2 - 2n^2 \pi^2)} 
= \frac{2}{\pi^3 n(9 - n^2)} \qquad \text{untuk } n \text{ ganjil dan } n \neq 3 .
\end{align*}
Artinya,
\begin{equation*}
x_p(t) =
\frac{-1}{9\pi^2}
\,
t \, \cos (3 \pi t)
+
\sum_{\substack{n=1 \\ n~\text{ganjil} \\ n \neq 3}}^\infty
\frac{2}{\pi^3 n(9 - n^2)}
\sin (n \pi t) . 
\end{equation*}
\end{example}

Ketika $c > 0$, Anda tidak perlu mengkhawatirkan resonansi murni.
Tidak pernah ada konflik, dan Anda tidak perlu mengalikan satu pun
suku dengan $t$.  Terdapat konsep yang bersesuaian, yaitu
\myindex{resonansi praktis,}
dan konsep tersebut sangat mirip dengan gagasan yang telah kita telaah dalam
\chapterref{ho:chapter}.
Pada dasarnya, yang terjadi dalam resonansi praktis adalah salah satu
koefisien dalam deret untuk $x_{sp}$ dapat menjadi sangat besar.  Kita tidak akan
membahas rinciannya di sini.

\subsection{Latihan}

\begin{exercise}
Misalkan $F(t) = \frac{1}{2} + \sum_{n=1}^\infty \frac{1}{n^2} \cos (n \pi t)$.
Temukan
solusi periodik tunak untuk
$x'' + 2 x = F(t)$.  Nyatakan solusi Anda sebagai deret Fourier.
\end{exercise}

\begin{exercise}
Misalkan $F(t) = \sum_{n=1}^\infty \frac{1}{n^3} \sin (n \pi t)$.  Temukan
solusi periodik tunak untuk
$x'' + x' + x = F(t)$.  Nyatakan solusi Anda sebagai deret Fourier.
\end{exercise}

\begin{exercise}
Misalkan $F(t) = \sum_{n=1}^\infty \frac{1}{n^2} \cos (n \pi t)$.  Temukan
solusi periodik tunak untuk
$x'' + 4 x = F(t)$.  Nyatakan solusi Anda sebagai deret Fourier.
\end{exercise}

\begin{exercise}
Misalkan $F(t) = t$ untuk $-1 < t < 1$ dan diperluas secara periodik.
Temukan solusi periodik tunak untuk
$x'' + x = F(t)$.  Nyatakan solusi Anda sebagai deret.
\end{exercise}

\begin{exercise}
Misalkan $F(t) = t$ untuk $-1 < t < 1$ dan diperluas secara periodik.
Temukan solusi periodik tunak untuk
$x'' + \pi^2 x = F(t)$.  Nyatakan solusi Anda sebagai deret.
\end{exercise}

\setcounter{exercise}{100}

\begin{exercise}
Misalkan $F(t) = \sin(2\pi t) + 0.1 \cos(10 \pi t)$.
Temukan solusi periodik tunak untuk $x'' + \sqrt{2}\, x = F(t)$.
Nyatakan solusi Anda sebagai deret Fourier.
\end{exercise}
\exsol{%
$x = \frac{1}{\sqrt{2}-4 \pi^2} \sin(2\pi t) + \frac{0.1}{\sqrt{2}-100 \pi^2} \cos(10 \pi t)$
}

\begin{exercise}
Misalkan $F(t) = \sum_{n=1}^\infty e^{-n} \cos(2 n t)$.
Temukan solusi periodik tunak untuk $x'' + 3 x = F(t)$.
Nyatakan solusi Anda sebagai deret Fourier.
\end{exercise}
\exsol{%
$x =
\sum\limits_{n=1}^\infty
\frac{e^{-n}}{3-{(2n)}^2} \cos(2n t)$
}

\begin{exercise}
Misalkan $F(t) = \lvert t \rvert$ untuk $-1 \leq t \leq 1$ yang diperluas secara periodik.
Temukan solusi periodik tunak untuk $x'' + \sqrt{3}\, x = F(t)$.
Nyatakan solusi Anda sebagai deret.
\end{exercise}
\exsol{%
$x =
\frac{1}{2\sqrt{3}} + 
\sum\limits_{\substack{n=1 \\ n \text{ ganjil}}}^\infty \frac{-4}{n^2 \pi^2
(\sqrt{3}-n^2 \pi^2)} \cos (n \pi t)$
}

\begin{exercise}
Misalkan $F(t) = \lvert t \rvert$ untuk $-1 \leq t \leq 1$ yang diperluas secara periodik.
Temukan solusi periodik tunak untuk $x'' + \pi^2 x = F(t)$.
Nyatakan solusi Anda sebagai deret.
\end{exercise}
\exsol{%
$x =
\frac{1}{2\pi^2} -
\frac{2}{\pi^3} t \sin(\pi t) + 
\sum\limits_{\substack{n=3 \\ n \text{ ganjil}}}^\infty \frac{-4}{n^2 \pi^4
(1-n^2)} \cos (n \pi t)$
}


%%%%%%%%%%%%%%%%%%%%%%%%%%%%%%%%%%%%%%%%%%%%%%%%%%%%%%%%%%%%%%%%%%%%%%%%%%%%%%

\sectionnewpage
\section{PDP, pemisahan variabel, dan persamaan panas}
\label{heateq:section}

%mbxINTROSUBSECTION

\sectionnotes{2 kuliah\EPref{, \S9.5 dalam \cite{EP}}\BDref{,
\S10.5 dalam \cite{BD}}}

Ingat bahwa \emph{\myindex{persamaan diferensial parsial}} atau
\emph{\myindex{PDP}} adalah persamaan yang memuat turunan parsial
terhadap \emph{beberapa} variabel bebas.  Menyelesaikan PDP
akan menjadi penerapan utama deret Fourier kita.

Suatu PDP disebut \emph{linear\index{PDP linear}} jika variabel terikat
dan turunannya muncul paling tinggi berpangkat satu dan tidak berada di dalam
fungsi.  Kita hanya akan membahas PDP linear.  Bersama dengan suatu PDP\@,
biasanya kita menetapkan beberapa
\emph{syarat batas\index{syarat batas untuk PDP}},
dengan nilai solusi atau turunannya diberikan sepanjang
batas suatu daerah, dan/atau
beberapa 
\emph{syarat awal\index{syarat awal untuk PDP}} dengan nilai
solusi atau turunannya diberikan untuk suatu waktu awal.
Kadang-kadang syarat-syarat semacam itu dicampurkan dan kita akan menyebutnya
secara umum sebagai 
\emph{syarat samping\index{syarat samping untuk PDP}}.

Kita akan mempelajari tiga
persamaan diferensial parsial khusus, yang masing-masing mewakili suatu
kelas umum persamaan.  Pertama, kita akan mempelajari
\emph{\myindex{persamaan panas}}, yang merupakan contoh
\emph{\myindex{PDP parabolik}}.  Selanjutnya, kita akan mempelajari
\emph{\myindex{persamaan gelombang}}, yang merupakan contoh
\emph{\myindex{PDP hiperbolik}}.  Terakhir, kita akan mempelajari
\emph{\myindex{persamaan Laplace}}, yang merupakan contoh
\emph{\myindex{PDP eliptik}}.  Setiap contoh akan menggambarkan
perilaku yang khas untuk seluruh kelasnya.

\subsection{Panas pada kawat terisolasi}

Kita mulai dengan persamaan panas.
Perhatikan sebuah kawat (atau batang logam tipis) dengan panjang $L$
yang terisolasi sepanjang kawat
kecuali pada titik-titik ujungnya.
Misalkan $x$ menyatakan posisi sepanjang kawat dan $t$ menyatakan waktu.
Lihat \figurevref{heat:wirefig}.

\begin{myfig}
\capstart
\inputpdft{heat-wire}{%
Diagram kawat lurus terisolasi yang terletak pada sumbu x, bermula di
x sama dengan 0 dan berakhir di x sama dengan L.  Kawat terisolasi sepanjang kawat kecuali pada
titik-titik ujungnya.  Sumbu vertikal diberi label `suhu u'.}
\caption{Kawat terisolasi.\label{heat:wirefig}}
\end{myfig}

Misalkan $u(x,t)$ menyatakan suhu pada titik $x$ pada waktu $t$.
Persamaan yang mengatur susunan ini adalah
apa yang disebut \emph{\myindex{persamaan panas satu dimensi}}\index{persamaan panas}:
\begin{equation*}
\mybxbg{~~
\frac{\partial u}{\partial t} =
k \frac{\partial^2 u}{\partial x^2} ,
~~}
\end{equation*}
dengan $k > 0$ suatu konstanta (\emph{\myindex{difusivitas termal}} bahan tersebut).
Artinya, perubahan panas terhadap waktu di suatu titik
sebanding dengan turunan kedua
panas dalam arah $x$---sepanjang kawat.
Ini masuk akal;
jika pada $t$ yang tetap
grafik distribusi panas memiliki suatu maksimum
(grafik cekung ke bawah dan turunan kedua terhadap $x$ negatif),
maka panas seharusnya mengalir menjauhi maksimum sehingga turunan terhadap $t$ juga
harus negatif.
Demikian pula di suatu minimum, panas cenderung mengalir masuk.

Umumnya kita menggunakan notasi yang lebih mudah untuk turunan parsial.
Kita menulis $u_t$ alih-alih $\frac{\partial u}{\partial t}$,
dan menulis $u_{xx}$ alih-alih $\frac{\partial^2 u}{\partial x^2}$.
Dengan notasi ini, persamaan panas menjadi
\begin{equation*}
u_t = k u_{xx} .
\end{equation*}

Daerah tempat kita akan menyelesaikan persamaan panas diberikan oleh 
\begin{equation*}
0 < x < L \quad \text{dan} \quad t > 0 .
\end{equation*}
%10 adalah jumlah baris, harus disesuaikan
\begin{mywrapfigsimp}[10]{2.25in}{2.55in}
\noindent
\inputpdft{heateqregion}{%
Diagram suatu daerah di bidang xt dengan sumbu x sebagai 
sumbu mendatar dan t sebagai sumbu tegak.  Daerah tersebut merupakan setengah lajur tak hingga
dengan t positif dan x berada antara 0 dan L, dan daerah ini diarsir
dalam gambar.  Daerah yang diarsir diberi label dengan
persamaan diferensial u sub t sama dengan k kali u sub xx.
Batas daerah ini diberikan oleh 3 garis.
Batas kiri adalah setengah garis tak hingga sepanjang sumbu t dari t sama dengan 0
hingga tak hingga,
dan garis batas ini diberi label `u sama dengan 0 atau u sub x sama dengan 0'.
Batas bawah adalah ruas sumbu x dari x sama dengan 0 hingga x sama dengan L,
dan ruas ini diberi label u sama dengan f dari x.
Batas kanan adalah setengah garis tegak tak hingga dengan x sama dengan L
dan t positif.
Garis ini diberi label `u sama dengan 0 atau u sub x sama dengan 0'.}
\end{mywrapfigsimp}
Kita juga harus memiliki beberapa syarat samping pada batas-batas daerah tersebut.
Kita mengandaikan bahwa ujung-ujung kawat terbuka 
dan menyentuh suatu benda bersuhu konstan, atau ujung-ujungnya terisolasi.
Jika ujung-ujung kawat dipertahankan pada suhu 0, maka
syaratnya adalah
\begin{equation*}
u(0,t) = 0 \quad \text{dan} \quad u(L,t) = 0
\quad \text{untuk } \; t > 0.
\end{equation*}
Jika, di sisi lain, ujung-ujungnya terisolasi, syaratnya adalah
\begin{equation*}
u_x(0,t) = 0 \quad \text{dan} \quad u_x(L,t) = 0
\quad \text{untuk } \; t > 0 .
\end{equation*}
Mari kita lihat mengapa demikian.
Jika $u_x$ positif di suatu titik $x_0$, maka pada waktu tertentu,
$u$ lebih kecil di sebelah kiri $x_0$ dan lebih tinggi di sebelah kanan $x_0$.
Panas mengalir dari suhu tinggi ke suhu rendah, yaitu ke kiri.
Di sisi lain, jika $u_x$ negatif, panas kembali mengalir
dari suhu tinggi ke suhu rendah,
yaitu ke kanan.  Jadi ketika $u_x$ nol,
kita berada di suatu titik tempat 
panas tidak mengalir ke kedua arah.
Dengan kata lain, $u_x(0,t) = 0$ berarti
tidak ada panas yang mengalir masuk atau keluar dari kawat di titik $x=0$.

Kita memiliki dua syarat sepanjang sumbu $x$ karena terdapat
dua turunan dalam arah $x$.
Syarat samping ini disebut
\emph{homogen\index{syarat samping homogen}}
(yaitu, $u$ atau suatu turunan dari $u$ ditetapkan sama dengan nol).

Kita juga memerlukan syarat awal---distribusi suhu
pada waktu $t=0$.  Yaitu,
\begin{equation*}
u(x,0) = f(x)
\quad \text{untuk } \; 0 < x < L ,
\end{equation*}
untuk suatu fungsi $f(x)$ yang diketahui.
Syarat awal ini bukan syarat samping homogen.

\subsection{Pemisahan variabel}

Persamaan panas linear karena $u$ dan turunannya tidak
muncul dalam pangkat apa pun atau di dalam fungsi apa pun,
dan persamaan ini homogen---tidak ada suku yang bebas dari $u$.
Dengan demikian prinsip \myindex{superposisi} masih berlaku untuk
persamaan panas
(tanpa syarat samping):
Jika $u_1$ dan $u_2$ merupakan
solusi dan $c_1$, $c_2$ merupakan konstanta, maka
$u = c_1 u_1 + c_2 u_2$ juga merupakan solusi.

\begin{exercise}
Verifikasi prinsip superposisi untuk persamaan panas.
\end{exercise}

Superposisi mempertahankan beberapa syarat samping.
Jika $u_1$ dan $u_2$ merupakan
solusi yang memenuhi $u(0,t) = 0$ dan $u(L,t) = 0$,
dan $c_1$, $c_2$ merupakan konstanta, maka
$u = c_1 u_1 + c_2 u_2$ tetap merupakan solusi
yang memenuhi $u(0,t) = 0$ dan $u(L,t) = 0$.  Demikian pula
untuk syarat samping $u_x(0,t) = 0$ dan $u_x(L,t) = 0$.  Secara umum,
superposisi mempertahankan semua syarat samping homogen.

Metode
\emph{pemisahan variabel\index{pemisahan variabel}} adalah
mencoba menemukan solusi yang merupakan hasil kali fungsi satu variabel.
Untuk persamaan panas, kita mencoba menemukan solusi berbentuk
\begin{equation*}
u(x,t) = X(x)T(t) .
\end{equation*}
Terlalu berlebihan untuk berharap bahwa solusi khusus yang diinginkan berbentuk demikian.
Namun, sangat wajar untuk meminta agar kita menemukan
cukup banyak solusi \myquote{blok pembangun} berbentuk
$u(x,t) = X(x)T(t)$ dengan menggunakan prosedur ini,
sehingga solusi yang diinginkan untuk PDP dapat dikonstruksi dari
blok-blok pembangun ini melalui superposisi.

Mari kita coba menyelesaikan masalah persamaan panas
\begin{equation*}
u_t = k u_{xx},
\qquad \text{dengan} \quad
u(0,t) = 0 ,\quad \quad u(L,t) = 0,
\quad \text{dan} \quad u(x,0) = f(x) .
\end{equation*}
Kita menebak $u(x,t) = X(x)T(t)$.  Kita akan mencoba membuat tebakan ini memenuhi
persamaan diferensial, $u_t = k u_{xx}$, dan syarat samping homogen,
$u(0,t) = 0$ dan $u(L,t) = 0$.  Kemudian, karena superposisi mempertahankan
persamaan diferensial dan syarat samping homogen, kita akan mencoba
membangun solusi dari blok-blok pembangun ini untuk menyelesaikan
syarat awal tak homogen $u(x,0) = f(x)$.

Pertama, kita menyubstitusikan $u(x,t) = X(x)T(t)$ ke dalam persamaan panas untuk
memperoleh
\begin{equation*}
X(x)T'(t) = k X''(x)T(t) .
\end{equation*}
Kita menulis ulang sebagai
\begin{equation*}
\frac{T'(t)}{k T(t)} =
\frac{X''(x)}{X(x)} .
\end{equation*}
Persamaan ini harus berlaku untuk semua $x$ dan semua $t$.  Namun
ruas kiri tidak bergantung pada $x$ dan ruas kanan tidak
bergantung pada $t$.  Karena itu, setiap ruas harus merupakan suatu konstanta.  Mari kita sebut
konstanta ini $-\lambda$ (tanda minus untuk kemudahan nanti).
Kita memperoleh dua persamaan
\begin{equation*}
\frac{T'(t)}{k T(t)} = -\lambda =
\frac{X''(x)}{X(x)} .
\end{equation*}
Dengan kata lain,
\begin{align*}
X''(x) + \lambda X(x) &= 0 , \\
T'(t) + \lambda k T(t) &= 0 .
\end{align*}
Syarat batas $u(0,t) = 0$ menyiratkan $X(0)T(t) = 0$.  Kita mencari
solusi tak trivial, sehingga kita dapat mengandaikan bahwa $T(t)$ tidak identik
nol.  Karena itu $X(0) = 0$.  Demikian pula, $u(L,t) = 0$ menyiratkan $X(L) = 0$.  Kita
mencari solusi tak trivial $X$ dari masalah nilai eigen
$X'' + \lambda X = 0$, $X(0) = 0$, $X(L) = 0$.  Sebelumnya kita telah menemukan bahwa
satu-satunya nilai eigen adalah $\lambda_n = \frac{n^2 \pi^2}{L^2}$, untuk bilangan bulat
$n \geq 1$,
dengan fungsi eigen $\sin \left(\frac{n \pi}{L} x\right)$.  Karena itu, mari kita pilih
solusi
\begin{equation*}
X_n (x) = \sin \left(\frac{n \pi}{L} x \right) .
\end{equation*}
$T_n$ yang bersesuaian harus memenuhi persamaan
\begin{equation*}
T_n'(t) + \frac{n^2 \pi^2}{L^2} k T_n(t) = 0 .
\end{equation*}
Ini merupakan salah satu
\hyperref[subsection:fourfundamental]{persamaan dasar} kita,
dan solusinya berupa
fungsi eksponensial:
\begin{equation*}
T_n(t) = e^{\frac{-n^2 \pi^2}{L^2} k t} ,
\end{equation*}
dengan memilih solusi khusus
yang memenuhi $T_n(0) = 1$ untuk kemudahan.
Solusi blok pembangun kita adalah
\begin{equation*}
u_n(x,t) = X_n(x)T_n(t) =
\sin \left( \frac{n \pi}{L} x \right)
e^{\frac{-n^2 \pi^2}{L^2} k t} .
\end{equation*}

Kita perhatikan bahwa $u_n(x,0) = \sin \left( \frac{n \pi}{L} x \right)$.
Kita menulis $f(x)$ sebagai deret sinus
\begin{equation*}
f(x) = \sum_{n=1}^\infty b_n \sin \left(\frac{n \pi}{L}  x \right) .
\end{equation*}
Artinya, kita menemukan deret Fourier dari perluasan periodik ganjil $f(x)$.
Kita menggunakan deret sinus karena deret tersebut bersesuaian dengan masalah nilai eigen untuk
$X(x)$ di atas.
Terakhir, kita menggunakan superposisi untuk menulis solusi sebagai
\begin{equation*}
\mybxbg{~~
u(x,t) = 
\sum_{n=1}^\infty
b_n
u_n(x,t)
=
\sum_{n=1}^\infty
b_n
\sin \left( \frac{n \pi}{L}  x \right)
e^{\frac{-n^2 \pi^2}{L^2} k t} .
~~}
\end{equation*}

Mengapa solusi ini berlaku?  Pertama, perhatikan bahwa solusi ini merupakan solusi untuk
persamaan panas berdasarkan superposisi.  Solusi tersebut memenuhi $u(0,t) = 0$
dan $u(L,t) = 0$, karena $x=0$ atau $x=L$ membuat semua sinus lenyap.
Terakhir, dengan menyubstitusikan $t=0$, kita melihat bahwa $T_n(0) = 1$, sehingga
\begin{equation*}
u(x,0) = 
\sum_{n=1}^\infty
b_n
u_n(x,0)
=
\sum_{n=1}^\infty
b_n
\sin \left( \frac{n \pi}{L} x \right)
=
f(x) .
\end{equation*}

\begin{example}
Perhatikan
kawat terisolasi dengan panjang 1 yang
ujung-ujungnya ditanam dalam es (suhu 0).
Misalkan $k=0.003$ dan distribusi panas awal
adalah $u(x,0) = 50\,x\,(1-x)$.
Lihat \figurevref{heat:wireexinitfig}.
Andaikan kita ingin menemukan fungsi suhu $u(x,t)$.
Secara khusus, andaikan kita
ingin menemukan kapan (pada waktu $t$ berapa)
suhu maksimum dalam kawat 
turun menjadi setengah maksimum awal, yaitu 12,5.

\begin{myfig}
\capstart
\diffyincludegraphics{width=3in}{width=4.5in}{heat-wireex-init}{%
Grafik contoh distribusi suhu pada kawat
pada waktu t sama dengan 0.  Grafik terhadap x pada interval dari 0 hingga 1
berupa parabola terbalik yang bermula di 0 ketika x sama dengan 0 dan kembali berakhir di 0 ketika
x sama dengan 1 serta mencapai maksimum 12,5 ketika x sama dengan 0,5.}
\caption{Distribusi awal suhu dalam
kawat.\label{heat:wireexinitfig}}
\end{myfig}

Kita sedang menyelesaikan masalah PDP berikut:
\begin{equation*}
\begin{array}{ll}
u_t = 0.003 \, u_{xx} & \qquad \text{untuk } \; 0 < x < 1 \text{ dan } t > 0, \\
u(0,t) = u(1,t) = 0 & \qquad \text{untuk } \; t > 0, \\
u(x,0) = 50\,x\,(1-x) & \qquad \text{untuk } \; 0 < x < 1 .
\end{array}
\end{equation*}
Tuliskan $f(x) = 50\,x\,(1-x)$ untuk $0 < x < 1$ sebagai deret sinus
$f(x) = \sum_{n=1}^\infty b_n \sin (n \pi x)$,
dengan
\begin{equation*}
b_n = 2 \int_0^1 50\,x\,(1-x) \sin (n \pi x) \,dx
= 
\frac{200}{{\pi }^{3}{n}^{3}}-\frac{200\,{\left( -1\right) }^{n}}{{\pi }^{3}{n}^{3}}
=
\begin{cases}
0 & \text{jika } n \text{ genap} , \\
\frac{400}{\pi^3 n^3} & \text{jika } n \text{ ganjil} .
\end{cases}
\end{equation*}
Kita menyubstitusikan koefisien ini ke dalam deret untuk $u(x,t)$ guna memperoleh
solusi
\begin{equation*}
u(x,t) = 
\sum_{\substack{n=1 \\ n \text{ ganjil}}}^\infty
\frac{400}{\pi^3 n^3}
\sin (n \pi x )
\, e^{-n^2 \pi^2 \, 0.003 \, t} .
\end{equation*}
Kita menggambar solusi
dalam \figurevref{heat:wireexfig} untuk $0 \leq t \leq 100$.

\begin{myfig}
\capstart
\diffyincludegraphics{width=5in}{width=7.5in}{heat-wireex}{%
Grafik fungsi u dari (x,t) di atas bidang xt
pada daerah dengan x antara 0 dan 1 serta t antara 0
dan 100.  Pada t sama dengan 0, fungsinya merupakan busur ke atas yang mencapai 12,5 di
tengah dan bernilai nol ketika x sama dengan 0 dan x sama dengan 1.
Permukaan kemudian melandai ke bawah dan mendekati bidang xt seiring berjalannya waktu,
sambil tetap bernilai 0 pada tepi dengan x sama dengan 0 atau x sama dengan 1.}
\caption{Grafik suhu $u(x,t)$ pada kawat di posisi $x$
pada waktu $t$ untuk $0 \leq t \leq 100$.  Perhatikan syarat samping
$u(0,t)=u(1,t)=0$ dan bagaimana fungsi eksponensial membuat suhu meluruh
seiring waktu.\label{heat:wireexfig}}
\end{myfig}

Terakhir, kita menjawab pertanyaan tentang suhu maksimum.  Relatif
mudah untuk melihat
bahwa suhu maksimum pada setiap waktu tetap selalu berada di $x=0.5$, yaitu
di tengah kawat.  Grafik $u(x,t)$ menegaskan intuisi ini.
Jika kita menyubstitusikan $x=0.5$, kita memperoleh
\begin{equation*}
u(0.5,t) = 
\sum_{\substack{n=1 \\ n \text{ ganjil}}}^\infty
\frac{400}{\pi^3 n^3}
\sin (n \pi\, 0.5 )
\, e^{-n^2 \pi^2 \, 0.003 \, t} .
\end{equation*}
Untuk $n=3$ dan seterusnya (ingat bahwa $n$ hanya ganjil), suku-suku
deret tersebut
tidak berarti dibandingkan dengan suku pertama.
Suku pertama dalam deret sudah merupakan pendekatan yang sangat baik
untuk fungsi tersebut.  
Karena itu 
\begin{equation*}
u(0.5,t) \approx
\frac{400}{\pi^3}
\, e^{-\pi^2 \, 0.003 \, t} .
\end{equation*}
Pendekatan tersebut semakin baik ketika $t$ semakin besar,
karena suku-suku lain
meluruh jauh lebih cepat.
Kita menggambar fungsi $u(0.5,t)$, yaitu suhu di titik tengah kawat
pada waktu $t$, serta pendekatannya oleh suku pertama
dalam \figurevref{heat:wireexmaxfig}.

\begin{myfig}
\capstart
\diffyincludegraphics{width=3in}{width=4.5in}{heat-wireex-max}{%
Grafik dua fungsi yang sangat mirip untuk t dari 0 hingga 1, keduanya meluruh
ke nilai yang relatif kecil ketika t sama dengan 1.  Kurva bawah bermula
tepat di 12,5 ketika t sama dengan 0, sedangkan kurva atas sedikit sekali berada di atasnya
ketika t sama dengan 0.  Keduanya semakin mendekat dan ketika t sama dengan 1 keduanya sulit
dibedakan.}
\caption{Suhu di titik tengah kawat (kurva bawah),
dan pendekatan suhu ini dengan hanya menggunakan suku pertama dalam
deret (kurva atas).\label{heat:wireexmaxfig}}
\end{myfig}

Setelah sekitar $t=5$,
akan sulit membedakan
antara suku pertama deret untuk $u(x,t)$ dan 
solusi sebenarnya $u(x,t)$.  Perilaku ini
merupakan ciri umum penyelesaian persamaan panas.
Jika Anda berkepentingan dengan perilaku untuk $t$ yang cukup besar, mungkin hanya
satu atau dua suku pertama yang diperlukan.

Kita kembali ke pertanyaan tentang kapan suhu maksimum menjadi setengah dari
suhu maksimum awal.  Artinya, kapan suhu
di titik tengah adalah $\nicefrac{12.5}{2} = 6.25$.  Grafik tersebut menunjukkan
bahwa pendekatan dengan suku pertama akan cukup dekat.  Kita
menyelesaikan
\begin{equation*}
6.25 =
\frac{400}{\pi^3}
\, e^{-\pi^2 \, 0.003 \, t} .
\end{equation*}
Artinya,
\begin{equation*}
t =
\frac{\ln \frac{6.25\,\pi^3}{400}}{-\pi^2 0.003}
\approx 24.5 .
\end{equation*}
Jadi suhu maksimum turun menjadi setengahnya pada sekitar $t=24.5$.
\end{example}

Kita menyebutkan suatu perilaku menarik dari solusi persamaan panas.
Persamaan panas
\myquote{memuluskan} fungsi $f(x)$ ketika $t$ membesar.  Untuk $t$ yang tetap,
solusinya merupakan deret Fourier dengan koefisien
$b_n e^{\frac{-n^2 \pi^2}{L^2} k t}$.  Jika $t > 0$, maka koefisien ini
menuju nol lebih cepat daripada setiap $\frac{1}{n^p}$ untuk sembarang pangkat $p$.  Dengan kata
lain, deret Fourier tersebut memiliki tak berhingga banyak turunan di mana-mana.
Dengan demikian, sekalipun fungsi $f(x)$ memiliki lompatan dan sudut, untuk
$t > 0$ yang tetap, solusi
$u(x,t)$ sebagai fungsi $x$ semulus apa pun yang kita
inginkan.

\begin{example}
Ketika syarat awal sudah berupa deret sinus, tidak ada yang perlu
dihitung; Anda cukup menyubstitusikannya.  Perhatikan
\begin{equation*}
u_t = 0.3 \, u_{xx}, \qquad u(0,t)=u(1,t)=0, \qquad u(x,0) = 0.1 \sin(\pi x) +
\sin(2\pi x) .
\end{equation*}
Solusinya kemudian adalah
\begin{equation*}
u(x,t) =
0.1 \sin(\pi x) e^{- 0.3 \pi^2 t}
+ 
\sin(2 \pi x) e^{- 1.2 \pi^2 t} .
\end{equation*}
\end{example}

\subsection{Ujung-ujung terisolasi}

Sekarang andaikan ujung-ujung kawat terisolasi.
Dalam hal ini, kita menyelesaikan
masalah
\begin{equation*}
u_t = k u_{xx}
\qquad \text{dengan} \quad
u_x(0,t) = 0, \quad u_x(L,t) = 0,
\quad \text{dan} \quad u(x,0) = f(x) .
\end{equation*}
Sekali lagi kita mencoba solusi berbentuk $u(x,t) = X(x)T(t)$.  Dengan prosedur yang
sama seperti sebelumnya, kita menyubstitusikannya ke dalam persamaan panas dan sampai pada
dua persamaan berikut:
\begin{align*}
X''(x) + \lambda X(x) &= 0 , \\
T'(t) + \lambda k T(t) &= 0 .
\end{align*}
Pada titik ini, ceritanya sedikit berubah.
Syarat batas $u_x(0,t) = 0$ menyiratkan $X'(0)T(t) = 0$.
Karena itu $X'(0) = 0$.  Demikian pula, $u_x(L,t) = 0$ menyiratkan $X'(L) = 0$.
Kita menginginkan solusi tak trivial $X$ dari masalah nilai eigen
$X'' + \lambda X = 0$, $X'(0) = 0$, $X'(L) = 0$.  Sebelumnya kita menemukan bahwa
satu-satunya nilai eigen adalah $\lambda_n = \frac{n^2 \pi^2}{L^2}$, untuk bilangan bulat
$n \geq 0$,
dengan fungsi eigen $\cos \left( \frac{n \pi}{L} x\right)$
(kita menyertakan fungsi eigen
konstan).  Kita memilih
solusi
\begin{equation*}
X_n (x) = \cos \left( \frac{n \pi}{L} x \right)
\qquad \text{dan} \qquad
X_0 (x) = 1.
\end{equation*}
$T_n$ yang bersesuaian harus memenuhi persamaan
\begin{equation*}
T_n'(t) + \frac{n^2 \pi^2}{L^2} k T_n(t) = 0 .
\end{equation*}
Untuk $n \geq 1$, seperti sebelumnya,
\begin{equation*}
T_n(t) = e^{\frac{-n^2 \pi^2}{L^2} k t} .
\end{equation*}
Untuk $n = 0$, kita memiliki $T_0'(t) = 0$ sehingga $T_0(t) = 1$.
Solusi blok pembangun kita adalah
\begin{equation*}
u_n(x,t) = X_n(x)T_n(t) =
\cos \left( \frac{n \pi}{L} x \right)
e^{\frac{-n^2 \pi^2}{L^2} k t}
\end{equation*}
dan
\begin{equation*}
u_0(x,t) = 1 .
\end{equation*}

Kita perhatikan bahwa $u_n(x,0) = \cos \left( \frac{n \pi}{L} x \right)$.
Kita menulis $f$ dengan menggunakan deret kosinus
\begin{equation*}
f(x) = \frac{a_0}{2} + \sum_{n=1}^\infty a_n \cos \left( \frac{n \pi}{L} x
\right) .
\end{equation*}
Artinya, kita menemukan deret Fourier dari perluasan periodik genap $f(x)$.

Kita menggunakan superposisi untuk menulis solusi sebagai
\begin{equation*}
\mybxbg{~~
u(x,t) = 
\frac{a_0}{2} + 
\sum_{n=1}^\infty
a_n
u_n(x,t)
=
\frac{a_0}{2} + 
\sum_{n=1}^\infty
a_n
\cos \left( \frac{n \pi}{L} x \right)
e^{\frac{-n^2 \pi^2}{L^2} k t} .
~~}
\end{equation*}

\begin{example}
Cobalah persamaan yang sama seperti sebelumnya, tetapi dengan ujung-ujung terisolasi.
Kita menyelesaikan masalah PDP berikut
\begin{equation*}
\begin{array}{ll}
u_t = 0.003 \, u_{xx} & \qquad \text{untuk } \; 0 < x < 1 \text{ dan } t > 0, \\
u_x(0,t) = u_x(1,t) = 0  & \qquad  \text{untuk } \; t > 0, \\
u(x,0) = 50\,x\,(1-x) & \qquad \text{untuk } \; 0 < x < 1 .
\end{array}
\end{equation*}


Untuk masalah ini, kita harus menemukan deret kosinus dari $u(x,0)$.
Untuk $0 < x < 1$, kita memiliki
\begin{equation*}
50\, x\,(1-x)
=
\frac{25}{3} +
\sum_{\substack{n=2 \\ n \text{ genap}}}^\infty
\left( \frac{-200}{\pi^2 n^2} \right)
\cos (n \pi x) .
\end{equation*}
Perhitungannya diserahkan kepada pembaca.
Karena itu, solusi masalah PDP, yang digambar dalam
\figurevref{heat:wireisolexfig}, diberikan oleh deret
\begin{equation*}
u(x,t)
=
\frac{25}{3} +
\sum_{\substack{n=2 \\ n \text{ genap}}}^\infty
\left( \frac{-200}{\pi^2 n^2} \right)
\cos ( n \pi x)
\, e^{-n^2 \pi^2 \, 0.003 \, t} .
\end{equation*}

\begin{myfig}
\capstart
\diffyincludegraphics{width=5in}{width=7.5in}{heat-wireisolex}{%
Grafik fungsi u dari (x,t) di atas bidang xt
pada daerah dengan x antara 0 dan 1 serta t antara 0
dan 30.  Pada t sama dengan 0, fungsinya merupakan busur ke atas yang mencapai 12,5 di
tengah dan bernilai nol ketika x sama dengan 0 dan x sama dengan 1.
Permukaan kemudian tampak berusaha mendekati bidang mendatar datar pada u sama dengan
sekitar 8,33 ketika t membesar.}
\caption{Grafik suhu kawat terisolasi pada posisi $x$
pada waktu $t$.\label{heat:wireisolexfig}}
\end{myfig}

Perhatikan pada grafik
bahwa seiring berjalannya waktu, suhu merata di sepanjang kawat.  Pada akhirnya, semua
suku kecuali suku konstan
lenyap, dan tersisa suhu seragam
$\frac{25}{3} \approx 8.33$ sepanjang seluruh kawat.
\end{example}

Mari kita uraikan lebih lanjut poin terakhir.  Suku konstan dalam deret adalah
\begin{equation*}
\frac{a_0}{2} = \frac{1}{L} \int_0^L f(x) \, dx .
\end{equation*}
Dengan kata lain, $\frac{a_0}{2}$ adalah nilai rata-rata $f(x)$, yaitu,
rata-rata suhu awal.  Karena kawat terisolasi
di mana-mana, tidak ada panas yang dapat keluar dan tidak ada panas yang dapat masuk.  Jadi suhu
berusaha terdistribusi secara merata seiring waktu, dan suhu rata-rata harus selalu
sama; khususnya, selalu $\frac{a_0}{2}$.  Ketika waktu menuju
tak hingga, suhu menuju konstanta $\frac{a_0}{2}$ di mana-mana.

\subsection{Latihan}

\begin{exercise}
Perhatikan kawat dengan panjang 2, dengan $k=0.001$ dan distribusi
suhu awal $u(x,0) = 50 x$.  Kedua ujungnya
ditanam dalam es (suhu 0).  Temukan solusinya sebagai deret.
\end{exercise}

\begin{exercise}
Temukan solusi deret dari
\begin{equation*}
\begin{array}{ll}
u_t =  u_{xx} & \qquad \text{untuk } \; 0 < x < 1 \text{ dan } t > 0, \\
u(0,t) = u(1,t) = 0 & \qquad \text{untuk } \; t > 0, \\
u(x,0) = 100 & \qquad \text{untuk } \; 0 < x < 1 .
\end{array}
\end{equation*}
\end{exercise}

\begin{exercise}
Temukan solusi deret dari
\begin{equation*}
\begin{array}{ll}
u_t =  u_{xx} & \qquad \text{untuk } \; 0 < x < \pi \text{ dan } t > 0, \\
u_x(0,t) = u_x(\pi,t) = 0 & \qquad \text{untuk } \; t > 0, \\
u(x,0) = 3\cos (x) + \cos (3x) & \qquad \text{untuk } \; 0 < x < \pi .
\end{array}
\end{equation*}
\end{exercise}

\begin{exercise} \label{heat:cosexr}
Temukan solusi deret dari
\begin{equation*}
\begin{array}{ll}
u_t = \frac{1}{3} u_{xx} & \qquad \text{untuk } \; 0 < x < \pi \text{ dan } t > 0, \\
u_x(0,t) = u_x(\pi,t) = 0 & \qquad \text{untuk } \; t > 0, \\
u(x,0) = \frac{10x}{\pi} & \qquad \text{untuk } \; 0 < x < \pi .
\end{array}
\end{equation*}
\end{exercise}

\begin{exercise} \label{heat:oneto100exr}
Temukan solusi deret dari
\begin{equation*}
\begin{array}{ll}
u_t =  u_{xx} & \qquad \text{untuk } \; 0 < x < 1 \text{ dan } t > 0, \\
u(0,t) = 0 , \quad u(1,t) = 100 & \qquad \text{untuk } \; t > 0, \\
u(x,0) = \sin (\pi x) & \qquad \text{untuk } \; 0 < x < 1 .
\end{array}
\end{equation*}
Petunjuk: Gunakan fakta bahwa $u(x,t) = 100 x$ merupakan solusi yang memenuhi
$u_t = u_{xx}$, $u(0,t) = 0$, $u(1,t) = 100$.  Kemudian gunakan superposisi.
\end{exercise}

\begin{exercise}
Temukan solusi \emph{\myindex{suhu keadaan tunak}} sebagai fungsi
$x$ saja,
dengan mengambil $t \to
\infty$ dalam solusi dari
latihan \ref{heat:cosexr} dan \ref{heat:oneto100exr}.
Verifikasi bahwa solusi tersebut memenuhi persamaan $u_{xx} = 0$.
\end{exercise}

\begin{exercise}
Gunakan pemisahan variabel untuk menemukan solusi tak trivial
dari $u_{xx} + u_{yy} = 0$, dengan $u(x,0) = 0$ dan $u(0,y) = 0$.
Petunjuk: Coba $u(x,y) = X(x)Y(y)$.
\end{exercise}

\begin{exercise}[menantang]
Andaikan salah satu ujung kawat terisolasi (misalnya di $x=0$) dan
ujung lainnya dipertahankan pada suhu nol.  Artinya,
temukan solusi deret dari
\begin{equation*}
\begin{array}{ll}
u_t = k u_{xx} & \qquad \text{untuk } \; 0 < x < L \text{ dan } t > 0, \\
u_x(0,t) = u(L,t) = 0  & \qquad \text{untuk } \; t > 0, \\
u(x,0) = f(x) & \qquad \text{untuk } \; 0 < x < L .
\end{array}
\end{equation*}
Nyatakan setiap koefisien dalam deret melalui integral $f(x)$.
\end{exercise}

\begin{exercise}[menantang]
Andaikan kawat tersebut melingkar dan terisolasi, sehingga tidak ada ujung. 
Anda dapat membayangkannya sebagai sekadar menghubungkan kedua ujung dan 
memastikan bahwa solusinya menyambung di kedua ujung.
Artinya, temukan solusi deret dari
\begin{equation*}
\begin{array}{ll}
u_t = k u_{xx} & \qquad \text{untuk } \; 0 < x < L \text{ dan } t > 0, \\
u(0,t) = u(L,t) , \quad
u_x(0,t) = u_x(L,t)  & \qquad \text{untuk } \; t > 0, \\
u(x,0) = f(x) & \qquad \text{untuk } \; 0 < x < L .
\end{array}
\end{equation*}
Nyatakan setiap koefisien dalam deret melalui integral $f(x)$.
\end{exercise}

\begin{exercise}
Perhatikan kawat yang terisolasi pada kedua ujung, $L=1$, $k=1$,
dan $u(x,0) = \cos^2(\pi x)$.
\begin{tasks}
\task
Temukan solusi $u(x,t)$.  Petunjuk: suatu identitas trigonometri.
\task
Temukan suhu rata-rata.
\task
Pada awalnya variasi suhu adalah 1 (maksimum dikurangi minimum).
Temukan waktu ketika variasinya $\nicefrac{1}{2}$.
\end{tasks}
\end{exercise}

\setcounter{exercise}{100}

%u(x,t) = 
%\sum_{n=1}^\infty
%b_n
%u_n(x,t)
%=
%\sum_{n=1}^\infty
%b_n
%\sin \left( \frac{n \pi}{L} x \right)
%e^{\frac{-n^2 \pi^2}{L^2} k t} .

\begin{exercise}
Temukan solusi deret dari
\begin{equation*}
\begin{array}{ll}
u_t =  3 u_{xx} & \qquad \text{untuk } \; 0 < x < \pi \text{ dan } t > 0, \\
u(0,t) = u(\pi,t) = 0 & \qquad \text{untuk } \; t > 0, \\
u(x,0) = 5\sin (x) + 2\sin (5x) & \qquad \text{untuk } \; 0 < x < \pi .
\end{array}
\end{equation*}
\end{exercise}
\exsol{%
$u(x,t) = 
5
\sin (x)
\, e^{- 3 t}
+
2
\sin (5x)
\, e^{-75 t}$
}

%u(x,t) = 
%\frac{a_0}{2} + 
%\sum_{n=1}^\infty
%a_n
%u_n(x,t)
%=
%\frac{a_0}{2} + 
%\sum_{n=1}^\infty
%a_n
%\cos \left( \frac{n \pi}{L} x \right)
%\, e^{\frac{-n^2 \pi^2}{L^2} k t} .

\begin{exercise}
Temukan solusi deret dari
\begin{equation*}
\begin{array}{ll}
u_t =  0.1 u_{xx} & \qquad \text{untuk } \; 0 < x < \pi \text{ dan } t > 0, \\
u_x(0,t) = u_x(\pi,t) = 0 &  \qquad \text{untuk } \; t > 0, \\
u(x,0) = 1 + 2\cos (x) & \qquad \text{untuk } \; 0 < x < \pi .
\end{array}
\end{equation*}
\end{exercise}
\exsol{%
$u(x,t) = 
1 + 
2
\cos (x)
\, e^{-0.1 t}$
}

\begin{exercise}
Gunakan pemisahan variabel untuk menemukan solusi tak trivial dari
$u_{xt} = u_{xx}$.
\end{exercise}
\exsol{%
$u(x,t) = e^{\lambda t} e^{\lambda x}$ untuk suatu $\lambda$
}

\begin{exercise}
Gunakan suatu variasi pemisahan variabel 
untuk menemukan solusi tak trivial dari
$u_{x} + u_{t} = u$.  Petunjuk: Coba $u(x,t) = X(x)+T(t)$.
\end{exercise}
\exsol{%
$u(x,t) = Ae^x + Be^t$
}

\begin{exercise}
Andaikan suhu pada kawat ditetapkan sebesar $0$
di ujung-ujungnya, $L=1$, $k=1$, dan $u(x,0) = 100\sin(2 \pi x)$.
\begin{tasks}
\task
Berapakah suhu di $x = \nicefrac{1}{2}$ pada setiap waktu?
\task
Berapakah suhu maksimum dan minimum pada kawat
ketika $t=0$?
\task
Pada waktu berapakah suhu maksimum pada kawat tepat
setengah dari maksimum awal ketika $t=0$?
\end{tasks}
\end{exercise}
\exsol{%
a) $0$,
\quad b) minimum $-100$, maksimum $100$,
\quad c) $t = \frac{\ln 2}{4 \pi^2}$.
}

%%%%%%%%%%%%%%%%%%%%%%%%%%%%%%%%%%%%%%%%%%%%%%%%%%%%%%%%%%%%%%%%%%%%%%%%%%%%%%

\sectionnewpage
\section{Persamaan gelombang satu dimensi} \label{we:section}

\sectionnotes{1 kuliah\EPref{, \S9.6 dalam \cite{EP}}\BDref{,
\S10.7 dalam \cite{BD}}}

Bayangkan senar gitar tegang dengan panjang $L$ yang dapat bergetar.
Kita hanya akan meninjau getaran dalam satu arah.  Misalkan
$x$ menyatakan posisi sepanjang senar, $t$ menyatakan waktu,
dan $y(x,t)$
menyatakan simpangan senar dari posisi diam.
Lihat
\figurevref{we:vibstrfig}.

\begin{myfig}
\capstart
\inputpdft{sps-vibstr}{%
%perhatikan bahwa diagram ini muncul lebih dari sekali dalam buku
Diagram senar gitar (gitar berarsir berada di latar belakang).
Jembatan gitar berada di x sama dengan 0 pada sumbu x dan nut gitar
berada di x sama dengan L.  Senar terpasang pada sumbu x di kedua titik ini.
Senar tampak dipetik ke atas di tengah sehingga membentuk segitiga
lebar.  Sumbu vertikal diberi label y, dan sebuah panah menunjukkan bahwa
y memberikan simpangan senar pada setiap x tertentu, yaitu,
grafik y merupakan bentuk senar.}
\caption{Senar bergetar dengan panjang $L$, $x$ adalah posisi, $y$ adalah simpangan.\label{we:vibstrfig}}
\end{myfig}

Persamaan yang mengatur susunan ini adalah apa yang disebut
\emph{\myindex{persamaan gelombang satu dimensi}}\index{persamaan gelombang}:
\begin{equation*}
\mybxbg{~~
y_{tt} =
a^2 y_{xx} ,
~~}
\end{equation*}
untuk suatu konstanta $a > 0$.
Intuisinya mirip dengan persamaan panas, dengan mengganti kecepatan oleh
percepatan: Percepatan di suatu titik tertentu sebanding dengan turunan kedua
bentuk senar.  Dengan kata lain,
ketika senar
cekung ke bawah, $y_{xx}$ negatif dan senar cenderung berakselerasi
ke bawah, sehingga $y_{tt}$ seharusnya negatif.  Demikian pula sebaliknya.
Persamaan gelombang merupakan contoh PDP hiperbolik\@.

Kita kembali akan menyelesaikan $y$ di daerah $0 < x < L$ dan $t > 0$.
Andaikan ujung-ujung senar terpasang tetap seperti pada gitar:
\begin{equation*}
y(0,t) = 0 \quad \text{dan} \quad y(L,t) = 0 \quad \text{untuk } \; t > 0.
\end{equation*}
Kita memiliki dua syarat sepanjang sumbu $x$ karena terdapat
dua turunan dalam arah $x$.
Ada pula dua turunan sepanjang arah $t$ sehingga kita memerlukan
dua syarat tambahan di sini.  Kita perlu mengetahui posisi awal
dan kecepatan awal senar.  Artinya,
untuk fungsi-fungsi $f(x)$ dan $g(x)$ yang diketahui, kita menetapkan
\begin{equation*}
y(x,0) = f(x)  \quad \text{dan} \quad y_t (x,0) = g(x)
\quad \text{untuk } \; 0 < x < L .
\end{equation*}

Persamaan tersebut linear dan homogen,
sehingga superposisi berlaku sebagaimana untuk
persamaan panas.
Superposisi juga mempertahankan syarat samping homogen 
$y(0,t)=0$ dan $y(L,t)=0$.
Sekali lagi, kita akan menggunakan pemisahan variabel untuk menemukan
cukup banyak solusi blok pembangun guna memperoleh solusi khusus
yang juga menyelesaikan syarat awal tak homogen.  Ada
satu perubahan.  Kita akan menyelesaikan dua masalah terpisah
dan menjumlahkan solusinya.

Dua masalah yang akan kita selesaikan adalah
\begin{equation} \label{wave:weq}
\begin{array}{ll}
w_{tt} = a^2 w_{xx} & \qquad \text{untuk } \; 0 < x < L \text{ dan } t > 0, \\
w(0,t) = w(L,t) = 0 & \qquad \text{untuk } \; t > 0,  \\
w(x,0) = 0 & \qquad \text{untuk } \; 0 < x < L , \\
w_t(x,0) = g(x) & \qquad \text{untuk } \; 0 < x < L ,
\end{array}
\end{equation}
dan
\begin{equation} \label{wave:zeq}
\begin{array}{ll}
z_{tt} = a^2 z_{xx} & \qquad \text{untuk } \; 0 < x < L \text{ dan } t > 0, \\
z(0,t) = z(L,t) = 0 & \qquad \text{untuk } \; t > 0,  \\
z(x,0) = f(x) & \qquad \text{untuk } \; 0 < x < L , \\
z_t(x,0) = 0 & \qquad \text{untuk } \; 0 < x < L .
\end{array}
\end{equation}

Prinsip superposisi menyiratkan bahwa
$y = w + z$ menyelesaikan persamaan gelombang dan syarat samping
homogen.
Lebih lanjut,
$y(x,0) = w(x,0) + z(x,0) = f(x)$ dan
$y_t(x,0) = w_t(x,0) + z_t(x,0) = g(x)$.  Karena itu, $y$ merupakan
solusi untuk
\begin{equation} \label{wave:yeq}
\begin{array}{ll}
y_{tt} = a^2 y_{xx}  & \qquad \text{untuk } \; 0 < x < L \text{ dan } t > 0,  \\
y(0,t) = y(L,t) = 0  & \qquad \text{untuk } \; t > 0,  \\
y(x,0) = f(x) & \qquad \text{untuk } \; 0 < x < L , \\
y_t(x,0) = g(x) & \qquad \text{untuk } \; 0 < x < L .
\end{array}
\end{equation}

Alasan semua kerumitan ini adalah bahwa superposisi hanya berlaku untuk
syarat homogen seperti
$y(0,t) = y(L,t) = 0$, $y(x,0) = 0$, atau $y_t(x,0) = 0$.  Oleh karena itu,
kita dapat
menggunakan pemisahan variabel untuk menemukan banyak solusi blok pembangun
yang menyelesaikan semua syarat homogen.  Kemudian kita dapat menggunakannya untuk
mengonstruksi solusi yang memenuhi syarat tak homogen yang tersisa.

Mari kita mulai dengan \eqref{wave:weq}.
Kita kembali mencoba solusi berbentuk $w(x,t) = X(x) T(t)$.  Kita menyubstitusikannya ke dalam
persamaan gelombang untuk memperoleh
\begin{equation*}
X(x)T''(t) = a^2 X''(x) T(t) .
\end{equation*}
Dengan menulis ulang, kita memperoleh
\begin{equation*}
\frac{T''(t)}{a^2 T(t)} = \frac{X''(x)}{X(x)} .
\end{equation*}
Sekali lagi, ruas kiri hanya bergantung pada $t$ dan ruas kanan hanya bergantung
pada $x$.  Jadi kedua ruas sama dengan suatu konstanta, yang kita nyatakan dengan
$-\lambda$:
\begin{equation*}
\frac{T''(t)}{a^2 T(t)} = -\lambda = \frac{X''(x)}{X(x)} .
\end{equation*}
Kita menyelesaikannya untuk memperoleh dua persamaan diferensial biasa
\begin{align*}
X''(x) + \lambda X(x) &= 0 , \\
T''(t) + \lambda a^2 T(t) &= 0 .
\end{align*}
Syarat $0 = w(0,t) = X(0) T(t)$ menyiratkan $X(0) = 0$ dan
$w(L,t) = 0$ menyiratkan bahwa $X(L) = 0$.  Oleh karena itu, satu-satunya solusi
tak trivial untuk persamaan pertama ada ketika
$\lambda = \lambda_n = \frac{n^2 \pi^2}{L^2}$ dan solusinya adalah
\begin{equation*}
X_n(x) = \sin \left( \frac{n \pi}{L} x \right) .
\end{equation*}
Solusi umum untuk $T$ bagi $\lambda_n$ khusus ini adalah
\begin{equation*}
T_n(t) = A \cos \left( \frac{n \pi a}{L} t \right)
+ B \sin \left( \frac{n \pi a}{L} t \right).
\end{equation*}
Kita juga memiliki syarat $w(x,0) = 0$ atau $X(x)T(0) = 0$.  Ini
menyiratkan bahwa $T(0) = 0$, yang kemudian memaksa $A = 0$.  Lebih
mudah memilih $B=\frac{L}{n \pi a}$ (sebentar lagi Anda akan melihat alasannya)
sehingga
\begin{equation*}
T_n(t) = \frac{L}{n \pi a} \sin \left( \frac{n \pi a}{L} t \right).
\end{equation*}
Solusi blok pembangun kita adalah
\begin{equation*}
w_n(x,t) = 
\frac{L}{n \pi a} 
\sin \left( \frac{n \pi}{L} x \right)
\sin \left( \frac{n \pi a}{L} t \right) .
\end{equation*}
Kita menurunkannya terhadap $t$:
\begin{equation*}
\frac{\partial w_n}{\partial t}(x,t) = 
\sin \left( \frac{n \pi}{L} x \right)
\cos \left( \frac{n \pi a}{L} t \right) .
\end{equation*}
Karena itu,
\begin{equation*}
\frac{\partial w_n}{\partial t}(x,0) =
\sin \left( \frac{n \pi}{L} x \right) .
\end{equation*}
Kita mengekspansi $g(x)$ dalam sinus-sinus ini sebagai
\begin{equation*}
g(x) =
\sum_{n=1}^\infty b_n \sin \left( \frac{n \pi}{L} x \right) .
\end{equation*}
Dengan menggunakan superposisi,
kita menulis solusi untuk \eqref{wave:weq} sebagai deret
\begin{equation*}
w(x,t) =
\sum_{n=1}^\infty
b_n
w_n(x,t)
=
\sum_{n=1}^\infty
b_n
\frac{L}{n \pi a}
\sin \left( \frac{n \pi}{L} x \right)
\sin \left( \frac{n \pi a}{L} t \right) .
\end{equation*}

\begin{exercise}
Periksa bahwa $w(x,0) = 0$ dan
$w_t(x,0) = g(x)$.
\end{exercise}

Kita menyelesaikan \eqref{wave:zeq} dengan cara serupa.  Kita kembali mencoba
$z(x,t) = X(x)T(t)$.  Mula-mula prosedurnya berlaku dengan cara yang persis sama.
Kita memperoleh
\begin{align*}
X''(x) + \lambda X(x) &= 0 , \\
T''(t) + \lambda a^2 T(t) &= 0 ,
\end{align*}
dan syarat $X(0) = 0$, $X(L) = 0$.  Sekali lagi,
$\lambda = \lambda_n = \frac{n^2 \pi^2}{L^2}$ dan
\begin{equation*}
X_n(x) = \sin \left( \frac{n \pi}{L} x \right) .
\end{equation*}
Kali ini,
syarat pada $T$ adalah $T'(0) = 0$.  Dengan demikian 
kita memperoleh $B = 0$, dan kita mengambil
\begin{equation*}
T_n(t) = \cos \left( \frac{n \pi a}{L} t \right).
\end{equation*}
Solusi blok pembangun kita adalah
\begin{equation*}
z_n(x,t) = 
\sin \left( \frac{n \pi}{L} x \right)
\cos \left( \frac{n \pi a}{L} t \right) .
\end{equation*}
Karena $z_n(x,0) = \sin \left( \frac{n \pi}{L} x \right)$,
kita mengekspansi $f(x)$ dalam sinus-sinus ini sebagai
\begin{equation*}
f(x) =
\sum_{n=1}^\infty c_n \sin \left( \frac{n \pi}{L} x \right) .
\end{equation*}
Kita
menuliskan solusi untuk \eqref{wave:zeq} sebagai deret
\begin{equation*}
z(x,t) =
\sum_{n=1}^\infty
c_n
z_n(x,t)
=
\sum_{n=1}^\infty
c_n
\sin \left( \frac{n \pi}{L} x \right)
\cos \left( \frac{n \pi a}{L} t \right) .
\end{equation*}

\begin{exercise}
Lengkapi rincian dalam penurunan solusi \eqref{wave:zeq}.
Periksa bahwa solusi tersebut memenuhi semua syarat samping.
\end{exercise}

Dengan menggabungkan kedua solusi ini, mari kita nyatakan hasilnya sebagai teorema.
\begin{theorem}
Ambil masalah
\begin{equation} \label{wave:tyeq}
\begin{array}{ll}
y_{tt} = a^2 y_{xx} & \qquad \text{untuk } \; 0 < x < L \text{ dan } t > 0,  \\
y(0,t) = y(L,t) = 0 & \qquad \text{untuk } \; t > 0,  \\
y(x,0) = f(x) & \qquad \text{untuk } \; 0 < x < L , \\
y_t(x,0) = g(x) & \qquad \text{untuk } \; 0 < x < L ,
\end{array}
\end{equation}
dengan
\begin{equation*}
f(x) =
\sum_{n=1}^\infty c_n \sin \left( \frac{n \pi}{L} x \right)
\qquad \text{dan} \qquad
g(x) =
\sum_{n=1}^\infty b_n \sin \left( \frac{n \pi}{L} x \right) .
\end{equation*}
Maka solusi $y(x,t)$ dapat ditulis sebagai jumlah solusi
dari \eqref{wave:weq} dan \eqref{wave:zeq}:
\begin{equation*}
\mybxbg{~~
\begin{aligned}
y(x,t)
& =
\sum_{n=1}^\infty
\left[
b_n
\frac{L}{n \pi a}
\sin \left( \frac{n \pi}{L} x \right)
\sin \left( \frac{n \pi a}{L} t \right) 
+
c_n
\sin \left( \frac{n \pi}{L} x \right)
\cos \left( \frac{n \pi a}{L} t \right) 
\right]
\\
& =
\sum_{n=1}^\infty
\sin \left( \frac{n \pi}{L} x \right)
\left[
b_n
\frac{L}{n \pi a}
\sin \left( \frac{n \pi a}{L} t \right) 
+
c_n
\cos \left( \frac{n \pi a}{L} t \right) 
\right] .
\end{aligned}
~~}
\end{equation*}
\end{theorem}

\begin{example} \label{example:pluckedstring}
Perhatikan suatu senar sepanjang 2 yang dipetik di tengah.
Senar tersebut mempunyai bentuk awal yang diberikan dalam \figurevref{wave:pluckedstrfig}.
Yaitu,
\begin{equation*}
f(x) = \begin{cases}
0.1\, x & \text{jika } \; 0 \leq x \leq 1 , \\
0.1\, (2-x) & \text{jika } \; 1 < x \leq 2 .
\end{cases}
\end{equation*}

\begin{myfig}
\capstart
\inputpdft{wave-pluckedstr}{%
Grafik suatu fungsi pada bidang $xy$ yang 
terdiri atas dua ruas garis lurus; ruas pertama bermula pada $x$ sama dengan 0
di sumbu $x$, lalu mencapai maksimum 0.1 di tengah,
kemudian ruas kedua turun kembali ke sumbu $x$ pada $x$ sama dengan 2.}
\caption{Bentuk awal senar yang dipetik dari
\exampleref{example:pluckedstring}.\label{wave:pluckedstrfig}}
\end{myfig}

Misalkan senar mulai dari keadaan diam ($g(x) =
0$), dan ambil $a=1$ untuk kesederhanaan.  Dengan kata lain, kita ingin
menyelesaikan masalah:
\begin{equation*}
\begin{array}{ll}
y_{tt} = y_{xx} & \qquad \text{untuk } \; 0 < x < 2 \text{ dan } t > 0, \\
y(0,t) = y(2,t)= 0 & \qquad \text{untuk } \; t > 0 , \\
y(x,0) = f(x) & \qquad \text{untuk } \; 0 < x < 2 , \\
y_t(x,0)= 0 & \qquad \text{untuk } \; 0 < x < 2.
\end{array}
\end{equation*}

Perincian penghitungan deret sinus dari $f(x)$
kita serahkan kepada pembaca.  Deret tersebut adalah
\begin{equation*}
f(x) = \sum_{n=1}^\infty
\frac{0.8}{n^2 \pi^2}
\sin \left( \frac{n \pi}{2} \right)
\sin \left( \frac{n \pi}{2} x \right) .
\end{equation*}
Perhatikan bahwa 
$\sin \left( \frac{n \pi}{2} \right)$
adalah barisan $1, 0, -1, 0, 1, 0, -1, \ldots$
untuk $n = 1,2,3,4,\ldots$.  Oleh karena itu,
\begin{equation*}
f(x) = 
\frac{0.8}{\pi^2}
\sin \left( \frac{\pi}{2} x \right)
-
\frac{0.8}{9 \pi^2}
\sin \left( \frac{3 \pi}{2} x \right)
+
\frac{0.8}{25 \pi^2}
\sin \left( \frac{5 \pi}{2} x \right)
- \cdots
\end{equation*}
Solusi $y(x,t)$ diberikan oleh
\begin{equation*}
\begin{split}
y(x,t) & = 
\sum_{n=1}^\infty
\frac{0.8}{n^2 \pi^2}
\sin \left( \frac{n \pi}{2} \right)
\sin \left( \frac{n \pi}{2} x \right)
\cos \left( \frac{n \pi}{2} t \right)
\\
& = 
\sum_{m=1}^\infty
\frac{0.8 {(-1)}^{m+1}}{{(2m-1)}^2 \pi^2}
\sin \left( \frac{(2m-1) \pi}{2} x \right)
\cos \left( \frac{(2m-1) \pi}{2} t \right)
\\
& =
\frac{0.8}{\pi^2} 
\sin \left( \frac{\pi}{2}  x \right)
\cos \left( \frac{\pi}{2}  t \right)
-
\frac{0.8}{9 \pi^2} 
\sin \left( \frac{3 \pi}{2}  x \right)
\cos \left( \frac{3 \pi}{2}  t \right)
\\
& \hspace{15em}
+
\frac{0.8}{25 \pi^2}
\sin \left( \frac{5 \pi}{2}  x \right)
\cos \left( \frac{5 \pi}{2}  t \right) 
- \cdots
\end{split}
\end{equation*}

Lihat 
\figurevref{wave:pluckedexfig} untuk grafik
bagi $0 < t < 3$.  Perhatikan
bahwa tidak seperti persamaan panas, solusinya tidak menjadi
\myquote{lebih mulus;}
\myquote{sudut-sudut tajamnya} tetap ada.  Kita akan melihat alasan perilaku ini dalam
bagian berikutnya, ketika kita menurunkan solusi persamaan gelombang dengan cara
yang berbeda.

\begin{myfig}
\capstart
\diffyincludegraphics{width=5in}{width=7.5in}{wave-pluckedex}{%
Grafik fungsi $y$ dari $(x,t)$ di atas bidang $xt$
pada daerah dengan $x$ di antara 0 dan 2 serta $t$ di antara 0
dan 3.  Pada $t$ sama dengan 0, fungsi tersebut berupa segitiga menghadap ke atas yang mencapai 0.1 di
tengah.  Permukaan bernilai 0 pada sisi-sisi tempat $x$ sama dengan 0 atau $x$ sama dengan 1.
Permukaan tampak seperti setengah piramida menghadap ke atas yang berakhir pada $t$ sama dengan 1,
kemudian terdapat piramida terbalik menghadap ke bawah yang mencapai minus 0.1
untuk $t$ di antara 1 dan 3.}
\caption{Bentuk senar yang dipetik untuk $0 < t < 3$.\label{wave:pluckedexfig}}
\end{myfig}

Pastikan Anda memahami apa yang disampaikan oleh grafik seperti yang ada dalam gambar.
Untuk setiap $t$ yang tetap, Anda dapat memandang fungsi 
$y(x,t)$ hanya sebagai fungsi dari $x$.  Fungsi ini memberikan bentuk
senar pada waktu $t$.  Lihat \figurevref{wave:pluckedtsfig} untuk grafik
$y$ sebagai fungsi dari $x$ pada beberapa nilai $t$ yang berbeda.
Pada grafik-grafik ini, Anda dapat melihat sudut-sudut tajam yang tetap ada dengan jauh lebih jelas.

\begin{myfig}
\capstart
%berkas asli wave-plucked-slicet0 wave-plucked-slicet0p4 wave-plucked-slicet0p8 wave-plucked-slicet1p2
\diffyincludegraphics{width=6.24in}{width=9in}{wave-plucked-t0-t0p4}{%
Dua grafik bentuk senar yang dipetik pada bidang $xy$,
dengan $x$ di antara 0 dan 2.
Di sebelah kiri adalah bentuk awal pada $t$ sama dengan 0, yakni segitiga menghadap ke atas dengan
kedua ujung pada 0 dan titik puncak di tengah pada 0.1.
Di sebelah kanan adalah bentuk pada $t$ sama dengan 0.4, yang awalnya sama seperti bentuk awal,
tetapi kini segitiganya terpotong pada bagian atas, sehingga berupa garis yang naik
dengan kemiringan sama seperti sebelumnya, kemudian suatu ruas horizontal, lalu garis
yang turun kembali ke 0 dengan kemiringan sama seperti sebelumnya.}
\\[5pt]
\diffyincludegraphics{width=6.24in}{width=9in}{wave-plucked-t0p8-t1p2}{%
Dua grafik bentuk senar yang dipetik pada bidang $xy$,
dengan $x$ di antara 0 dan 2.
Di sebelah kiri adalah bentuk pada $t$ sama dengan 0.8, yang mirip dengan bentuk pada $t$
sama dengan 0.4
di atas, tetapi kini ruas horizontal di bagian atas lebih panjang dan muncul
lebih awal.  Di sebelah kanan adalah bentuk pada $t$ sama dengan 1.2, yang merupakan versi terbalik
dari bentuk $t$ sama dengan 0.8 di sebelah kiri, yaitu, garis-garis miring sekarang berada
di bawah sumbu $x$.}
\caption{Senar yang dipetik untuk $t=0$, $t=0.4$, $t=0.8$, dan
$t=1.2$.%
\label{wave:pluckedtsfig}}
\end{myfig}
\end{example}

Satu hal yang dapat dipetik dari semua ini adalah cara sebuah gitar menghasilkan suara.  Perhatikan bahwa
frekuensi (sudut) yang muncul dalam solusi adalah
$n \frac{\pi a}{L}$.  Artinya, terdapat suatu
\emph{\myindex{frekuensi dasar}} $\frac{\pi a}{L}$, dan kemudian kita juga
memperoleh semua kelipatan frekuensi ini, yang dalam musik disebut
\emph{\myindex{nada atas}}.  Nada atas mana yang muncul dan dengan amplitudo berapa
merupakan apa yang oleh musisi disebut \emph{\myindex{warna nada}}.
Matematikawan biasanya menyebutnya \emph{\myindex{spektrum}}.
Karena semua frekuensi merupakan kelipatan dari satu frekuensi (frekuensi
dasar),
kita memperoleh bunyi yang enak dan menyenangkan.

Frekuensi dasar $\frac{\pi a}{L}$ meningkat ketika kita memperkecil panjang $L$.  Artinya, jika
kita meletakkan jari pada papan jari lalu memetik senar, kita memperoleh nada yang lebih
tinggi.  Konstanta $a$ diberikan oleh
\begin{equation*}
a = \sqrt{\frac{T}{\rho}} ,
\end{equation*}
di mana $T$ adalah tegangan dan $\rho$ adalah kerapatan linear senar.
Mengencangkan senar (memutar pasak penyetem pada gitar) meningkatkan $a$ dan
karena itu menghasilkan frekuensi dasar yang lebih tinggi (nada yang lebih tinggi).
Di sisi lain, menggunakan senar yang lebih berat 
mengurangi $a$ dan menghasilkan frekuensi dasar yang lebih rendah (nada yang lebih rendah).
Gitar bas mempunyai senar yang lebih panjang dan tebal, sedangkan ukulele mempunyai senar pendek
yang terbuat dari bahan lebih ringan.

Hal yang cukup menarik adalah kesimetrian yang hampir berlaku antara ruang dan waktu.
Dalam bentuk paling sederhananya, kita melihat kesimetrian ini dalam solusi
\begin{equation*}
\sin \left( \frac{n \pi}{L} x \right)
\sin \left( \frac{n \pi a}{L} t \right)  .
\end{equation*}
Kecuali konstanta $a$, solusi ini tampak sama jika kita
menukar waktu dan ruang.
Secara umum, solusi untuk $x$ tetap merupakan deret Fourier dalam $t$, sedangkan untuk
$t$ tetap, solusi itu merupakan deret Fourier dalam $x$, dan koefisien-koefisiennya saling berkaitan.
Jika bentuk $f(x)$ atau kecepatan awal $g(x)$ mempunyai banyak sudut, maka
gelombang bunyinya akan mempunyai banyak sudut.  Hal itu karena koefisien Fourier
dari bentuk awal meluruh ke nol (ketika $n \to \infty$) dengan laju yang sama seperti koefisien Fourier
dari gelombang terhadap waktu (untuk suatu $x$ tetap).  Jadi, jika Anda menggunakan benda tajam untuk
memetik senar, Anda memperoleh bunyi yang lebih tajam dengan banyak komponen
berfrekuensi tinggi, sedangkan jika Anda menggunakan ibu jari, Anda memperoleh bunyi lebih lembut tanpa
begitu banyak nada atas tinggi.
Demikian pula, jika Anda memetik dekat jembatan (dekat
salah satu ujung senar), petikan yang Anda hasilkan
lebih menyerupai gelombang gigi gergaji, dan Anda memperoleh bunyi yang bahkan
lebih tajam.

Sesungguhnya, jika Anda melihat rumus solusinya, Anda mendapati bahwa untuk setiap
$x$ tetap, kita memperoleh deret Fourier yang hampir sebarang dalam $t$, yakni semuanya
kecuali suku konstanta.  Secara teori, Anda dapat memperoleh warna nada apa pun yang Anda inginkan
dengan memetik senar tepat dengan cara yang sesuai.
Tentu saja, kita sedang meninjau senar ideal tanpa kekakuan dan tanpa hambatan
udara.  Variabel-variabel tersebut jelas juga memengaruhi bunyi.

\subsection{Latihan}

\begin{exercise}
Selesaikan
\begin{equation*}
\begin{array}{ll}
y_{tt} = 9 y_{xx}  & \qquad \text{untuk } \; 0 < x < 1 \text{ dan } t > 0,  \\
y(0,t) = y(1,t) = 0 & \qquad \text{untuk } \; t > 0 ,  \\
y(x,0) = \sin (3\pi x) + \frac{1}{4} \sin (6 \pi x) & \qquad \text{untuk } \; 0 < x < 1 , \\
y_t(x,0) = 0 & \qquad \text{untuk } \; 0 < x < 1 .
\end{array}
\end{equation*}
\end{exercise}

\begin{exercise}
Selesaikan
\begin{equation*}
\begin{array}{ll}
y_{tt} = 4 y_{xx} & \qquad \text{untuk } \; 0 < x < 1 \text{ dan } t > 0, \\
y(0,t) = y(1,t) = 0 & \qquad \text{untuk } \; t > 0,  \\
y(x,0) = \sin (3\pi x) + \frac{1}{4} \sin (6 \pi x) & \qquad \text{untuk } \; 0 < x < 1 , \\
y_t(x,0) = \sin (9 \pi x) & \qquad \text{untuk } \; 0 < x < 1 .
\end{array}
\end{equation*}
\end{exercise}

\begin{exercise}
Turunkan solusi untuk senar dipetik yang umum dengan panjang $L$ dan
sebarang konstanta $a$ (dalam persamaan $y_{tt} = a^2 y_{xx}$), dengan senar
kita angkat sejauh $b$ pada titik tengah lalu kita lepaskan.
\end{exercise}

\begin{samepage}
\begin{exercise}
Bayangkan sebuah alat musik berdawai jatuh ke lantai.  Andaikan
panjang senarnya adalah 1 dan $a=1$.  Ketika alat musik tersebut menghantam
lantai, senar berada pada posisi diam sehingga $y(x,0) = 0$.  Namun,
senar bergerak dengan suatu kecepatan ketika tumbukan terjadi ($t=0$),
misalnya $y_t(x,0) = -1$.  Carilah
solusi
$y(x,t)$ untuk bentuk senar pada waktu $t$.
\end{exercise}
\end{samepage}

\begin{exercise}[menantang]
Andaikan Anda mempunyai suatu senar bergetar dan
terdapat hambatan udara yang sebanding dengan kecepatan.  Artinya, Anda mempunyai
\begin{equation*}
\begin{array}{ll}
y_{tt} = a^2 y_{xx} - k y_t  & \qquad \text{untuk } \; 0 < x < 1 \text{ dan } t > 0,  \\
y(0,t) = y(1,t) = 0 & \qquad \text{untuk } \; t > 0, \\
y(x,0) = f(x) & \qquad \text{untuk } \; 0 < x < 1 , \\
y_t(x,0) = 0 & \qquad \text{untuk } \; 0 < x < 1 .
\end{array}
\end{equation*}
Andaikan $0 < k < 2 \pi a$.
Turunkan suatu solusi deret untuk masalah tersebut.  Setiap koefisien dalam deret
harus dinyatakan sebagai integral dari $f(x)$.
\end{exercise}

\begin{exercise}
Andaikan Anda menyentuh senar gitar tepat di tengah untuk
memastikan syarat lain $y(\nicefrac{L}{2},t) = 0$ untuk semua waktu.
Kelipatan mana dari frekuensi dasar $\frac{\pi a}{L}$
yang muncul dalam solusi?
\end{exercise}

\setcounter{exercise}{100}

\begin{exercise}
Selesaikan
\begin{equation*}
\begin{array}{ll}
y_{tt} = y_{xx}  & \qquad \text{untuk } \; 0 < x < \pi, t > 0,  \\
y(0,t) = y(\pi,t) = 0 & \qquad \text{untuk } \; t > 0,  \\
y(x,0) = \sin(x) & \qquad \text{untuk } \; 0 < x < \pi , \\
y_t(x,0) = \sin(x) & \qquad \text{untuk } \; 0 < x < \pi .
\end{array}
\end{equation*}
\end{exercise}
\exsol{%
$
y(x,t)
=
\sin(x)
\bigl(\sin(t) + \cos(t)\bigr)
$
}

\begin{exercise}
Selesaikan
\begin{equation*}
\begin{array}{ll}
y_{tt} = 25 y_{xx} & \qquad \text{untuk } \; 0 < x < 2 \text{ dan } t > 0,  \\
y(0,t) = y(2,t) = 0 & \qquad \text{untuk } \; t > 0, \\
y(x,0) = 0 & \qquad \text{untuk } \; 0 < x < 2 , \\
y_t(x,0) = \sin(\pi x) + 0.1 \sin(2\pi x) & \qquad \text{untuk } \; 0 < x < 2 .
\end{array}
\end{equation*}
\end{exercise}
\exsol{%
$y(x,t)
=
\frac{1}{5 \pi}
\sin (\pi x)
\sin (5 \pi t)
+
\frac{1}{100 \pi}
\sin(2\pi x)
\sin(10\pi t)$
}
%$
%y(x,t)
%=
%\sum\limits_{n=1}^\infty
%\sin \left( \frac{n \pi}{L} x \right)
%\left[
%b_n
%\frac{L}{n \pi a}
%\sin \left( \frac{n \pi a}{L} t \right) 
%+
%c_n
%\cos \left( \frac{n \pi a}{L} t \right) 
%\right] .

\begin{exercise}
Selesaikan
\begin{equation*}
\begin{array}{ll}
y_{tt} = 2 y_{xx} & \qquad \text{untuk } \; 0 < x < \pi \text{ dan } t > 0,  \\
y(0,t) = y(\pi,t) = 0 & \qquad \text{untuk } \; t > 0 , \\
y(x,0) = x & \qquad \text{untuk } \; 0 < x < \pi , \\
y_t(x,0) = 0 & \qquad \text{untuk } \; 0 < x < \pi .
\end{array}
\end{equation*}
\end{exercise}
\exsol{%
$
y(x,t)
=
\sum\limits_{n=1}^\infty
\frac{2{(-1)}^{n+1}}{n}
\sin(nx)
\cos( n \sqrt{2}\,t ) 
$
%$
%y(x,t)
%=
%\sum\limits_{n=1}^\infty
%\sin \left( \frac{n \pi}{L} x \right)
%\left[
%b_n
%\frac{L}{n \pi a}
%\sin \left( \frac{n \pi a}{L} t \right) 
%+
%c_n
%\cos \left( \frac{n \pi a}{L} t \right) 
%$
}

\begin{exercise}
Apa yang terjadi ketika $a=0$?  Carilah solusi untuk
$y_{tt} = 0$, $y(0,t) = y(\pi,t) = 0$,
$y(x,0) = \sin(2x)$,
$y_t(x,0) = \sin(x)$.
\end{exercise}
\exsol{%
$y(x,t) = \sin(2x)+t\sin(x)$
}

%%%%%%%%%%%%%%%%%%%%%%%%%%%%%%%%%%%%%%%%%%%%%%%%%%%%%%%%%%%%%%%%%%%%%%%%%%%%%%

\sectionnewpage
\section{Solusi d'Alembert untuk persamaan gelombang}
\label{dalemb:section}

%mbxINTROSUBSECTION

\sectionnotes{1 kuliah\EPref{, berbeda dari \S9.6 dalam \cite{EP}}\BDref{,
bagian dari \S10.7 dalam \cite{BD}}}

Kita telah menyelesaikan persamaan gelombang dengan menggunakan deret Fourier.  Namun, sering kali
lebih mudah menggunakan apa yang disebut
\emph{\myindex{solusi d'Alembert untuk persamaan gelombang}}%
\footnote{Dinamai menurut matematikawan Prancis
\href{https://en.wikipedia.org/wiki/D\%27Alembert}{Jean le Rond d'Alembert}
(1717--1783).}.
Walaupun solusi ini juga dapat diturunkan dengan menggunakan deret Fourier,
hal itu sebenarnya merupakan penggunaan konsep-konsep tersebut yang canggung.  Cara yang lebih mudah dan lebih
instruktif adalah menurunkan
solusi ini dengan melakukan perubahan variabel yang tepat untuk memperoleh persamaan yang
dapat diselesaikan melalui integrasi sederhana.

Andaikan kita ingin menyelesaikan \myindex{persamaan gelombang}
\begin{equation} \label{dalemb:weq}
y_{tt} = a^2 y_{xx}
\end{equation}
pada daerah $0 < x < L$ dan $t > 0$
dengan syarat-syarat samping
\begin{equation} \label{dalemb:weqside}
\begin{array}{ll}
y(0,t) =  y(L,t) = 0 &
\qquad \text{untuk } \; t > 0 , \\
y(x,0) = f(x) & \qquad \text{untuk } \; 0 < x < L , \\
y_t(x,0) = g(x) &  \qquad \text{untuk } \; 0 < x < L .
\end{array}
\end{equation}

\subsection{Perubahan variabel}

Kita akan mentransformasikan persamaan tersebut ke bentuk yang lebih sederhana sehingga dapat diselesaikan melalui
integrasi sederhana.
Kita mengubah variabel menjadi $\xi = x - at$, $\eta = x + at$.
Aturan rantai menyatakan:
\begin{align*}
& \frac{\partial}{\partial x}
=
\frac{\partial \xi}{\partial x}
\frac{\partial}{\partial \xi}
+
\frac{\partial \eta}{\partial x}
\frac{\partial}{\partial \eta}
=
\frac{\partial}{\partial \xi}
+
\frac{\partial}{\partial \eta} , \\
& \frac{\partial}{\partial t}
=
\frac{\partial \xi}{\partial t}
\frac{\partial}{\partial \xi}
+
\frac{\partial \eta}{\partial t}
\frac{\partial}{\partial \eta}
=
-a
\frac{\partial}{\partial \xi}
+
a
\frac{\partial}{\partial \eta} .
\end{align*}
Kita hitung
\begin{align*}
& y_{xx} = \frac{\partial^2 y}{\partial x^2}
=
\left(
\frac{\partial}{\partial \xi}
+
\frac{\partial}{\partial \eta}
\right)
\left(
\frac{\partial y}{\partial \xi}
+
\frac{\partial y}{\partial \eta}
\right)
=
\frac{\partial^2 y}{\partial \xi^2}
+
2 \frac{\partial^2 y}{\partial \xi \partial \eta}
+
\frac{\partial^2 y}{\partial \eta^2} ,
\\
& y_{tt} = \frac{\partial^2 y}{\partial t^2}
=
\left(
-a
\frac{\partial}{\partial \xi}
+ a
\frac{\partial}{\partial \eta}
\right)
\left(
-a
\frac{\partial y}{\partial \xi}
+
a
\frac{\partial y}{\partial \eta}
\right)
=
a^2
\frac{\partial^2 y}{\partial \xi^2}
-
2 a^2 \frac{\partial^2 y}{\partial \xi \partial \eta}
+
a^2
\frac{\partial^2 y}{\partial \eta^2} .
\end{align*}
Dalam penghitungan di atas, kita menggunakan fakta dari kalkulus bahwa
$\frac{\partial^2 y}{\partial \xi \partial \eta} = 
\frac{\partial^2 y}{\partial \eta \partial \xi}$.
Kita substitusikan hasil yang diperoleh ke persamaan gelombang,
\begin{equation*}
0 = a^2 y_{xx} - y_{tt} =
4 a^2 \frac{\partial^2 y}{\partial \xi \partial \eta} = 4 a^2 y_{\xi\eta} .
\end{equation*}
Oleh karena itu, persamaan gelombang \eqref{dalemb:weq} berubah menjadi
$y_{\xi\eta} = 0$.
Solusi umum persamaan baru ini mudah ditemukan dengan mengintegralkan
dua kali.  Dengan mempertahankan $\xi$ konstan, pertama-tama kita mengintegralkan terhadap $\eta$%
\footnote{Tidak ada yang istimewa mengenai $\eta$; Anda dapat lebih dahulu mengintegralkan terhadap
$\xi$, jika menghendakinya.}
dan memperhatikan bahwa
konstanta integrasi bergantung pada $\xi$; untuk setiap $\xi$, kita dapat memperoleh
konstanta integrasi yang berbeda.  Kita memperoleh
$y_{\xi} = C(\xi)$.
Berikutnya, kita mengintegralkan terhadap $\xi$ dan memperhatikan bahwa konstanta
integrasi bergantung pada $\eta$.
Dengan demikian,
$y = \int C(\xi) \, d\xi + B(\eta)$.
Oleh karena itu, solusi harus berbentuk berikut untuk suatu fungsi
$A(\xi)$ dan $B(\eta)$:
\begin{equation*}
y = A(\xi) + B(\eta) = A(x-at) + B(x+at) .
\end{equation*}
Solusi tersebut merupakan superposisi dua fungsi (gelombang) yang merambat dengan laju
$a$
ke arah yang berlawanan.  Koordinat $\xi$ dan $\eta$ disebut
\emph{\myindex{koordinat karakteristik}}, dan teknik serupa dapat
diterapkan pada PDE hiperbolik yang lebih rumit.
Dalam \sectionref{fopde:section}, teknik ini digunakan untuk menyelesaikan
PDE linear orde pertama.
Pada dasarnya, untuk menyelesaikan persamaan gelombang (atau persamaan
hiperbolik yang lebih umum), kita mencari kurva karakteristik tertentu yang sepanjang kurva tersebut
persamaannya sebenarnya hanya berupa ODE\@, atau sepasang ODE.  Dalam kasus ini,
kurva-kurva tersebut adalah kurva tempat $\xi$ dan $\eta$ konstan.

\subsection{Rumus d'Alembert}

Kita mengetahui bentuk setiap solusi, tetapi kita perlu menyelesaikannya untuk
syarat-syarat samping yang diberikan.  Kita cukup memberikan rumusnya dan melihat bahwa rumus itu berlaku.
Misalkan $F(x)$
menyatakan perluasan periodik ganjil dari $f(x)$, dan misalkan $G(x)$ menyatakan
perluasan periodik ganjil dari $g(x)$.  Definisikan
\begin{equation*}
A(x) = \frac{1}{2} F(x) - \frac{1}{2a} \int_0^x G(s) \,ds ,
\qquad
B(x) = \frac{1}{2} F(x) + \frac{1}{2a} \int_0^x G(s) \,ds .
\end{equation*}
Kita mengklaim bahwa $A(x)$ dan $B(x)$ ini memberikan solusi.  Secara eksplisit,
solusinya adalah $y(x,t) = A(x-at) + B(x+at)$, atau dengan kata lain:
\begin{equation} \label{dalemb:form}
\mybxbg{~~
\begin{aligned}
y(x,t) & =
\frac{1}{2} F(x-at) - \frac{1}{2a} \int_0^{x-at} G(s) \,ds 
+
\frac{1}{2} F(x+at) + \frac{1}{2a} \int_0^{x+at} G(s) \,ds \\
& =
\frac{F(x-at) + F(x+at)}{2} + \frac{1}{2a} \int_{x-at}^{x+at} G(s) \,ds .
\end{aligned}
~~}
\end{equation}

Mari kita periksa bahwa rumus d'Alembert benar-benar berlaku.
Pertama,
\begin{equation*}
y(x,0) =
\frac{F(x) + F(x)}{2} + \frac{1}{2a} \int_{x}^{x} G(s) \,ds 
=
F(x) .
\end{equation*}
Sejauh ini baik.  Untuk kesederhanaan, asumsikan $F$ terdiferensialkan.
Kita menggunakan bentuk pertama dari \eqref{dalemb:form} karena bentuk tersebut
lebih mudah didiferensialkan.
Menurut teorema dasar kalkulus, kita mempunyai
\begin{equation*}
y_t(x,t) =
\frac{-a}{2} F'(x-at) + \frac{1}{2} G(x-at)
+
\frac{a}{2} F'(x+at) + \frac{1}{2} G(x+at) .
\end{equation*}
Jadi,
\begin{equation*}
y_t(x,0) =
\frac{-a}{2} F'(x) + \frac{1}{2} G(x)
+
\frac{a}{2} F'(x) + \frac{1}{2} G(x) = G(x) .
\end{equation*}
Bagus!  Sekarang semuanya berjalan lancar.  Baiklah, kini syarat-syarat batas.  Perhatikan
bahwa $F(x)$ dan $G(x)$ ganjil.  Jadi,
\begin{equation*}
%mbxlatex \begin{split}
y(0,t)
%mbxlatex &
=
\frac{F(-at) + F(at)}{2} + \frac{1}{2a} \int_{-at}^{at} G(s) \,ds 
%mbxlatex \\
%mbxlatex &
=
\frac{-F(at) + F(at)}{2} + \frac{1}{2a} \int_{-at}^{at} G(s) \,ds 
= 0 + 0 = 0.
%mbxlatex \end{split}
\end{equation*}
Sekarang $F(x)$ ganjil dan periodik-$2L$, sehingga
\begin{equation*}
F(L-at)+F(L+at)
= 
F(-L-at)+F(L+at)
=
-F(L+at)+F(L+at)
= 0 .
\end{equation*}
Berikutnya, $G(s)$ ganjil dan periodik-$2L$, sehingga kita melakukan perubahan
variabel $v = s-L$.  Kemudian kita perhatikan bahwa $G(v+L)=G(v-L)=-G(-v+L)$,
sehingga $G(v+L)$ ganjil sebagai fungsi dari $v$:
\begin{equation*}
\int_{L-at}^{L+at}
G(s)\,ds
=
\int_{-at}^{at}
G(v+L)\,dv
=
0 .
\end{equation*}
Oleh karena itu,
\begin{equation*}
y(L,t) =
\frac{F(L-at) + F(L+at)}{2} + \frac{1}{2a} \int_{L-at}^{L+at} G(s) \,ds 
=
0 + 0  = 0.
\end{equation*}
Dan voil\`a, rumus tersebut berlaku.

\begin{example} \label{dalemb:implusexample}
D'Alembert mengatakan bahwa solusi merupakan superposisi dari
dua fungsi (gelombang) yang bergerak ke arah berlawanan dengan \myquote{laju} $a$.
Untuk memperoleh gambaran tentang
cara kerjanya, kita uraikan contoh berikut.
Perhatikan susunan yang lebih sederhana
\begin{equation*}
\begin{array}{ll}
y_{tt} = y_{xx} & \qquad \text{untuk } \; 0 < x < 1 \text{ dan } t > 0, \\
y(0,t) = y(1,t) = 0 & \qquad \text{untuk } \; t > 0, \\
y(x,0) = f(x)  & \qquad \text{untuk } \; 0 < x < 1, \\
y_t(x,0) = 0   & \qquad \text{untuk } \; 0 < x < 1 .
\end{array}
\end{equation*}
Misalkan $f(x)$ adalah impuls segitiga berikut dengan tinggi 1 yang berpusat pada $x=0.5$:
\begin{equation*}
f(x) =
\begin{cases}
0 & \text{jika } \; \phantom{0.5}0 \leq x < 0.45, \\
20\,(x-0.45) & \text{jika } \; 0.45 \leq x < 0.5, \\
20\,(0.55-x) & \text{jika } \; \phantom{5}0.5 \leq x < 0.55, \\
0 & \text{jika } \; 0.55 \leq x \leq 1 .
\end{cases}
\end{equation*}
Grafik impuls ini adalah grafik kiri atas dalam
\figurevref{dalemb:impulsfig}.

Misalkan $F(x)$ adalah perluasan periodik ganjil dari $f(x)$.  Maka
\eqref{dalemb:form} menyatakan bahwa
solusinya adalah
\begin{equation*}
y(x,t) = \frac{F(x-t) + F(x+t)}{2} .
\end{equation*}
Tidak sulit menghitung nilai-nilai tertentu dari $y(x,t)$.
Sebagai contoh, untuk
menghitung $y(0.1,0.6)$, kita perhatikan bahwa $x-t = -0.5$ dan $x+t = 0.7$.
Sekarang $F(-0.5) =
-f(0.5) = - 20\,(0.55 - 0.5) = -1$
dan $F(0.7) = f(0.7) = 0$.  Oleh karena itu,
$y(0.1,0.6) = \frac{-1 + 0}{2} = -0.5$.  Seperti dapat Anda lihat, solusi d'Alembert
jauh lebih mudah dihitung dan digambarkan daripada solusi deret Fourier.
Lihat \figurevref{dalemb:impulsfig} untuk grafik solusi $y$
bagi beberapa $t$ yang berbeda.
\begin{myfig}
\capstart
%berkas asli dalemb-impuls1 dalemb-impuls2 dalemb-impuls3 dalemb-impuls4
\diffyincludegraphics{width=6.24in}{width=9in}{dalemb-impuls-1-2}{%
Dua grafik bentuk solusi persamaan gelombang pada
bidang $xy$ untuk nilai $t$ tetap yang berbeda, dengan $x$ di antara 0 dan 1.
Di sebelah kiri adalah bentuk awal pada $t$ sama dengan 0, yang tersusun atas
ruas-ruas garis dan berupa ruas horizontal pada $y$ sama dengan 0 hingga hampir mencapai tengah
grafik, lalu naik dan
turun membentuk puncak tajam hingga $y$ sama dengan 1.  Setelah itu, grafik berupa ruas
horizontal pada $y$ sama dengan 0 hingga ujung kanan.
Di sebelah kanan adalah gambar serupa untuk bentuk pada $t$ sama dengan 0.2,
tetapi kini terdapat dua puncak yang lebih kecil,
yang hanya naik hingga 0.5, masing-masing pada satu sisi (secara simetris) dari pusat.}
\\[5pt]
\diffyincludegraphics{width=6.24in}{width=9in}{dalemb-impuls-3-4}{%
Dua grafik bentuk solusi persamaan gelombang pada
bidang $xy$ untuk nilai $t$ tetap yang berbeda, dengan $x$ di antara 0 dan 1.
Di sebelah kiri adalah bentuk pada $t$ sama dengan 0.4, yang mirip
dengan bentuk pada $t$ sama dengan 0.2 di atas, tetapi kini kedua puncak lebih dekat ke sisi-sisi.
Puncak-puncak tersebut masih naik hingga $y$ sama dengan 0.5.
Di sebelah kanan adalah bentuk pada $t$ sama dengan 0.6, yang merupakan bentuk terbalik dari
gambar pada $t$ sama dengan 0.4, yaitu,
puncak-puncaknya kini mengarah turun hingga $y$ sama dengan minus 0.5.}
\caption{Grafik solusi d'Alembert untuk $t=0$, $t=0.2$, $t=0.4$, dan
$t=0.6$.%
\label{dalemb:impulsfig}}
\end{myfig}
\end{example}

\subsection{Cara lain menyelesaikan syarat-syarat samping}

Barangkali prosedurnya lebih mudah dan lebih berguna
untuk diingat daripada rumusnya sendiri dihafalkan.
Hal yang penting adalah bahwa suatu solusi persamaan
gelombang merupakan superposisi dua gelombang yang merambat ke arah berlawanan.
Yaitu,
\begin{equation*}
y(x,t) = A(x-at) + B(x+at) .
\end{equation*}
Jika Anda memikirkannya, rumus eksak untuk $A$ dan $B$ tidak sulit
ditebak setelah Anda menyadari syarat samping seperti apa yang harus
dipenuhi $y(x,t)$.  Mari kita telusuri kembali rumus tersebut, tetapi dengan cara yang agak berbeda.
Pendekatan terbaik adalah melakukannya secara bertahap.  Ketika $g(x) = 0$ (dan karenanya
$G(x) = 0$), solusinya adalah
\begin{equation*}
\frac{ F(x-at) + F(x+at) }{2} .
\end{equation*}
Di sisi lain,
ketika $f(x) = 0$ (dan karenanya $F(x) = 0$), kita definisikan
\begin{equation*}
H(x) = \int_0^x G(s) \,ds .
\end{equation*}
Solusi dalam kasus ini adalah
\begin{equation*}
\frac{1}{2a} \int_{x-at}^{x+at} G(s) \,ds
=
\frac{ -H(x-at) + H(x+at) }{2a} .
\end{equation*}
Dengan superposisi, kita memperoleh suatu solusi untuk syarat samping umum
\eqref{dalemb:weqside} (ketika baik $f(x)$ maupun $g(x)$ tidak identik dengan nol).
\begin{equation} \label{dalemb:altform}
y(x,t) = \frac{ F(x-at) + F(x+at) }{2} +
\frac{ -H(x-at) + H(x+at) }{2a} .
\end{equation}
Perhatikan tanda minus di depan $H$ dan $a$ pada penyebut kedua.

\begin{exercise}
Periksa bahwa rumus baru \eqref{dalemb:altform} memenuhi
syarat-syarat samping
\eqref{dalemb:weqside}.
\end{exercise}

\textbf{Peringatan:}
Pastikan Anda menggunakan perluasan periodik ganjil $F(x)$ dan $G(x)$
ketika Anda mempunyai rumus untuk $f(x)$ dan $g(x)$.
Masalahnya, rumus-rumus tersebut secara umum hanya berlaku
untuk $0 < x < L$ dan biasanya tidak sama dengan $F(x)$ dan $G(x)$
untuk $x$ lainnya.

\subsection{Beberapa catatan}

Kita mencatat bahwa rumus $y(x,t) = A(x-at) + B(x+at)$ merupakan alasan
mengapa solusi persamaan gelombang tidak menjadi \myquote{lebih bagus} seiring
berjalannya waktu, yaitu, mengapa dalam contoh-contoh dengan syarat awal
yang mempunyai sudut, solusinya juga mempunyai sudut pada setiap waktu $t$.

\medskip

Sudut-sudut tersebut membawa kita pada catatan menarik lainnya.  Pada awalnya, tidak seorang pun menyadari
bahwa solusi contoh kita bahkan tidak terdiferensialkan (solusi tersebut mempunyai sudut):
Dalam \exampleref{dalemb:implusexample} di atas, solusinya tidak
terdiferensialkan setiap kali $x=t+0.5$ atau $x=-t+0.5$, misalnya.
Sesungguhnya, untuk dapat menghitung $y_{xx}$ atau $y_{tt}$, Anda tidak memerlukan satu, melainkan dua
turunan.  Jangan khawatir, kita dapat memikirkan suatu bentuk yang sangat mendekati
$F(x)$ tetapi memang mempunyai dua turunan dengan sedikit membulatkan sudut-sudutnya,
dan kemudian solusinya akan sangat mendekati
$\frac{F(x-t)+F(x+t)}{2}$ dan tidak seorang pun akan menyadari penggantian tersebut.

\medskip

Satu catatan terakhir adalah apa yang diberitahukan solusi d'Alembert kepada kita mengenai
bagian-bagian syarat awal yang memengaruhi solusi pada suatu titik tertentu.
Kita dapat mengetahui hal ini dengan \myquote{bergerak mundur sepanjang
karakteristik.}  Untuk kesederhanaan, andaikan senarnya sangat panjang (mungkin
tak hingga).  Karena solusi pada waktu $t$ adalah
\begin{equation*}
y(x,t) =
\frac{F(x-at) + F(x+at)}{2} + \frac{1}{2a} \int_{x-at}^{x+at} G(s) \,ds ,
\end{equation*}
kita memperhatikan bahwa kita hanya menggunakan syarat awal pada interval
$[x-at,x+at]$.  Titik-titik ujung interval ini disebut
\emph{\myindex{muka gelombang}}, sebab di sanalah muka gelombang berada, jika diberikan
suatu gangguan awal ($t=0$) pada $x$.
Jika $a=1$, seorang pengamat yang berada pada $x=0$ pada waktu $t=1$ hanya melihat
syarat awal untuk $x$ dalam interval $[-1,1]$
dan sama sekali tidak mengetahui hal lainnya.
Inilah sebabnya, sebagai contoh, kita tidak mengetahui bahwa sebuah supernova telah terjadi di
alam semesta sampai kita melihat cahayanya, jutaan tahun setelah
peristiwa tersebut terjadi.

\subsection{Latihan}

\begin{exercise}
Dengan menggunakan solusi d'Alembert, selesaikan $y_{tt} = 4y_{xx}$, $0 < x < \pi$, $t >
0$,
$y(0,t) = y(\pi, t) = 0$, $y(x,0) = \sin x$, dan
$y_t(x,0) = \sin x$.  Petunjuk: Perhatikan bahwa $\sin x$ adalah perluasan periodik ganjil dari
$y(x,0)$ dan $y_t(x,0)$.
\end{exercise}

\begin{exercise}
Dengan menggunakan solusi d'Alembert, selesaikan $y_{tt} = 2y_{xx}$, $0 < x < 1$, $t > 0$,
$y(0,t) = y(1, t) = 0$, $y(x,0) = \sin^5 (\pi x)$, dan
$y_t(x,0) = \sin^3 (\pi x)$.
\end{exercise}

\begin{exercise}
Ambil
$y_{tt} = 4y_{xx}$, $0 < x < \pi$, $t > 0$,
$y(0,t) = y(\pi, t) = 0$, $y(x,0) = x(\pi-x)$, dan
$y_t(x,0) = 0$.
\begin{tasks}
\task Selesaikan dengan menggunakan rumus d'Alembert.  Petunjuk: Anda dapat menggunakan deret sinus
untuk $y(x,0)$.
\task Carilah solusi sebagai fungsi dari $x$ untuk $t=0.5$, $t=1$, dan
$t=2$ yang tetap.  Jangan gunakan deret sinus di sini.
\end{tasks}
\end{exercise}

\begin{exercise}
Turunkan solusi d'Alembert untuk $y_{tt} = a^2 y_{xx}$, $0 < x < \pi$, $t >
0$,
$y(0,t) = y(\pi, t) = 0$, $y(x,0) = f(x)$, dan
$y_t(x,0) = 0$, dengan menggunakan solusi deret Fourier dari persamaan gelombang,
melalui penerapan suatu identitas trigonometri yang sesuai.
Petunjuk: Lakukan terlebih dahulu untuk satu suku dari solusi deret Fourier,
khususnya lakukan ketika $y$ adalah
$\sin(n  x )\cos(n a t )$.
\end{exercise}

\begin{exercise}
Solusi d'Alembert masih berlaku jika tidak ada syarat batas dan
syarat awal didefinisikan pada seluruh garis real.  Andaikan
$y_{tt} = y_{xx}$ (untuk semua $x$ pada garis real dan $t \geq 0$),
$y(x,0) = f(x)$, dan
$y_t(x,0) = 0$, dengan
\begin{equation*}
f(x) =
\begin{cases}
0 & \text{jika } \; \phantom{{-1} \leq {} }x < -1, \\
x+1 & \text{jika } \; {-1} \leq x < 0, \\
-x+1 & \text{jika } \; \phantom{-}0 \leq x < 1, \\
0 & \text{jika } \; \phantom{-}1 < x .
\end{cases}
\end{equation*}
Selesaikan dengan menggunakan solusi d'Alembert. Artinya, tuliskan definisi
bagian demi bagian untuk solusinya.  Kemudian sketsalah solusi untuk $t=0$,
$t=\nicefrac{1}{2}$, $t=1$, dan $t=2$.
\end{exercise}

\setcounter{exercise}{100}

\begin{exercise}
Dengan menggunakan solusi d'Alembert, selesaikan $y_{tt} = 9y_{xx}$, $0 < x < 1$, $t >
0$,
$y(0,t) = y(1, t) = 0$, $y(x,0) = \sin (2 \pi x)$, dan
$y_t(x,0) = \sin (3 \pi x)$.
\end{exercise}
\exsol{%
$y(x,t)=
\frac{\sin(2 \pi (x-3 t))+\sin(2 \pi (3 t+x))}{2}
+
\frac{\cos(3 \pi (x-3 t))-\cos(3 \pi (3 t+x))}{18\pi}$
}

\begin{exercise}
Ambil $y_{tt} = 4y_{xx}$, $0 < x < 1$, $t > 0$,
$y(0,t) = y(1, t) = 0$, $y(x,0) = x-x^2$, dan
$y_t(x,0) = 0$.  Dengan menggunakan solusi d'Alembert, carilah
solusi pada
\begin{tasks}(3)
\task $t=0.1$,
\task $t=\nicefrac{1}{2}$,
\task $t=1$.
\end{tasks}
Anda mungkin perlu membagi jawaban Anda menjadi beberapa kasus.
\end{exercise}
\exsol{%
a)
$y(x,0.1) =
\begin{cases}
%\frac{(x-0.2)-{(x-0.2)}^2+(x+0.2)-{(x+0.2)}^2}{2}
x-x^2-0.04
& \text{jika } \; 0.2 \leq x \leq 0.8 \\
%\frac{-(x+0.8)+{(x+0.8)}^2+(x+0.2)-{(x+0.2)}^2}{2}
0.6x
& \text{jika } \; x \leq 0.2 \\
%\frac{(x-0.2)-{(x-0.2)}^2-(x-0.8)+{(x-0.8)}^2}{2}
0.6-0.6x
& \text{jika } \; x \geq 0.8 \\
\end{cases}$
\\
b)
$y(x,\nicefrac{1}{2}) = -x+x^2$
\quad
c)
$y(x,1) = x-x^2$
}

\begin{exercise}
Ambil $y_{tt} = 100y_{xx}$, $0 < x < 4$, $t > 0$,
$y(0,t) = y(4, t) = 0$, $y(x,0) = F(x)$, dan
$y_t(x,0) = 0$.  Andaikan bahwa
$F(0)=0$,
$F(1)=2$,
$F(2)=3$,
$F(3)=1$.
Dengan menggunakan solusi d'Alembert, carilah
\begin{tasks}(3)
\task $y(1,1)$,
\task $y(4,3)$,
\task $y(3,9)$.
\end{tasks}
\end{exercise}
\exsol{%
a) $y(1,1) = -\nicefrac{1}{2}$
\quad
b) $y(4,3) = 0$
\quad
c) $y(3,9) = \nicefrac{1}{2}$
}

%%%%%%%%%%%%%%%%%%%%%%%%%%%%%%%%%%%%%%%%%%%%%%%%%%%%%%%%%%%%%%%%%%%%%%%%%%%%%%

\sectionnewpage
\section{Suhu keadaan tunak dan Laplasian}
\label{dirich:section}

\sectionnotes{1 kuliah\EPref{, \S9.7 dalam \cite{EP}}\BDref{,
\S10.8 dalam \cite{BD}}}

Perhatikan suatu kawat terisolasi, pelat, atau objek tiga dimensi.
Kita menerapkan
suhu-suhu tetap tertentu pada ujung-ujung kawat, tepi-tepi pelat,
atau semua sisi objek tiga dimensi tersebut.  Kita ingin mencari distribusi
\emph{\myindex{suhu keadaan tunak}}.
Artinya, kita ingin mengetahui
suhu setelah jangka waktu yang cukup lama.

Sesungguhnya kita mencari solusi persamaan panas yang tidak
bergantung pada waktu.
Pertama-tama, kita menyelesaikan masalah tersebut dalam satu variabel ruang.
Kita mencari suatu fungsi $u$ yang memenuhi
\begin{equation*}
u_t = k u_{xx} ,
\end{equation*}
tetapi sedemikian sehingga $u_t = 0$ untuk semua $x$ dan $t$.  Oleh karena itu, kita mencari
fungsi dari $x$ saja yang memenuhi $u_{xx} = 0$.  Persamaan ini mudah diselesaikan
dengan integrasi,
dan kita melihat bahwa $u = Ax+B$ untuk suatu konstanta $A$ dan $B$.

Perhatikan suatu kawat terisolasi dengan suhu konstan $T_1$ yang kita terapkan
pada salah satu ujung (misalnya tempat $x=0$) dan $T_2$ pada ujung lainnya (pada $x=L$, dengan $L$
sebagai panjang kawat).  Solusi keadaan tunak kita adalah
\begin{equation*}
u(x) = \frac{T_2-T_1}{L} x + T_1 .
\end{equation*}
Solusi tersebut hanyalah garis lurus dari satu ujung ke ujung lainnya.
Solusi ini sesuai dengan intuisi
akal sehat mengenai bagaimana panas seharusnya
terdistribusi dalam kawat.  Jadi, dalam satu dimensi, solusi keadaan tunak
hanyalah garis-garis lurus.

Situasinya lebih rumit dalam dua atau lebih dimensi ruang.
Untuk kesederhanaan, kita membatasi diri pada dua dimensi ruang.
Persamaan panas dalam dua
variabel ruang adalah
\begin{equation} \label{dirich:heateq}
u_t = k(u_{xx} + u_{yy}) ,
\end{equation}
atau lebih lazim ditulis sebagai
$u_t = k \Delta u$ atau
$u_t = k \nabla^2 u$.  Simbol $\Delta$ dan $\nabla^2$
keduanya berarti $\frac{\partial^2}{\partial x^2} +
\frac{\partial^2}{\partial y^2}$.  Di sini kita akan menggunakan $\Delta$.
Alasan penggunaan notasi tersebut adalah bahwa Anda
dapat mendefinisikan $\Delta$ sebagai hal yang tepat untuk sebarang jumlah dimensi
ruang, dan kemudian persamaan panas selalu berupa
$u_t = k \Delta u$.  Operator $\Delta$ disebut \emph{\myindex{Laplasian}}.

Baiklah, setelah urusan notasi selesai, mari kita lihat bentuk persamaan
untuk solusi keadaan tunak.  Kita mencari solusi
\eqref{dirich:heateq} yang tidak bergantung pada $t$, yaitu,
$u_t = 0$.
Oleh karena itu, kita mencari suatu
fungsi $u(x,y)$ sedemikian sehingga
\begin{equation*}
\mybxbg{~~
\Delta u = 
u_{xx} + u_{yy} = 0 .
~~}
\end{equation*}
Persamaan ini disebut \emph{\myindex{persamaan Laplace}}%
\footnote{Dinamai menurut matematikawan Prancis
\href{https://en.wikipedia.org/wiki/Laplace}{Pierre-Simon, marquis de Laplace}
(1749--1827).} dan merupakan contoh persamaan eliptik.
Solusi persamaan Laplace
disebut \emph{fungsi harmonik\index{fungsi harmonik}}
dan memiliki banyak sifat yang baik serta
penerapan yang jauh melampaui masalah panas keadaan tunak.

Fungsi harmonik dalam dua variabel tidak lagi hanya linear
(grafiknya tidak hanya berupa bidang).  Sebagai contoh, Anda dapat memeriksa bahwa fungsi
$x^2-y^2$ dan $xy$ harmonik.  Namun, perhatikan bahwa jika $u_{xx}$ positif, $u$ cekung
ke atas dalam arah $x$, maka $u_{yy}$ harus negatif dan $u$ harus
cekung ke bawah dalam arah $y$.  Suatu fungsi harmonik tidak pernah dapat
mempunyai \myquote{puncak bukit} atau \myquote{lembah} pada grafiknya.  Pengamatan ini
sesuai dengan gagasan intuitif kita tentang distribusi panas keadaan tunak;
titik terpanas atau terdingin semestinya tidak berada di bagian dalam.

Umumnya persamaan Laplace merupakan bagian dari apa yang disebut
\emph{\myindex{masalah Dirichlet}}%
\footnote{Dinamai menurut matematikawan Jerman
\href{https://en.wikipedia.org/wiki/Dirichlet}{Johann Peter Gustav Lejeune Dirichlet}
(1805--1859).}.
Artinya,
kita mempunyai suatu daerah dalam bidang $xy$ dan menetapkan nilai-nilai tertentu sepanjang
batas daerah tersebut.  Kemudian kita mencoba mencari suatu solusi $u$ dari
persamaan Laplace yang didefinisikan pada
daerah ini sedemikian sehingga $u$ sesuai dengan nilai-nilai yang kita tetapkan pada
batasnya.

Dalam bagian ini, kita meninjau suatu daerah persegi panjang.  Untuk kesederhanaan,
kita menetapkan nilai batas sama dengan nol pada tiga dari keempat tepi dan hanya
menetapkan suatu fungsi sebarang pada satu tepi.
Dengan prinsip superposisi, kita dapat menggunakan solusi yang lebih sederhana ini
untuk menurunkan solusi umum
bagi nilai batas sebarang dengan menyelesaikan empat masalah yang berbeda,
satu untuk setiap tepi, dan menjumlahkan solusi-solusi tersebut.
Susunan ini diserahkan sebagai latihan.

Kita ingin menyelesaikan masalah berikut.  Misalkan $h$ dan $w$
adalah tinggi dan lebar persegi panjang kita, dengan salah satu sudut di titik asal dan
terletak dalam kuadran pertama, sehingga daerah tersebut
diberikan oleh $0 < x < w$ dan $0 < y < h$.
Perhatikan masalah

%mbx <md>
%mbx   <mrow xml:id="dirich_eq1" number="yes" eqnumber="%MBXEQNNUMBER%">
%mbx     &amp; \Delta u = 0 ,
%mbx   </mrow>
%mbx   <mrow xml:id="dirich_eq2" number="yes" eqnumber="%MBXEQNNUMBER%">
%mbx     &amp; u(0,y) = 0 \quad \text{untuk } 0 &lt; y &lt; h,
%mbx   </mrow>
%mbx   <mrow xml:id="dirich_eq3" number="yes" eqnumber="%MBXEQNNUMBER%">
%mbx     &amp; u(x,h) = 0 \quad \text{untuk } 0 &lt; x &lt; w,
%mbx   </mrow>
%mbx   <mrow xml:id="dirich_eq4" number="yes" eqnumber="%MBXEQNNUMBER%">
%mbx     &amp; u(w,y) = 0 \quad \text{untuk } 0 &lt; y &lt; h,
%mbx   </mrow>
%mbx   <mrow xml:id="dirich_eq5" number="yes" eqnumber="%MBXEQNNUMBER%">
%mbx     &amp; u(x,0) = f(x) \quad \text{untuk } 0 &lt; x &lt; w.
%mbx   </mrow>
%mbx </md>

\begin{center}
%mbxSTARTIGNORE
\begin{minipage}[b]{2.8in}
\vspace{\fill}
\begin{align}
& \Delta u = 0 , & &  \label{dirich:eq1} \\
& u(0,y) = 0 & & \text{untuk }  0 < y < h,\label{dirich:eq2} \\
& u(x,h) = 0 & & \text{untuk }  0 < x < w,\label{dirich:eq3} \\
& u(w,y) = 0 & & \text{untuk }  0 < y < h,\label{dirich:eq4} \\
& u(x,0) = f(x) & & \text{untuk }  0 < x < w.\label{dirich:eq5}
\end{align}
\vspace{\fill}
\end{minipage}
\qquad
%mbxENDIGNORE
\inputpdft{dirichsetup}{%
Diagram persegi panjang berarsir dalam bidang $xy$ dengan titik-titik sudut
$(0,0)$, $(w,0)$, $(w,h)$, dan $(0,h)$.
Bagian dalam persegi panjang diarsir dan diberi label $\Delta u$ sama dengan 0.
Bagian atas dan sisi-sisinya diberi label $u$ sama dengan 0.  Sisi bawahnya
diberi label $u$ sama dengan $f$ dari $x$.}
%mbxSTARTIGNORE
\qquad
%mbxENDIGNORE
\end{center}

Metode yang kita terapkan adalah pemisahan variabel.  Sekali lagi, kita akan
memperoleh cukup banyak solusi
blok pembangun yang memenuhi semua syarat batas homogen
(semua syarat kecuali \eqref{dirich:eq5}).  Kita memperhatikan bahwa superposisi
tetap berlaku bagi persamaan tersebut dan semua syarat homogen.
Oleh karena itu,
kita dapat menggunakan deret Fourier untuk $f(x)$ guna mencari
suatu solusi $u$ yang juga menyelesaikan \eqref{dirich:eq5}.

Kita mencoba $u(x,y) = X(x)Y(y)$.  Kita substitusikan $u$ ke persamaan untuk memperoleh
\begin{equation*}
X''Y + XY'' = 0 .
\end{equation*}
Kita tempatkan suku-suku $X$ pada satu ruas dan suku-suku $Y$ pada ruas lainnya untuk memperoleh
\begin{equation*}
- \frac{X''}{X} = \frac{Y''}{Y} .
\end{equation*}
Ruas kiri hanya bergantung pada $x$ dan ruas kanan hanya bergantung
pada $y$.  Oleh karena itu, terdapat suatu konstanta $\lambda$ sedemikian sehingga
$\lambda = \frac{-X''}{X} = \frac{Y''}{Y}$.
Kita memperoleh dua persamaan
\begin{align*}
& X'' + \lambda X = 0 , \\
& Y'' - \lambda Y = 0 .
\end{align*}
Selain itu, syarat-syarat batas homogen menyiratkan bahwa
$X(0) = X(w) = 0$ dan $Y(h) = 0$.  Dengan menggunakan persamaan untuk $X$,
kita telah melihat bahwa terdapat solusi tak trivial jika dan hanya jika
$\lambda = \lambda_n = \frac{n^2 \pi^2}{w^2}$ dan solusinya merupakan
kelipatan dari
\begin{equation*}
X_n(x) = \sin \left( \frac{n \pi}{w} x \right) .
\end{equation*}
Untuk $\lambda_n$ yang diberikan ini,
solusi umum bagi $Y$ (satu untuk setiap $n$) adalah
\begin{equation} \label{dirich:Yngensol}
Y_n(y) = A_n \cosh \left( \frac{n \pi}{w} y \right)
+ B_n \sinh \left( \frac{n \pi}{w} y \right) .
\end{equation}
Hanya terdapat satu syarat pada $Y_n$, sehingga kita dapat memilih salah satu dari $A_n$
atau $B_n$
menjadi sesuatu yang mudah digunakan.
Akan berguna jika $Y_n(0) = 1$, jadi ambil $A_n=1$.
Dengan menetapkan $Y_n(h) = 0$ dan menyelesaikan untuk $B_n$, kita memperoleh
\begin{equation*}
B_n = \frac{- \cosh \left( \frac{n \pi h }{w} \right)}%
{\sinh \left( \frac{n \pi h }{w} \right)} .
\end{equation*}
Setelah kita substitusikan $A_n$ dan $B_n$
ke \eqref{dirich:Yngensol} dan menyederhanakannya dengan menggunakan
identitas $\sinh(\alpha-\beta) =
\sinh(\alpha) \cosh(\beta) -
\cosh(\alpha) \sinh(\beta)$, kita memperoleh
\begin{equation*}
Y_n(y) =
\frac{\sinh \left( \frac{n \pi (h-y) }{w} \right)}%
{\sinh \left( \frac{n \pi h }{w} \right)} .
\end{equation*}
Kita definisikan $u_n(x,y) = X_n(x)Y_n(y)$.
Perhatikan bahwa $u_n$
memenuhi \eqref{dirich:eq1}--\eqref{dirich:eq4}.
Amati
\begin{equation*}
u_n(x,0) = X_n(x)Y_n(0) = \sin \left( \frac{n \pi}{w} x \right) .
\end{equation*}
Andaikan
\begin{equation*}
f(x) =
%\frac{a_0}{2} +
\sum_{n=1}^\infty
%a_n \cos \left( \frac{n \pi x }{w} \right)
%+
b_n \sin \left( \frac{n \pi x }{w} \right) .
\end{equation*}
Maka kita memperoleh solusi \eqref{dirich:eq1}--\eqref{dirich:eq5} dalam
bentuk berikut.
\begin{equation*}
\mybxbg{
~~
u(x,y) =
\sum_{n=1}^\infty
b_n u_n(x,y)
=
\sum_{n=1}^\infty
b_n 
\sin \left( \frac{n \pi}{w} x \right)
\left( \frac{\sinh \left( \frac{n \pi (h-y) }{w} \right)}%
{\sinh \left( \frac{n \pi h }{w} \right)} \right)
.
~~
}
\avoidbreak
\end{equation*}
Karena $u_n$ memenuhi \eqref{dirich:eq1}--\eqref{dirich:eq4}, kombinasi linear
(hingga ataupun tak hingga) dari $u_n$, dan karenanya $u$, memenuhi 
\eqref{dirich:eq1}--\eqref{dirich:eq4}. 
Kita substitusikan $y=0$ untuk melihat bahwa $u$
juga memenuhi 
\eqref{dirich:eq5}.

\begin{example}
Ambil $w=h=\pi$ dan misalkan $f(x) = \pi$.  Mari kita hitung deret
sinus untuk fungsi $\pi$ (sama dengan deret untuk gelombang persegi).
Untuk $0 < x < \pi$, kita mempunyai
\begin{equation*}
f(x) =
\sum_{\substack{n=1 \\ n \text{ ganjil}}}^\infty
\frac{4}{n}
\sin (n x) .
\end{equation*}
Solusi $u(x,y)$, lihat \figurevref{dirichsquareplot:fig},
untuk masalah Dirichlet yang bersesuaian
diberikan sebagai
\begin{equation*}
u(x,y) =
\sum_{\substack{n=1 \\ n \text{ ganjil}}}^\infty
\frac{4}{n}
\sin (n x)
\left( \frac{\sinh \bigl( n (\pi-y) \bigr) }{\sinh (n \pi)} \right)
.
\end{equation*}

\begin{myfig}
\capstart
\diffyincludegraphics{width=5in}{width=7.5in}{dirichsquareplot}{%
Grafik fungsi $u$ dari $(x,y)$ di atas bidang $xy$
pada daerah persegi dengan $x$ di antara 0 dan pi serta $y$ di antara 0
dan pi.  Pada sisi persegi tempat $y$ sama dengan 0, fungsi sama dengan pi, sedangkan pada
ketiga sisi lainnya fungsi bernilai 0.  Di tengah, tampak hampir seolah-olah
seseorang merentangkan lembaran karet dan mengikatnya pada sisi-sisi di nilai yang diberikan.}
\caption{Suhu keadaan tunak suatu pelat persegi, dengan tiga sisi
dipertahankan pada nol dan satu sisi dipertahankan pada $\pi$.\label{dirichsquareplot:fig}}
\end{myfig}
\end{example}

Skenario ini
bersesuaian dengan suhu keadaan tunak pada suatu pelat persegi berlebar $\pi$
dengan tiga sisi dipertahankan pada 0 derajat dan satu sisi dipertahankan pada $\pi$ derajat.
Jika kita mempunyai data awal sebarang pada semua sisi, maka kita menyelesaikan empat masalah,
masing-masing menggunakan satu bagian data tak homogen.  Kemudian kita menggunakan prinsip
superposisi untuk menjumlahkan keempat solusi agar memperoleh solusi bagi
masalah semula.

Cara lain
untuk memvisualisasikan solusi persamaan Laplace adalah dengan
mengambil seutas kawat dan membengkokkannya
sedemikian sehingga bentuknya bersesuaian dengan grafik
suhu di atas batas daerah Anda.  Potong selembar karet sesuai
bentuk daerah Anda---persegi dalam kasus kita---dan rentangkan,
dengan mengikat tepi-tepi lembaran pada kawat.
Lembaran karet tersebut merupakan hampiran yang baik bagi grafik solusi
persamaan Laplace dengan data batas yang diberikan.

\subsection{Latihan}

\begin{exercise}
Pada daerah yang diberikan oleh $0 < x < \pi$ dan $0 < y < \pi$,
selesaikan masalah
\begin{equation*}
\Delta u = 0, \quad u(x,0) = \sin x, \quad u(x,\pi) = 0,
\quad u(0,y) = 0, 
\quad u(\pi,y) = 0 .
\end{equation*}
\end{exercise}

\begin{exercise}
Pada daerah yang diberikan oleh $0 < x < 1$ dan $0 < y < 1$,
selesaikan masalah
\begin{equation*}
\begin{array}{ll}
u_{xx} + u_{yy} = 0, & \\
u(x,0) = \sin (\pi x) - \sin (2\pi x), \quad & u(x,1) = 0, \\
u(0,y) = 0, \quad & u(1,y) = 0 .
\end{array}
\end{equation*}
\end{exercise}

\begin{exercise}
Pada daerah yang diberikan oleh $0 < x < 1$ dan $0 < y < 1$,
selesaikan masalah
\begin{align*}
& u_{xx} + u_{yy} = 0, \\
& u(x,0) = u(x,1) = u(0,y) = u(1,y) = C .
\end{align*}
untuk suatu konstanta $C$.  Petunjuk: Tebak, lalu periksa intuisi Anda.
\end{exercise}

\begin{exercise} \label{dirich:diffsepexr}
Pada daerah yang diberikan oleh $0 < x < \pi$ dan $0 < y < \pi$,
selesaikan
\begin{equation*}
\Delta u = 0,
\quad u(x,0) = 0,
\quad u(x,\pi) = \pi,
\quad u(0,y) = y,
\quad u(\pi,y) = y .
\end{equation*}
Petunjuk: Cobalah solusi berbentuk $u(x,y) = X(x) + Y(y)$ (pemisahan
variabel yang berbeda).
\end{exercise}

\begin{exercise}
Gunakan solusi \exerciseref{dirich:diffsepexr} untuk menyelesaikan
\begin{equation*}
\Delta u = 0,
\quad u(x,0) = \sin x,
\quad u(x,\pi) = \pi,
\quad u(0,y) = y,
\quad u(\pi,y) = y .
\end{equation*}
Petunjuk: Gunakan superposisi.
\end{exercise}

\begin{exercise}
Pada daerah yang diberikan oleh $0 < x < w$ dan $0 < y < h$,
selesaikan masalah
\begin{equation*}
\begin{array}{ll}
u_{xx} + u_{yy} = 0, & \\
u(x,0) = 0, \quad & u(x,h) = f(x), \\
u(0,y) = 0, \quad & u(w,y) = 0.
\end{array}
\end{equation*}
Solusinya harus berbentuk deret yang menggunakan koefisien deret Fourier
dari $f(x)$.
\end{exercise}

\begin{exercise}
Pada daerah yang diberikan oleh $0 < x < w$ dan $0 < y < h$,
selesaikan masalah
\begin{equation*}
\begin{array}{ll}
u_{xx} + u_{yy} = 0, & \\
u(x,0) = 0, \quad & u(x,h) = 0, \\
u(0,y) = f(y), \quad & u(w,y) = 0.
\end{array}
\end{equation*}
Solusinya harus berbentuk deret yang menggunakan koefisien deret Fourier
dari $f(y)$.
\end{exercise}

\begin{exercise}
Pada daerah yang diberikan oleh $0 < x < w$ dan $0 < y < h$,
selesaikan masalah
\begin{equation*}
\begin{array}{ll}
u_{xx} + u_{yy} = 0, & \\
u(x,0) = 0, \quad & u(x,h) = 0, \\
u(0,y) = 0, \quad & u(w,y) = f(y).
\end{array}
\end{equation*}
Solusinya harus berbentuk deret yang menggunakan koefisien deret Fourier
dari $f(y)$.
\end{exercise}

\begin{exercise}
Pada daerah yang diberikan oleh $0 < x < 1$ dan $0 < y < 1$,
selesaikan masalah
\begin{equation*}
\begin{array}{ll}
u_{xx} + u_{yy} = 0, & \\
u(x,0) = \sin (9 \pi x), \quad & u(x,1) = \sin (2 \pi x), \\
u(0,y) = 0, \quad & u(1,y) = 0 .
\end{array}
\end{equation*}
Petunjuk: Gunakan superposisi.
\end{exercise}

\begin{exercise}
Pada daerah yang diberikan oleh $0 < x < 1$ dan $0 < y < 1$,
selesaikan masalah
\begin{align*}
& u_{xx} + u_{yy} = 0, \\
& u(x,0) = \sin (\pi x), \quad u(x,1) = \sin (\pi x), \\
& u(0,y) = \sin (\pi y), \quad u(1,y) = \sin (\pi y) .
\end{align*}
Petunjuk: Gunakan superposisi.
\end{exercise}

\begin{exercise}[menantang]
Dengan hanya menggunakan intuisi Anda, carilah $u(\nicefrac{1}{2},\nicefrac{1}{2})$
untuk masalah
$\Delta u = 0$, dengan $u(0,y) = u(1,y) = 100$ untuk $0 < y < 1$, dan
$u(x,0) = u(x,1) = 0$ untuk $0 < x < 1$.  Jelaskan.
\end{exercise}

\setcounter{exercise}{100}

\begin{exercise}
Pada daerah yang diberikan oleh $0 < x < 1$ dan $0 < y < 1$,
selesaikan masalah
\begin{equation*}
\Delta u = 0, \quad u(x,0) = \sum_{n=1}^\infty \frac{1}{n^2} \sin (n \pi x),
\quad u(x,1) = 0,
\quad u(0,y) = 0, 
\quad u(1,y) = 0 .
\end{equation*}
\end{exercise}
\exsol{%
$u(x,y) =
\sum\limits_{n=1}^\infty
\frac{1}{n^2}
\sin ( n \pi x )
\left( \frac{\sinh ( n \pi (1-y) )}{\sinh ( n \pi )} \right)
$
}

\begin{exercise}
Pada daerah yang diberikan oleh $0 < x < 1$ dan $0 < y < 2$,
selesaikan masalah
\begin{equation*}
\Delta u = 0, \quad u(x,0) = 0.1 \sin (\pi x),
\quad u(x,2) = 0,
\quad u(0,y) = 0, 
\quad u(1,y) = 0 .
\end{equation*}
\end{exercise}
\exsol{%
$u(x,y) =
0.1
\sin ( \pi x )
\left( \frac{\sinh ( \pi (2-y)  )}%
{\sinh ( 2 \pi )} \right)$
}

%\begin{equation*}
%f(x) =
%%\frac{a_0}{2} +
%\sum_{n=1}^\infty
%%a_n \cos \left( \frac{n \pi x }{w} \right)
%%+
%b_n \sin \left( \frac{n \pi x }{w} \right) .
%\end{equation*}
%
%u(x,y) =
%\sum_{n=1}^\infty
%b_n u_n(x,y)
%=
%\sum_{n=1}^\infty
%b_n 
%\sin \left( \frac{n \pi}{w} x \right)
%\left( \frac{\sinh \left( \frac{n \pi (h-y) }{w} \right)}%
%{\sinh \left( \frac{n \pi h }{w} \right)} \right)

%%%%%%%%%%%%%%%%%%%%%%%%%%%%%%%%%%%%%%%%%%%%%%%%%%%%%%%%%%%%%%%%%%%%%%%%%%%%%%

\sectionnewpage
\section{Masalah Dirichlet pada lingkaran dan kernel Poisson}
\label{dirichdisc:section}

%mbxINTROSUBSECTION

\sectionnotes{2 kuliah\EPref{, \S9.7 dalam \cite{EP}}\BDref{,
\S10.8 dalam \cite{BD}}}

\subsection{Laplace dalam koordinat polar}

\begin{mywrapfigsimp}[5]{1.3in}{1.5in}
\diffypdfversion{\vspace*{5pt}}
\noindent
\inputpdft{polarcoords}{%
Diagram pada bidang dengan sumbu horizontal (sumbu $x$) dan vertikal (sumbu $y$).
Suatu titik dalam kuadran pertama ditandai sebagai $(r,\theta)$, dan sebuah garis
digambar dari titik asal ke titik ini.  Jarak dari titik asal ke
titik tersebut ditandai sebagai $r$, dan sudut yang dibentuknya dengan sumbu $x$ positif
diberi label $\theta$.}
\diffypdfversion{\vspace*{5pt}}
\end{mywrapfigsimp}
Lingkaran, bukan persegi panjang, merupakan latar yang lebih alami untuk persamaan Laplace $\Delta u = 0$.
Di sisi lain, yang membuat
masalahnya agak lebih sulit adalah bahwa kita memerlukan koordinat polar.
Ingatlah bahwa koordinat polar untuk bidang $(x,y)$ adalah $(r,\theta)$: 
\begin{equation*}
x = r \cos \theta , \quad y = r \sin \theta ,
\end{equation*}
dengan $r \geq 0$ dan $-\pi < \theta \leq \pi$.  Jadi, titik $(x,y)$ berjarak
$r$ dari titik asal pada sudut $\theta$ dari
sumbu $x$ positif.

\begin{mywrapfigsimp}{2.4in}{2.7in}
\noindent
\inputpdft{dirichdiscsetup}{%
Diagram lingkaran berjari-jari 1 yang berpusat pada titik asal dalam bidang $xy$.
Bagian dalam lingkaran diarsir dan ditandai dengan $\Delta u$ sama dengan 0.
Lingkaran batas berupa garis tebal gelap yang diberi label
$u$ dari $(1,\theta)$ sama dengan $g$ dari $\theta$.}
\end{mywrapfigsimp}
Setelah kita mengetahui koordinat kita, mari kita nyatakan masalah yang ingin
kita selesaikan.  Kita mempunyai daerah lingkaran berjari-jari 1 dan tertarik
pada masalah Dirichlet untuk persamaan Laplace pada daerah ini.  Misalkan
$u(r,\theta)$ menyatakan suhu di titik $(r,\theta)$ dalam koordinat
polar.
Kita akan menetapkan suhu pada lingkaran sebagai $g(\theta)$.

Kita mempunyai masalah:
\begin{equation} \label{dirichdisc:theprobeq}
\begin{aligned}
& \Delta u = 0  & & \text{untuk } \; r < 1, \\
& u(1,\theta) = g(\theta) & & \text{untuk } \; {-\pi} < \theta \leq \pi.
\end{aligned}
\end{equation}

Persoalan pertama yang kita hadapi adalah bahwa kita tidak mengetahui Laplasian
dalam koordinat polar.
Biasanya, kita akan mencari $u_{xx}$ dan $u_{yy}$ dalam suku-suku
turunan terhadap $r$ dan $\theta$.  Kita perlu menyelesaikan
$r$ dan $\theta$ dalam suku-suku $x$ dan $y$.  Dalam kasus ini,
lebih mudah bekerja dengan
arah sebaliknya.  Kita menghitung turunan terhadap $r$ dan $\theta$ dalam suku-suku
turunan terhadap $x$ dan $y$, lalu kita menyelesaikannya.  Dengan cara ini,
penghitungannya lebih mudah.  Pertama,
\begin{equation*}
\begin{aligned}
& x_r = \cos \theta, & &
x_\theta = - r \sin \theta, \\
& y_r = \sin \theta, & &
y_\theta = r \cos \theta.
\end{aligned}
\end{equation*}
Menurut aturan rantai, kita memperoleh
\begin{align*}
u_r & = u_x x_r + u_y y_r = \cos(\theta) u_x + \sin(\theta) u_y ,
\\
u_{rr} & =
\cos(\theta) ( u_{xx} x_r +u_{xy} y_r )
+ \sin(\theta) ( u_{yx} x_r +u_{yy} y_r )
\\
&
=
\cos^2(\theta) u_{xx} +
2 \cos(\theta)\sin(\theta) u_{xy} +
\sin^2(\theta) u_{yy} .
\end{align*}
Demikian pula untuk turunan terhadap $\theta$.  Perhatikan bahwa kita harus menggunakan
aturan hasil kali untuk turunan kedua.
\begin{align*}
u_\theta & = u_x x_\theta + u_y y_\theta =
-r\sin(\theta) u_x + r\cos(\theta) u_y ,
\\
u_{\theta\theta} & =
-r\cos(\theta) u_x
-r\sin(\theta) (u_{xx} x_\theta + u_{xy} y_\theta)
-r\sin(\theta) u_y
+
r\cos(\theta) (u_{yx} x_\theta + u_{yy} y_\theta)
%\\
%& = 
%-r\cos(\theta) u_x
%-r\sin(\theta) (u_{xx} (-r\sin(\theta)) + u_{xy} (r \cos(\theta)))
%-r\sin(\theta) u_y
%+
%r\cos(\theta) (u_{yx} (-r\sin(\theta)) + u_{yy} (r \cos (\theta)))
\\
& = 
-r\cos(\theta) u_x
-r\sin(\theta) u_y
+r^2 \sin^2(\theta) u_{xx}
-r^2 2\sin(\theta)\cos(\theta) u_{xy}
+r^2 \cos^2(\theta) u_{yy} .
\end{align*}
Sekarang mari kita mencoba mencari $u_{xx} + u_{yy}$.  Kita mulai dengan
$\frac{1}{r^2} u_{\theta\theta}$ untuk menghilangkan faktor-faktor $r^2$ yang mengganggu.
Jika kita menambahkan $u_{rr}$
dan menggunakan fakta bahwa $\cos^2(\theta) +\sin^2(\theta) = 1$, kita memperoleh
\begin{equation*}
\frac{1}{r^2} u_{\theta\theta}
+
u_{rr}
=
u_{xx} + u_{yy} - \frac{1}{r} \cos(\theta) u_x - \frac{1}{r} \sin(\theta)
u_y .
\end{equation*}
Kita belum sepenuhnya sampai, tetapi yang kurang hanyalah 
$\frac{1}{r} u_r$.  Kita menambahkannya untuk memperoleh
\emph{\myindex{Laplasian dalam koordinat polar}}:
\begin{equation*}
\mybxbg{~~
\Delta u 
=
u_{xx} + u_{yy} =
\frac{1}{r^2} u_{\theta\theta}
+
\frac{1}{r} u_{r}
+
u_{rr} .
~~}
\end{equation*}

Perhatikan bahwa Laplasian dalam koordinat polar tidak lagi mempunyai koefisien
konstan.
%Bentuknya tampak lebih rumit daripada bentuk semula, tetapi sebenarnya
%penghitungannya menjadi lebih mudah mulai sekarang.

\subsection{Solusi deret}

Mari kita pisahkan variabel seperti biasa.  Artinya, kita mencoba
$u(r,\theta) = R(r)\Theta(\theta)$.  Maka
\begin{equation*}
0 = \Delta u = 
\frac{1}{r^2} R \Theta''
+
\frac{1}{r} R' \Theta
+
R'' \Theta .
\end{equation*}
Kita tempatkan $R$ pada satu ruas dan $\Theta$ pada ruas lainnya, lalu menyimpulkan
bahwa kedua ruas harus konstan.
\begin{align*}
\frac{1}{r^2} R \Theta''
& =
-
\left(\frac{1}{r} R' + R''\right) \Theta 
\\
\frac{\Theta''}{\Theta}
& =
-
\frac{r R' + r^2 R''}{R} = -\lambda
\end{align*}
Kita memperoleh dua persamaan:
\begin{align*}
& \Theta'' + \lambda \Theta = 0 ,
\\
& r^2 R'' + r R' -\lambda R = 0.
\end{align*}
Pertama-tama, kita berfokus pada $\Theta$.  Kita mengetahui bahwa $u(r,\theta)$ seharusnya
periodik-$2\pi$ dalam $\theta$, yaitu,
$u(r,\theta) = u(r,\theta+2\pi)$.  Oleh karena itu, solusi
$\Theta'' + \lambda \Theta = 0$ harus periodik-$2\pi$.
Kita telah melihat masalah semacam ini dalam \exampleref{bvp-periodic:example}.
Kita menyimpulkan
bahwa
%$\lambda = 0,1,4,9,\ldots$.  Yaitu,
$\lambda = n^2$ untuk suatu
bilangan bulat tak negatif $n=0,1,2,3,\ldots$.  Persamaannya menjadi
$\Theta'' + n^2 \Theta = 0$.  Ketika $n=0$, persamaan tersebut hanyalah
$\Theta'' = 0$, sehingga kita mempunyai solusi umum $A \theta + B$.  Karena
$\Theta$ periodik,
$A=0$.
Untuk memudahkan, kita tuliskan solusi ini sebagai
\begin{equation*}
\Theta_0 = \frac{a_0}{2}
\end{equation*}
untuk suatu konstanta $a_0$.  Untuk $n$ positif,
solusi
$\Theta'' + n^2 \Theta = 0$ adalah
\begin{equation*}
\Theta_n = a_n \cos(n\theta) + b_n \sin(n\theta) ,
\end{equation*}
untuk suatu konstanta $a_n$ dan $b_n$.

Berikutnya, kita meninjau persamaan untuk $R$,
\begin{equation*}
r^2 R'' + r R' - n^2 R = 0.
\end{equation*}
Persamaan ini telah muncul dalam latihan sebelumnya---kita
menyelesaikannya dalam \exerciseref{sol:eulerex}
dan \exercisevref{sol:eulerexln}.  Gagasannya adalah mencoba suatu solusi
$r^s$, dan jika itu tidak memberikan dua solusi, juga mencoba solusi berbentuk
$r^s \ln r$.  Seperti biasa, kita namai solusinya $R_n$.  Ketika $n=0$, kita memperoleh
\begin{equation*}
R_0 = A r^0 + B r^0 \ln r = A + B \ln r ,
\end{equation*}
dan jika $n > 0$, kita memperoleh
\begin{equation*}
R_n = A r^n + B r^{-n} .
\end{equation*}
Fungsi $u(r,\theta)$ harus berhingga (tidak boleh menjadi tak terbatas) di titik asal, yaitu, ketika $r=0$.
Jadi, $B=0$ dalam kedua
kasus, sebab jika tidak, $r^{-n}$ atau $\ln r$ menjadi tak terbatas ketika $r\to 0$.
Tetapkan $A=1$ dalam kedua kasus; konstanta dalam $\Theta_n$
akan menyerap faktor tersebut sehingga tidak ada yang hilang.  Yaitu,
\begin{equation*}
R_0 = 1
\qquad \text{dan} \qquad
R_n = r^n .
\end{equation*}
Solusi-solusi blok pembangun kita adalah
\begin{equation*}
u_0(r,\theta) = \frac{a_0}{2}
\qquad \text{dan} \qquad
u_n(r,\theta) = a_n r^n \cos(n \theta) + b_n r^n \sin(n \theta) .
\end{equation*}
Dengan menggabungkan semuanya, solusi kita adalah
\begin{equation*}
\mybxbg{~~
u(r,\theta)
=
\frac{a_0}{2} +
\sum_{n=1}^\infty
\Bigl[
a_n r^n \cos(n \theta) + b_n r^n \sin(n \theta)
\Bigr]
.
~~}
\end{equation*}

Kita melihat syarat batas dalam \eqref{dirichdisc:theprobeq},
\begin{equation*}
g(\theta) = u(1,\theta)
=
\frac{a_0}{2} +
\sum_{n=1}^\infty
\Bigl[
a_n \cos(n \theta) + b_n \sin(n \theta)
\Bigr].
\end{equation*}
Oleh karena itu, untuk menyelesaikan \eqref{dirichdisc:theprobeq},
kita uraikan $g(\theta)$, yang merupakan
fungsi periodik-$2\pi$, sebagai deret Fourier, kemudian
kalikan suku ke-$n$ dengan $r^n$.  Untuk
mencari $a_n$ dan $b_n$, kita hitung
\begin{equation*}
a_n =
\frac{1}{\pi} \int_{-\pi}^\pi g(\theta) \cos (n\theta) \, d\theta \qquad
\text{dan} \qquad
b_n =
\frac{1}{\pi} \int_{-\pi}^\pi g(\theta) \sin (n\theta) \, d\theta.
\end{equation*}

\begin{example}
Andaikan kita ingin menyelesaikan
\begin{align*}
& \Delta u = 0 , \qquad 0 \leq r < 1, \quad -\pi < \theta \leq \pi,\\
& u(1,\theta) = \cos(10\,\theta), \qquad -\pi < \theta \leq \pi.
\end{align*}

$g(\theta)$ sudah terurai, sehingga
solusinya adalah
\begin{equation*}
u(r,\theta) = r^{10} \cos(10\,\theta) .
\end{equation*}
Lihat grafik dalam \figurevref{dirichdisc:tenspeedfig}.
Hal yang perlu diperhatikan dalam contoh ini adalah bahwa pengaruh frekuensi tinggi
sebagian besar terasa pada batas.  Di tengah cakram, solusinya
sangat dekat dengan nol.  Hal itu karena $r^{10}$ cukup kecil ketika $r$
dekat dengan 0.
\begin{myfig}
\capstart
\diffyincludegraphics{width=5in}{width=7.5in}{dirichdisc-tenspeed}{%
Grafik fungsi $u$ dari $(x,y)$ di atas bidang $xy$
pada bagian dalam lingkaran berjari-jari 1 yang berpusat di titik asal.
Pada batas, permukaan berosilasi naik dan turun 10 kali
mengelilingi lingkaran.  Kemudian, di bagian dalam, permukaan dengan cepat mendekati 0
ketika $(x,y)$ semakin mendekati titik asal.}
\caption{Solusi masalah Dirichlet dalam cakram dengan
$\cos(10\,\theta)$ sebagai data batas.\label{dirichdisc:tenspeedfig}}
\end{myfig}
\end{example}

\begin{example}
Mari kita selesaikan masalah yang lebih sulit.
Perhatikan suatu batang panjang
dengan penampang melintang lingkaran berjari-jari 1.  Andaikan kita ingin menyelesaikan
masalah panas keadaan tunak dalam batang tersebut.
Jika batang cukup panjang, kita hanya perlu menyelesaikan
persamaan Laplace dalam dua dimensi.  Kita tempatkan pusat batang di
titik asal, dan kita memperoleh tepat daerah yang sedang
kita kaji---lingkaran berjari-jari 1.  Untuk syarat batas, andaikan dalam
koordinat Kartesius $x$ dan
$y$, suhu pada batas adalah 0 ketika $y < 0$, dan $2y$ ketika $y > 0$.

Mari kita susun masalahnya.
Karena $y = r\sin(\theta)$, maka pada
lingkaran berjari-jari 1, yaitu, tempat $r=1$, kita mempunyai $2y = 2\sin(\theta)$.  Jadi,
%Masalahnya menjadi
\begin{align*}
& \Delta u = 0 , \qquad 0 \leq r < 1, \quad -\pi < \theta \leq \pi,\\
& u(1,\theta) = 
\begin{cases}
2\sin(\theta) & \text{jika } \; \phantom{-}0 \leq \theta \leq \pi, \\
0 & \text{jika } \; {-\pi} < \theta < 0.
\end{cases}
\end{align*}

Sekarang kita harus menghitung deret Fourier untuk syarat
batas tersebut.  Pada tahap ini pembaca telah mempunyai banyak pengalaman dalam menghitung
deret Fourier, sehingga kita cukup menyatakan bahwa 
\begin{equation*}
u(1,\theta) = 
\frac{2}{\pi}
+
\sin(\theta)
+
\sum_{n=1}^\infty \frac{-4}{\pi(4n^2-1)} \cos(2n\theta) .
\end{equation*}

\begin{exercise}
Hitung deret untuk $u(1,\theta)$ dan periksa bahwa deret tersebut benar-benar seperti
yang baru saja kita klaim.  Petunjuk: Berhati-hatilah; pastikan tidak membagi dengan nol.
\end{exercise}

Untuk memperoleh solusi (lihat \figurevref{dirichdisc:zero2yfig}), kita
menuliskan deretnya dengan mengalikan suku-suku dalam deret $g(\theta)$
dengan $r^n$ di tempat yang tepat:
\begin{equation*}
u(r,\theta) = 
\frac{2}{\pi}
+
r\sin(\theta)
+
\sum_{n=1}^\infty \frac{-4r^{2n}}{\pi(4n^2-1)} \cos(2n\theta) .
\end{equation*}
\begin{myfig}
\capstart
\diffyincludegraphics{width=5in}{width=7.5in}{dirichdisc-zero2y}{%
Grafik fungsi $u$ dari $(x,y)$ di atas bidang $xy$
pada bagian dalam lingkaran berjari-jari 1 yang berpusat di titik asal.
Pada batas, permukaan berada pada nol ketika $y$ negatif
atau naik secara linear terhadap $y$ untuk $y$ positif.  Di dalam lingkaran,
permukaan tampak seolah-olah selembar karet telah ditempatkan pada
nilai-nilai di batas.}
\caption{Solusi masalah Dirichlet dengan
data batas 0 untuk $y < 0$ dan $2y$ untuk $y > 0$.\label{dirichdisc:zero2yfig}}
\end{myfig}
%Grafik solusi diberikan dalam \figurevref{dirichdisc:zero2yfig}.
\end{example}

%Perhatikan bahwa rumus-rumus tersebut tidak sulit digeneralisasi
%untuk lingkaran dengan
%sebarang jari-jari, dan generalisasi ini diserahkan kepada pembaca dalam latihan.

\subsection{Kernel Poisson}

Terdapat cara lain untuk menyelesaikan masalah Dirichlet---dengan bantuan
kernel integral.  Artinya, kita akan mencari suatu fungsi $P(r,\theta,\alpha)$
yang disebut \emph{\myindex{kernel Poisson}}\footnote{%
Dinamai menurut matematikawan Prancis
\href{https://en.wikipedia.org/wiki/Sim\%C3\%A9on_Denis_Poisson}{Sim\'eon
Denis Poisson}
(1781--1840).} sedemikian sehingga
\begin{equation*}
u(r,\theta) = 
\frac{1}{2\pi}
\int_{-\pi}^{\pi}
P(r,\theta,\alpha) \, g(\alpha) \,d\alpha .
\end{equation*}
Walaupun integral tersebut secara umum tidak dapat diselesaikan secara analitik, integral itu dapat
dievaluasi secara numerik.  Bahkan, kecuali jika data batas sudah diberikan
sebagai deret Fourier, barangkali jauh lebih mudah mengevaluasi
rumus ini secara numerik karena hanya ada satu integral yang perlu dievaluasi.

Rumus tersebut juga mempunyai penerapan teoretis.
Sebagai contoh, karena $P(r,\theta,\alpha)$ 
akan mempunyai tak hingga banyak turunan, maka
dengan diferensiasi di bawah integral, kita mendapati
bahwa solusi $u(r,\theta)$ mempunyai tak hingga banyak turunan, setidaknya
ketika berada di dalam lingkaran, $r < 1$.  Dengan \myquote{mempunyai tak hingga banyak
turunan,} yang perlu Anda
bayangkan adalah bahwa $u(r,\theta)$ \myquote{tidak mempunyai sudut} dan semua
turunan parsialnya dari semua orde ada dan juga \myquote{tidak mempunyai sudut.}

%Rumus integral dan kernel integral yang serupa
%selalu ada, apa pun bentuk daerahnya,
%tetapi jauh lebih sulit dihitung ketika daerah tersebut
%sangat tak beraturan.  Untuk lingkaran, rumus yang pada akhirnya akan kita peroleh
%cukup sederhana dan mempunyai interpretasi geometris yang bagus.

Kita akan menghitung
rumus untuk $P(r,\theta,\alpha)$ dari solusi
deret, dan gagasan ini dapat diterapkan kapan pun Anda mempunyai
solusi deret yang mudah digunakan, dengan koefisien yang diperoleh melalui integrasi.
Oleh karena itu, Anda dapat menerapkan penalaran ini untuk memperoleh kernel integral semacam itu
bagi persamaan lain, seperti persamaan panas.
Penghitungannya panjang dan \myindex{menjemukan}, tetapi tidak terlalu sulit.
Karena gagasan-gagasannya sering diterapkan dalam konteks serupa, penting untuk
memahami cara kerja penghitungan ini.

Yang kita lakukan adalah memulai dengan solusi deret dan mengganti koefisien-koefisiennya
dengan integral yang menghitung koefisien tersebut.  Kemudian kita mencoba menuliskan semuanya sebagai
satu integral.  Kita harus menggunakan variabel semu yang berbeda untuk
integrasi, sehingga kita menggunakan $\alpha$, bukan $\theta$.
\begin{equation*}
\begin{split}
u(r,\theta)
& =
\frac{a_0}{2} +
\sum_{n=1}^\infty
\Bigl[
a_n r^n \cos(n \theta) + b_n r^n \sin(n \theta)
\Bigr]
\\
& =
\underbrace{
 \left(
  \frac{1}{2\pi} \int_{-\pi}^\pi g(\alpha) \, d\alpha
 \right)
}_{\frac{a_0}{2}}
+
%\\
%& ~~~~~~~~ +
\sum_{n=1}^\infty
\Biggl[
\underbrace{
 \left(
  \frac{1}{\pi} \int_{-\pi}^\pi g(\alpha) \cos (n\alpha) \, d\alpha
 \right)
}_{a_n}
r^n \cos(n \theta) +
\\
& ~~~~~~~~ +
\underbrace{
 \left(
  \frac{1}{\pi} \int_{-\pi}^\pi g(\alpha) \sin (n\alpha) \, d\alpha
 \right)
}_{b_n}
r^n \sin(n \theta)
\Biggr]
\\
& =
\frac{1}{2\pi}
\int_{-\pi}^\pi
\left(  g(\alpha)
+
2
\sum_{n=1}^\infty
g(\alpha) \cos (n\alpha) 
\, r^n \cos(n \theta) +
g(\alpha) \sin (n\alpha)
\, r^n \sin(n \theta)
\right) \,d\alpha
\\
& =
\frac{1}{2\pi}
\int_{-\pi}^\pi
\underbrace{
 \left( 1
 +
 2
 \sum_{n=1}^\infty
 r^n 
 \bigl(
 \cos (n\alpha) 
 \cos(n \theta) +
 \sin (n\alpha)
 \sin(n \theta) \bigr)
 \right)
}_{P(r,\theta,\alpha)}
g(\alpha) \,d\alpha .
%\\
%& =
%\frac{1}{2\pi}
%\int_{-\pi}^\pi
%\left( 1
%+
%2
%\sum_{n=1}^\infty
%r^n 
%\cos \bigl(n(\theta-\alpha)\bigr)
%\right) g(\alpha) \,d\alpha .
\end{split}
\end{equation*}
Baiklah, kita telah memperoleh yang diinginkan: Ekspresi dalam tanda kurung adalah
kernel Poisson, $P(r,\theta,\alpha)$.  Namun, kita masih dapat memperbaikinya.  Ekspresi tersebut masih diberikan sebagai
deret, sedangkan kita benar-benar menginginkan ekspresi sederhana yang bagus untuknya. 
Kita harus bekerja sedikit lebih keras.  Triknya adalah menulis ulang semuanya dalam suku-suku
eksponensial kompleks.  Mari kita kerjakan
kernel itu saja.
\begin{equation*}
\begin{split}
P(r,\theta,\alpha)
& =
1
+
2
\sum_{n=1}^\infty
r^n 
\bigl(
\cos (n\alpha) 
\cos(n \theta) +
\sin (n\alpha)
\sin(n \theta) \bigr)
\\
& =
1
+
2
\sum_{n=1}^\infty
r^n 
\cos \bigl(n(\theta-\alpha)\bigr)
\\
& =
1
+
\sum_{n=1}^\infty
r^n 
\bigl(
e^{in(\theta-\alpha)} +
e^{-in(\theta-\alpha)} \bigr)
\\
& =
1
+
\sum_{n=1}^\infty
{\bigl(
re^{i(\theta-\alpha)}\bigr)}^{n}
+
\sum_{n=1}^\infty
{\bigl(
re^{-i(\theta-\alpha)}\bigr)}^{n} .
\end{split}
\end{equation*}
Dalam ekspresi di atas, kita mengenali
\emph{\myindex{deret geometri}}.
Ingat dari kalkulus bahwa jika $z$ adalah bilangan kompleks dengan $\lvert z \rvert < 1$, maka
\begin{equation*}
\sum_{n=1}^\infty z^n = \frac{z}{1-z} .
\end{equation*}
Perhatikan bahwa $n$ bermula dari $1$, dan itulah sebabnya kita mempunyai $z$ pada pembilang.
Deret tersebut adalah deret geometri standar yang dikalikan dengan $z$.
Kita dapat menggunakan $z = re^{i(\theta-\alpha)}$, sebab
ternyata $\lvert re^{i(\theta-\alpha)} \rvert = r < 1$.
Kita lanjutkan penghitungannya.
\begin{equation*}
\begin{split}
P(r,\theta,\alpha)
& =
1
+
\sum_{n=1}^\infty
{\bigl(
re^{i(\theta-\alpha)}\bigr)}^{n}
+
\sum_{n=1}^\infty
{\bigl(
re^{-i(\theta-\alpha)}\bigr)}^{n}
\\
& =
1
+
\frac{re^{i(\theta-\alpha)}}{1-re^{i(\theta-\alpha)}}
+
\frac{re^{-i(\theta-\alpha)}}{1-re^{-i(\theta-\alpha)}}
\\
& = 
\frac{
\bigl(1-re^{i(\theta-\alpha)}\bigr)\bigl(1-re^{-i(\theta-\alpha)}\bigr)
+
\bigl(1-re^{-i(\theta-\alpha)}\bigr)re^{i(\theta-\alpha)} +
\bigl(1-re^{i(\theta-\alpha)}\bigr)re^{-i(\theta-\alpha)}}
{\bigl(1-re^{i(\theta-\alpha)}\bigr)\bigl(1-re^{-i(\theta-\alpha)}\bigr)}
\\
& = 
\frac{1 -r^2}{1 - re^{i(\theta-\alpha)} - re^{-i(\theta-\alpha)} +r^2}
\\
& = 
\frac{1 -r^2}{1 - 2r\cos(\theta-\alpha) +r^2} .
\end{split}
\end{equation*}
Itulah rumus yang dapat kita gunakan.  Solusi
masalah Dirichlet dengan menggunakan kernel Poisson adalah
\begin{equation*}
\mybxbg{~~
u(r,\theta) = 
\frac{1}{2\pi} \int_{-\pi}^{\pi}
\frac{1 -r^2}{1 - 2r\cos(\theta-\alpha) +r^2} g(\alpha) \, d\alpha .
~~}
\end{equation*}
Kadang-kadang rumus kernel Poisson
diberikan bersama konstanta $\frac{1}{2\pi}$; dalam hal itu, tentu saja kita
tidak boleh menyisakannya di depan integral.
Kadang-kadang batas
integralnya diberikan dari 0 hingga $2\pi$; semua yang berada di dalamnya
periodik-$2\pi$ dalam $\alpha$, sehingga ini tidak mengubah integral.

%11 adalah jumlah baris, harus disesuaikan
\begin{mywrapfigsimp}[11]{2.1in}{2.4in}
\diffypdfversion{\vspace*{5pt}}
\noindent
\inputpdft{poisson}{%
Diagram sebuah lingkaran.  Terdapat suatu titik bertanda $(r,\theta)$
di dalam lingkaran, dan suatu titik $(1,\alpha)$ bertanda pada lingkaran.
Jarak dari pusat lingkaran ke $(1,\alpha)$ panjangnya 1,
jarak dari pusat ke $(r,\theta)$ adalah $r$, dan
jarak antara $(1,\alpha)$ dan $(r,\theta)$ adalah $s$.}
\end{mywrapfigsimp}
Janganlah kita meninggalkan kernel Poisson tanpa menjelaskan makna
geometrisnya.
Kernel yang bersesuaian dengan titik $(r,\theta)$ dalam cakram
dan titik $(1,\alpha)$ pada batas hanya bergantung
pada jarak $(r,\theta)$ dari pusat dan
jarak antara $(r,\theta)$ dan $(1,\alpha)$.
Misalkan $s$ adalah jarak dari $(r,\theta)$ ke
$(1,\alpha)$.
Jarak $s$ dalam koordinat polar ini diberikan tepat oleh akar kuadrat
dari $1 - 2r\cos(\theta-\alpha) +r^2$.  Artinya, kernel Poisson sebenarnya
adalah rumus
\begin{equation*}
\frac{1-r^2}{s^2} .
\end{equation*}
%Lihat gambar di sebelah kanan.

Satu catatan terakhir mengenai rumus tersebut adalah bahwa rumus itu sebenarnya
merupakan rata-rata berbobot dari nilai-nilai batas.
Pertama, kita melihat
apa yang terjadi di titik asal,
yakni ketika $r=0$: %(dan kemudian $\theta$ dapat berupa apa pun)
\begin{equation*}
\begin{split}
u(0,0) &= 
\frac{1}{2\pi} \int_{-\pi}^{\pi}
\frac{1 -0^2}{1 - 2(0)\cos(0-\alpha) +0^2} g(\alpha) \, d\alpha
\\
& =
\frac{1}{2\pi} \int_{-\pi}^{\pi}
g(\alpha) \, d\alpha .
\end{split}
\end{equation*}
Jadi, $u(0,0)$ tepat merupakan nilai rata-rata $g(\theta)$ dan
karena itu nilai rata-rata $u$ pada batas.  Ini
merupakan sifat umum fungsi harmonik: nilai pada suatu titik $p$
sama dengan rata-rata nilai pada lingkaran yang berpusat di $p$.

Yang dinyatakan rumus tersebut pada titik-titik lain di dalam lingkaran
adalah bahwa nilai solusi
merupakan rata-rata berbobot dari data batas $g(\theta)$.  Kernel tersebut
lebih besar ketika $(1,\alpha)$ lebih dekat ke $(r,\theta)$.  Oleh karena itu, ketika
menghitung $u(r,\theta)$, kita
memberikan bobot lebih besar kepada nilai-nilai $g(\alpha)$ ketika $(1,\alpha)$ lebih dekat ke $(r,\theta),$
dan bobot lebih kecil kepada nilai-nilai $g(\alpha)$ ketika $(1,\alpha)$ jauh dari $(r,\theta)$.

\subsection{Latihan}

\begin{exercise}
Dengan menggunakan deret, selesaikan
$\Delta u = 0$, $u(1,\theta) = \lvert \theta \rvert$, untuk $-\pi < \theta
\leq \pi$.
\end{exercise}

\begin{exercise}
Dengan menggunakan deret, selesaikan $\Delta u = 0$, $u(1,\theta) = g(\theta)$ untuk
data berikut.  Petunjuk: identitas trigonometri.
\begin{tasks}(2)
\task
$g(\theta) = 
\nicefrac{1}{2} + 3\sin(\theta) + \cos(3\theta)$
\task
$g(\theta) = 
3\cos(3\theta) + 3\sin(3\theta) + \sin(9\theta)$
\task
$g(\theta) = 2 \cos(\theta+1)$
\task
$g(\theta) = \sin^2(\theta)$
\end{tasks}
\end{exercise}

\begin{exercise}
Dengan menggunakan kernel Poisson, berikan solusi untuk
$\Delta u = 0$, dengan $u(1,\theta)$ bernilai nol untuk $\theta$ di luar
interval $[-\nicefrac{\pi}{4},\nicefrac{\pi}{4}]$ dan 
$u(1,\theta)$ bernilai 1 untuk $\theta$ pada interval
$[-\nicefrac{\pi}{4},\nicefrac{\pi}{4}]$.
\end{exercise}

\begin{exercise}
\pagebreak[2]
\leavevmode
\begin{tasks}
\task Gambarlah grafik kernel Poisson sebagai fungsi dari $\alpha$
ketika $r=\nicefrac{1}{2}$ dan $\theta = 0$.
\task Jelaskan apa yang terjadi pada grafik ketika Anda memperbesar $r$ (ketika $r$
mendekati 1).
\task Dengan mengetahui bahwa solusi $u(r,\theta)$ merupakan rata-rata berbobot
dari $g(\theta)$ dengan kernel Poisson sebagai bobotnya, jelaskan arti jawaban Anda
untuk bagian b).
\end{tasks}
\end{exercise}

\begin{exercise} \label{exercise:dirichproblemxy}
Misalkan $g(\theta)$ adalah fungsi $xy = \cos \theta \sin
\theta$ pada batas.  Gunakan solusi deret untuk mencari suatu solusi
masalah Dirichlet $\Delta u = 0$, $u(1,\theta) = g(\theta)$.  Sekarang
ubahlah solusi ke koordinat Kartesius $x$ dan $y$.  Apakah
solusi ini mengejutkan?  Petunjuk: Gunakan identitas trigonometri Anda.
\end{exercise}

\begin{exercise}
Laksanakan penghitungan yang kita perlukan dalam pemisahan variabel dan selesaikan
$r^2 R'' + r R' - n^2 R = 0$, untuk $n=0,1,2,3,\ldots$.
\end{exercise}

\begin{exercise}[menantang]
Turunkan solusi deret bagi masalah Dirichlet jika daerahnya merupakan
lingkaran berjari-jari $\rho$, bukan
1.
Artinya, selesaikan $\Delta u = 0$, $u(\rho,\theta) = g(\theta)$.
\end{exercise}

\begin{exercise}[menantang]
\leavevmode
\begin{tasks}
\task
Carilah solusi untuk
$\Delta u = 0$, $u(1,\theta) = x^2y^3 + 5 x^2$.  Tuliskan jawabannya dalam koordinat Kartesius.
\task
Sekarang selesaikan
$\Delta u = 0$, $u(1,\theta) = x^k y^\ell$.
Tuliskan solusinya dalam koordinat Kartesius.
\task
Diberikan suatu polinom $P(x,y) = \sum_{j=0}^m \sum_{k=0}^n c_{j,k}
x^j y^k$, selesaikan $\Delta u = 0$, $u(1,\theta) = P(x,y)$ (yaitu, tuliskan
rumus jawabannya).  Tuliskan jawaban
dalam koordinat Kartesius.
\end{tasks}
Perhatikan bahwa jawabannya kembali berupa polinom dalam $x$ dan $y$.
Lihat juga \exerciseref{exercise:dirichproblemxy}.
\end{exercise}

\setcounter{exercise}{100}

\begin{exercise}
Dengan menggunakan deret, selesaikan
$\Delta u = 0$, $u(1,\theta) = 1+ \sum\limits_{n=1}^\infty \frac{1}{n^2}\sin(n\theta)$.
\end{exercise}
\exsol{%
$u = 1+ \sum\limits_{n=1}^\infty \frac{1}{n^2}r^n\sin(n\theta)$
}

\begin{exercise}
Dengan menggunakan solusi deret, carilah solusi untuk
$\Delta u = 0$, $u(1,\theta) = 1- \cos(\theta)$.  Nyatakan solusinya
dalam koordinat Kartesius (yaitu, dengan menggunakan $x$ dan $y$).
\end{exercise}
\exsol{%
$u = 1-x$
}

\begin{exercise}
\leavevmode
\begin{tasks}
\task
Cobalah menebak suatu solusi untuk $\Delta u = -1$, $u(1,\theta) = 0$.
Petunjuk: Cobalah solusi yang hanya bergantung pada~$r$.  Selain itu, mula-mula jangan khawatir
mengenai syarat batas.
\task
Sekarang selesaikan $\Delta u = -1$, $u(1,\theta) = \sin(2\theta)$ dengan menggunakan
superposisi.
\end{tasks}
\end{exercise}
\exsol{%
a) $u = \frac{-1}{4} r^2 + \frac{1}{4}$
b) $u = \frac{-1}{4} r^2 + \frac{1}{4} + r^2 \sin(2\theta)$
}

\begin{exercise}[menantang]
Turunkan solusi kernel Poisson
jika daerahnya merupakan lingkaran berjari-jari $\rho$, bukan
1.  Artinya, selesaikan $\Delta u = 0$, $u(\rho,\theta) = g(\theta)$.
\end{exercise}
\exsol{%
$\displaystyle
u(r,\theta) = 
\frac{1}{2\pi} \int_{-\pi}^{\pi}
\frac{\rho^2 -r^2}{\rho^2 - 2r\rho\cos(\theta-\alpha) +r^2} g(\alpha) \, d\alpha$
}
